% A scalar finite-volume criterion for a continuum mass gap in SU(2) lattice gauge theory.
% Version 1.3, 12 November 2026: adds Section 7 (flux sectors and 't Hooft twists), Section 8 (the criterion fails at every fixed lattice size as beta -> infinity), two abstract Lean files (Appendix B) and the corresponding statements of the abstract, introduction, discussion and use of AI. Theory only; no simulation data.
% Version 1.2 (doi:10.5281/zenodo.24153901), 12 November 2026: adds the reverse implication (Section 4, Theorem R), a conditional bridge corollary, a second Lean file, and rewrites the comparison with version 5.0.0 of Bhattacharjee et al.
% Versions 1.2, 1.1 (doi:10.5281/zenodo.24152406) and 1.0 (doi:10.5281/zenodo.24150807) are unchanged as records. CritereScalaire.lean, main.lean, Reciproque.lean and main_v1_2.lean are byte for byte those of version 1.2.
% Compile: pdflatex, three passes.
\documentclass[11pt,a4paper]{article}
\usepackage[T1]{fontenc}
\usepackage[utf8]{inputenc}
\usepackage{lmodern}
\usepackage[english]{babel}
\usepackage{amsmath,amssymb,amsthm}
\usepackage{booktabs,array}
\usepackage{geometry}
\geometry{margin=2.2cm}
\setlength{\emergencystretch}{2em}
\usepackage[colorlinks=true,linkcolor=blue,citecolor=blue,urlcolor=blue]{hyperref}
\usepackage{xurl}

\hypersetup{pdftitle={A scalar finite-volume criterion for a continuum mass gap in SU(2) lattice gauge theory },pdfauthor={K\'evin R\'emondi\`ere}}

\theoremstyle{plain}
\newtheorem{theorem}{Theorem}[section]
\newtheorem{proposition}[theorem]{Proposition}
\newtheorem{lemma}[theorem]{Lemma}
\newtheorem{corollary}[theorem]{Corollary}
\theoremstyle{definition}
\newtheorem{example}[theorem]{Example}
\theoremstyle{remark}
\newtheorem{remark}[theorem]{Remark}
\numberwithin{equation}{section}

\newcommand{\R}{\mathbb{R}}
\newcommand{\Z}{\mathbb{Z}}
\newcommand{\N}{\mathbb{N}}
\newcommand{\Tr}{\operatorname{Tr}}
\newcommand{\Hinv}{\mathcal{H}_{\mathrm{inv}}}
\newcommand{\Kh}{\widehat K}
\newcommand{\st}[1]{\textsf{[#1]}}
\newcommand{\cA}{\mathcal{A}}
% --- numbers filled in from the sandbox outputs (paquet/instance, paquet/lean) ---
\newcommand{\BETASTAR}{0.005}
\newcommand{\DELTABOUND}{1.741\cdot10^{-3}}
\newcommand{\BETARATIO}{38}
\newcommand{\COUNTS}{the following. On $2^4$ (degree $18$): $c_m=0$ for $m\le3$ and for $m=5$, $c_4=1$, $c_6=4$, i.e.\ $N_4=24$ (the sheets) and $N_6=64$ (the unit cubes), and no other connected closed set of at most six plaquettes. On $4^4$ (degree $20$): no closed set of at most five plaquettes, and $N_6=1024$ (the unit cubes, $c_6=4$)}
\newcommand{\REBUILDJOB}{20261111-005922-02de92}
\newcommand{\MUTANTS}{All eight made the build fail (job identifiers in \path{lean/mutants_chantier/resultats.json}). They were
run on the text of an intermediate build, before the non-vacuity guard was added, and we did not check that each failure is located in
the targeted declaration. A re-run on the final text with localisation of the errors (\path{lean/mutants_paquet.py}) was started on
11 November 2026 but could not run: the sandbox had no free slot.}
\newcommand{\REVIEWSTATUS}{reviewed the assembly by reading, without running programs (11 November 2026): it confirmed by hand the constants of the chain, the torus--cylinder lemma, the diagonal choice, the bracket, the $\Z_2$ counterexample and the role of the factor $16$, pointed out the unlisted bilinearity hypothesis and the untreated compact $U(1)$ Coulomb phase (both now stated), and objected that (W2) assumes a simple $\lambda_0$; this objection is not retained, since degenerate eigenvalues are allowed as $r_i=1$, as a later reviewer (GPT-6.1 Sol) also found. A second, line-by-line review of an earlier draft by Qwen, again without running programs, found that the torus--cylinder lemma needs a positivity-preserving kernel (now hypothesis (W2$^+$)), that the rounded constants of the chain must be read as upper bounds (now stated), and left the numerical inputs of the strong-coupling instance and of the witnesses unchecked.}

\title{A scalar finite-volume criterion for a continuum mass gap\\ in $SU(2)$ lattice gauge theory:\\ the doubling test in dimensions four and three,\\ its continuum assembly and its converse}
\author{K\'evin R\'emondi\`ere\thanks{Independent researcher, Oloron-Sainte-Marie, France. ORCID
\href{https://orcid.org/0009-0008-2443-7166}{0009-0008-2443-7166}. The proofs, the programs, the Lean files and this note were
prepared with the assistance of AI models (Claude Opus~5.5, Claude Sonnet~5.5 and Claude Fable~5.2, Anthropic; Qwen3.8, Qwen team; GPT-6 Astra and GPT-6.1 Sol, OpenAI; DeepSeek V4 Pro, DeepSeek); their roles are described in the
statement on the use of AI at the end of the note.}}
\date{12 November 2026 (version 1.3)}

\begin{document}
\maketitle

\begin{abstract}
This technical note transposes the two-dimensional finite-volume doubling criterion of OpenAI (Proposition~2.2 of \emph{Sharp mass bounds
for the two-dimensional $O(4)$ model}, 2026) to lattice gauge theories in dimensions four and three. It is predominantly a transposition,
not a new mass-gap mechanism. For the $SU(2)$ Wilson theory on periodic tori set
$\Delta_\beta(n)=16\log Z_\beta(n^4)-\log Z_\beta((2n)^4)$.
(i)~Under hypercubic symmetry and a positive trace-class transfer representation (a standard hypothesis, not formalised),
$\Delta_\beta(n)\le2^{-9}$ implies $\Delta_\beta(2n)\le256\,\Delta_\beta(n)^2$ and
$(\lambda_i/\lambda_0)^n\le256\,\Delta_\beta(n)\le\frac12$ for every excited eigenvalue of the transfer operator on every spatial torus
$(2^jn)^3$; in dimension three the threshold is $2^{-7}$ and the constant $64$.
(ii)~Fix $\ell>0$ and $0<\eta\le2^{-9}$. Assume (W1)--(W3) and $\Delta_{\beta_k}(n_k)\le\eta$ for every $k$, at spacing $a_k=\ell/n_k$, and assume (W4): the lattice
Osterwalder--Schrader forms converge, which is the existence of the continuum limit and is \emph{assumed}. Then the limit satisfies
$\langle v,e^{-\ell H}v\rangle\le256\eta\|v\|^2$ on the complement of the vacuum, with the same constant for every $k$; given the
Osterwalder--Schrader reconstruction (also assumed), it has a mass gap at least $-\log(256\eta)/\ell\ge(\log2)/\ell$ and a unique vacuum.
This implication is formalised in Lean~4 against our own Lean statement of the problem (ClayYM v1.1), with the transfer representation, the
convergence and the reconstruction as named hypotheses (in the Lean statement the hypothesis on the transfer structure, \texttt{SpecLink}, is weaker than (W3), and the specialisation to Wilson's partition function requires a symmetric extension to periods $1$ satisfying the spectral identities, which is an additional assumption, as stated under Table~\ref{tab:hyp}); the Lean sources were rebuilt on 11 November 2026 by non-authors on a clean
machine, first from hashed copies of the sources and then from the git archive of the package tree (Appendix~\ref{app:lean}). For square tori the passage from cylinders uses a further hypothesis (W2$^+$) and is proved on paper for one
diagonal choice of the volumes only.
(iii)~$\Delta$ is bracketed by excited-state traces on spatial tori with sides in $\{\ell,2\ell\}$; a finite-volume gap including the
electric-flux sectors, together with a trace bound, gives an explicit sufficient inequality for the criterion; examples show that the criterion can fail when either ingredient is missing, but the example for the flux sectors has a trivial local limit, and the example for the trace bound consists of decoupled copies of an abstract theory, so that strictness within non-trivial theories, or within pure $SU(2)$, is not established.
(iv)~For a stipulated Maxwell spectrum (whose identification with a normalised Villain measure is not derived here) the defect exceeds the
threshold at every size $n\ge3$ and every coupling; a massive free scalar passes numerically for $M\ell\ge11.47$.
(v)~A single implementation reports the sufficient bound $\Delta_\beta(2)\le1.741\cdot10^{-3}$ for $\beta\le0.005$. This interval is conditional on the finite enumeration (ENC) and rational comparison (RAT), neither independently certified here.

\emph{Relation to earlier public work.} Version~5.0.0 of the series of Bhattacharjee et al.\ \cite{BB26v5}, whose Zenodo record was
created on 11 November 2026, before the release of version 1.0 of this note, and which we read after that version was written, contains for four-dimensional
tori the same doubling defect, its bracket by thermal traces, the purity reading, and a one-scale criterion whose hypothesis (B$^*$), a family of
inequalities between partition functions of finite tori, implies ours (\S\ref{sec:bb26v5}). \emph{Version~1.2 proves the reverse implication} (Theorem~\ref{thm:R}):
under symmetry and a positive transfer representation, $\Delta_\beta(n)\le2^{-9}$ implies (B$^*$) at the scale $16n$, and the gap
$(\log2)/n$ on every integer cube of side $N\ge n$. The two criteria are therefore equivalent up to a change of scale: the reverse implication loses a factor $16$ (not optimal), and a round trip, (B$^*$) at $T_0$ $\Rightarrow$ $\Delta(2T_0)\le2^{-9}$ $\Rightarrow$ (B$^*$) at $32T_0$, a factor $32$. The proof is paper-only for the aspect-ratio statement, and the main steps are checked by the Lean kernel in an abstract form (Appendix~\ref{app:lean}).
This does not answer a question left open in \cite{BB26v5}, whose data differ from ours, and it is not a proof of the mass gap. A corollary linking our criterion to their
Theorem~F is stated as conditional on their hypothesis (W) and their Theorem~F as we read them. No priority is claimed for the one-scale criterion, the defect or the monotonicity of the free energy in a period (\S\ref{sec:bb26v5}).

\emph{What the note does not do.} It gives no estimate of $\Delta$ at weak coupling, no existence or non-triviality of the continuum limit,
and no reconstruction. It does not apply as it stands to the formalised construction of OpenAI (Section~\ref{sec:openai-construction}).
Reading $\Delta$ as a combination of R\'enyi-2 entropies of the thermal weights is a reformulation, not a result. The criterion is a sufficient
condition; its equivalence with the Yang--Mills mass-gap problem is not claimed. The proofs have been checked by AI models only.

\emph{Version 1.3 (12 November 2026), theory only.} Two sections and two abstract Lean files are added. (1)~Section~\ref{sec:thooft}: for the Wilson $SU(2)$ theory the periodic trace splits exactly over the electric-flux sectors; the criterion $\Delta(n)\le2^{-9}$ at one size implies disorder of the 't~Hooft twists, $0\le1-Z_k/Z_0\le q^{N/n}/124$, on every hypercube of side $N\ge n$ (a length $N/n$, not an area: no area law in infinite volume follows); the criterion is equivalent, up to constants, to two inequalities at one scale, twist disorder (V1) and a neutral-sector bound (V2), each failed by a different model; a necessary growth of $n$ with $\beta$ is derived \emph{conditionally} on Theorem~7.59 of \cite{BB26v5}. (2)~Section~\ref{sec:fixed}: at every fixed lattice size the criterion \emph{fails} as $\beta\to\infty$: the twist ratio of $SU(2)$ is bounded by a power of $\beta$ (Laplace's method on isolated twisted vacua), while that of $U(1)$ and $\Z_2$ decays exponentially, and $\liminf_{\beta\to\infty}\Delta_\beta(n)\ge45\log2$. This is a statement of where the criterion fails, not a mass-gap result; the constants are not explicit; the ingredients are classical and we do not claim novelty for the asymptotics. The proofs are on paper, checked by AI models only; the abstract Lean files check the finite algebra and nothing analytic (Appendix~\ref{app:lean13}).

\emph{Version 1.1 (12 November 2026): literature corrections only.} The principle that an inequality between partition functions of finite periodic tori (twisted against untwisted) implies decay is classical (Tomboulis--Yaffe, Kov\'acs--Tomboulis, Kanazawa), and the comparison with these papers is recorded in \S\ref{sec:compare}.

\emph{Version 1.2 (12 November 2026; doi:10.5281/zenodo.24153901).} Adds Section~\ref{sec:reverse} (Theorem~\ref{thm:R}, its remarks and the conditional Corollary~\ref{cor:bridge}), the second Lean file \path{lean/Reciproque.lean} and its combination with the published file \path{lean/main_v1_2.lean}, an example showing that positivity of the transfer operator cannot be dropped (Remark~\ref{rem:Rused}), and rewrites \S\ref{sec:bb26v5} and the related statements. The statements, proofs and Lean files of versions 1.0 and 1.1 are not changed: \path{lean/CritereScalaire.lean} and \path{lean/main.lean} are byte for byte those of version 1.0.
\end{abstract}

\paragraph{Status labels.}
\st{proved}: complete argument written in this note; \st{Lean}: also checked by the Lean~4 kernel (Appendix~\ref{app:lean}), modulo
the named hypotheses listed there; \st{computed}: floating-point or exact computation in the sandbox described in
\S\ref{sec:repro}, never used in a proof unless stated; \st{read}: statement read in the cited source, location given;
\st{hypothesis}: a named assumption, not proved here; \st{heuristic}.

{\small\setcounter{tocdepth}{1}\tableofcontents}

%=====================================================================================================================
\section{Introduction}\label{sec:intro}

\subsection{The doubling test and where it comes from}
Let $Z(n_1,n_2)$ be the periodic partition function of a two-dimensional spin model whose transfer operator is self-adjoint, positive
and trace class. OpenAI \cite[\S2]{OAI26} introduce the \emph{square-box test} $\Delta(n)=4\log Z(n,n)-\log Z(2n,2n)$, in which the
coefficients cancel every term proportional to the area, and prove \cite[Prop.~2.2]{OAI26} \st{read}: if $n\ge4$ is even and
$\Delta(n)\le\eta_*:=2^{-10}$, then the transfer operators on circumferences $2^kn$ satisfy $\lambda_2/\lambda_1\le e^{-(\log 2)/n}$,
every local limit of the square tori of sides $2^kn$ has full gap at least $(\log2)/n$, and $m_{\mathrm{lat}}(\beta)\ge(\log 2)/n$ if
the periodic infinite-plane state is unique. The mechanism has four parts: the identity $e^{-p}=\sum_it_i^2$ for the normalised trace
weights and the squaring lemma \cite[Lemma~2.1]{OAI26}; the rectangle identity \cite[(2.7)]{OAI26}; the quadratic recursion
$\Delta(2r)\le44\Delta(r)^2$ \cite[(2.12)]{OAI26}; and the comparison of square tori with infinite cylinders.

\textbf{The present work is predominantly a transposition of this mechanism} to lattice gauge theories in dimensions $d=4$ and $d=3$,
followed by a conditional continuum statement; it is not a new mass-gap mechanism. We claim no priority for the mechanism. The transposition
to $d=4$ was first written in internal, unpublished work of the same project (25 October 2026: threshold $2^{-12}$, constant $512$, gap
$(\log8)/n$; and 10 November 2026: all dimensions, constants in Table~\ref{tab:const}); the present note supersedes those notes.

The identity $e^{-p}=\sum_it_i^2$ says that $p$ is the R\'enyi-2 entropy of the normalised thermal weights $t_i$; OpenAI already read it as
a purity \cite[\S2]{OAI26}. Reading $\Delta_\beta$ as a combination of such R\'enyi-2 entropies of the transfer spectrum
is a reformulation of the same quantities: it adds no hypothesis and no result (it was written in our internal notes of 25 October 2026),
and nothing below depends on it. The same reading, as the purity $\Tr\rho^2$ of a thermal state, is made for four-dimensional tori in
\cite[after Proposition~7.57]{BB26v5}.

\subsection{Results}
Let $F_\beta=\log Z_\beta$ on rectangular tori $\Lambda(\vec n)=\prod_{\mu=1}^4\Z/n_\mu\Z$ (\S\ref{sec:setting}) and
\begin{equation}\label{eq:Delta}
\Delta_\beta(n)=16\,F_\beta(n,n,n,n)-F_\beta(2n,2n,2n,2n),\qquad
\Delta^{(3)}_\beta(n)=8\,F_\beta(n,n,n)-F_\beta(2n,2n,2n).
\end{equation}
\begin{itemize}
\item \emph{Lattice part} (Theorem~\ref{thm:lattice}, \st{proved}, \st{Lean}). Under (W1)--(W2): $\Delta_\beta(n)\le2^{-9}$ implies
$\Delta_\beta(2n)\le256\,\Delta_\beta(n)^2$ and, on every spatial torus $(2^jn)^3$, $(\lambda_i/\lambda_0)^n\le256\,\Delta_\beta(n)$ for
every excited eigenvalue of the transfer operator. In $d=3$: threshold $2^{-7}$, constant $64$ (Theorem~\ref{thm:3d}).
\item \emph{Reverse implication} (Theorem~\ref{thm:R}, \st{proved}, \st{Lean} for its main steps in abstract form). $\Delta_\beta(n)\le2^{-9}$ implies the one-scale bound (B$^*$) of \cite{BB26v5} at the scale $16n$, and the gap $(\log 2)/n$ on every integer cube of side $N\ge n$ (hence not only the dyadic cubes). Positivity of the transfer operator is used and cannot be dropped. A corollary linking the criterion to Theorem~F of \cite{BB26v5} is \emph{conditional on their hypothesis (W) and their Theorem~F as we read them, and not verified by us}.
\item \emph{Conditional assembly} (Theorem~\ref{thm:main}, \st{proved}, \st{Lean} modulo (W2)--(W4), (W6)). If $0<\eta\le2^{-9}$ and
$\Delta_{\beta_k}(n_k)\le\eta$ uniformly in $k$, with $a_k=\ell/n_k$, then every limit of the cylinder OS forms, in the sense of (W3)--(W4), has
the $q$-form with $q=256\eta$. Under the reconstruction hypothesis (W6) its mass gap is at least $-\log(256\eta)/\ell\ge(\log2)/\ell$ and its
vacuum is unique. For square tori the same conclusion needs (W5) for the family actually considered; Corollary~\ref{cor:tori} proves it for one
diagonal choice of volumes, not for an arbitrary sequence $j_k\to\infty$.
\item \emph{Thermal-trace converse} (Proposition~\ref{prop:bracket}, Corollary~\ref{cor:converse}, \st{proved}). $\Delta$ is bracketed,
up to the factor $2(2^d-1)$, by the excited trace at time $\ell$ on $d$ spatial tori with sides in $\{\ell,2\ell\}$. A finite-volume gap
(flux sectors included) and a trace bound give an explicit sufficient inequality for the criterion at a fixed physical size; eventual
smallness as $\ell\to\infty$ needs control of the dependence of the trace bound on $\ell$. Examples~\ref{ex:copies} and~\ref{ex:z2} show that
neither the finite-volume gap nor the trace bound can be dropped.
\item \emph{Spectral witnesses} (\S\ref{sec:witness}). For the stipulated Maxwell (photon, rotor, flux) spectrum, whose identification
with a normalised monopole-free Villain measure is not derived here, the defect exceeds the threshold by a factor $\ge90$ at every size $n\ge3$
and every coupling \st{proved} given that spectrum (numerically by $\ge1.7\cdot10^4$ \st{computed}). A massive free scalar passes numerically
for $M\ell\ge11.47$ ($d=4$) and $M\ell\ge8.47$ ($d=3$) \st{computed}; these crossings come from a truncated spectrum and carry no tail
certificate.
\item \emph{Reported computational instance} (Proposition~\ref{prop:instance}). Conditional on (ENC) and (RAT), $\Delta_\beta(2)\le1.741\cdot10^{-3}<2^{-9}$ for the $SU(2)$ Wilson theory and $0\le\beta\le0.005$. The analytic implication is proved on paper; the two computational inputs are reported from one implementation and remain uncertified here. The gap is already known in this regime.
\item \emph{Flux sectors and 't~Hooft twists} (Section~\ref{sec:thooft}, version~1.3). Exact decomposition of the excited trace into a neutral part and a twist part (Proposition~\ref{prop:sectors}, \st{proved}); $\Delta(n)\le2^{-9}$ implies $1-Z_k/Z_0\le q^{N/n}/124$ on every hypercube $N\ge n$ (Theorem~\ref{thm:twistall}, \st{proved}, resting on Theorem~\ref{thm:R}); equivalence of the criterion with two one-scale conditions (Theorem~\ref{thm:twosub}, \st{proved}); a necessary scale $n\gtrsim(\beta/(8\log\beta))^{1/4}$ (Proposition~\ref{prop:scale}, \st{conditional} on Theorem~7.59 of \cite{BB26v5}). The bound is a length, not an area.
\item \emph{The criterion fails at fixed size} (Section~\ref{sec:fixed}, version~1.3). For $SU(2)$, $Z_k/Z_0\le C\beta^{-(d-1)/2}$ at fixed torus (Theorem~\ref{thm:L}, \st{proved} on paper, constants not explicit), against $e^{-c\beta}$ for $U(1)$ and $\Z_2$; hence $\liminf_{\beta\to\infty}\Delta_\beta(n)\ge45\log2$ at every $n\ge3$ (Corollary~\ref{cor:L1}, \st{proved} under (W1), (W2), (C)). The ingredients are classical (zero-action twisted configurations, the vacuum valley of the untwisted torus, the marginal constant-mode integral of $SU(2)$ in dimension four); the bound itself was not found in the texts we could read, and six relevant papers could not be obtained (Section~\ref{sec:fixed}).
\end{itemize}

\subsection{What is new and what is not}\label{subsec:notnew}
\emph{Not new.} The doubling mechanism, in $d=2$ for spins \cite{OAI26}. The positive, self-adjoint transfer operator of lattice gauge
theories \cite{Lus77,OS78}. The closure lemma ``same $q$ along any filter'' and the passage from the $q$-form to a gap through a named
reconstruction hypothesis, which are part of our archived Lean statement ClayYM v1.1 \cite{ClayYM}. The form of the result, ``a local
hypothesis at a fixed physical scale $\ell$, uniform in the lattice spacing, plus existence of a limit, implies a gap of every limit'',
which is that of Theorem~B of Bhattacharjee et al.\ \cite{BB26} and of Theorem~5.3 of Rajabu \cite{Raj26}, with different hypotheses
(\S\ref{sec:compare}). The doubling defect of four-dimensional tori, its decomposition over the four directions, its bracket by thermal
traces and a one-scale criterion with a gap conclusion, which are in version~5.0.0 of Bhattacharjee et al.\ \cite{BB26v5}
(\S\ref{sec:bb26v5}). The unconditional squaring bound $p(2m)\le2(e^{p(m)}-1)^2$ (Lemma~\ref{lem:squaring}) is implicit in the chain of
inequalities of \cite[proof of Lemma~2.1]{OAI26}. Step~1 of the proof of Theorem~\ref{thm:R} rests on the elementary monotonicity of the counting-measure $\ell^p$ norms: for $0<L\le L'$, whenever the traces are finite, $(\Tr T^{L'})^{1/L'}\le(\Tr T^{L})^{1/L}$ for a positive operator $T$. Lemmas~7.12 and~7.44 of \cite{BB26v5} contain related monotonicity statements (of the thermal trace in the time, and of the tube energy in a side); the precise statement on the pressure used here follows from the spectral representation. The same concavity argument, with dyadic enlargement in every direction, is used long before by Brydges, Fr\"ohlich and Seiler \cite[proof of (2), p.~38 of the preprint]{BFS79}, to bound a lattice gauge partition function below by a fractional power of its value on larger periodic rectangles: $\Tr(T^{L_0})\ge\bigl(\Tr T^{2^{m_0}}\bigr)^{L_0/2^{m_0}}$ for a positive semi-definite transfer matrix $T$ and $2^{m_0}\ge L_0$, then the same in the other directions, which gives $\langle e^{-A_\Lambda}\rangle_0\ge\langle e^{-A_{\widetilde\Lambda}}\rangle_0^{|\Lambda|/|\widetilde\Lambda|}$ in their notation (read on the image of the page; the optical character recognition of our copy garbles the exponents). We claim no novelty for the monotonicity.

The principle that an inequality between the partition functions of finite periodic tori, twisted against untwisted, implies a decay: for Wilson loops of $SU(2)$, $\langle W[C]\rangle\le2(e^{-F_{\mathrm{el}}})^{|A_C|/|A_{12}|}$ with $e^{-F_{\mathrm{el}}}=\langle\frac12(1-\mathcal O[\mathcal V])\rangle$, in Kov\'acs and Tomboulis \cite[eqs.~(6) and~(11), pp.~4--5]{KT02}, and restated as Theorem~1 of Kanazawa \cite[p.~5]{Kan}, who attributes it to Tomboulis and Yaffe \cite{TY85} and to \cite{KT02}; for $SU(N)$ and $\mathbb Z_N$ fluxes, in Kanazawa \cite[Theorem~2, p.~7]{Kan}; and, for correlations of classical spin systems, in Kanazawa \cite[Theorem~5, p.~17]{Kan}, for $L_\mu=2^kn$: $|\langle\Gamma_\mu(x;n)\rangle|/\langle\Gamma_\mu(x;0)\rangle\le2\{\langle\Gamma_\mu(x;0)^2\rangle/\langle\Gamma_\mu(x;0)\rangle^2\}^{n/L_\mu}\{1-\int dg\,\langle\mathcal O(g)[\mathcal V]\rangle\}^{n/L_\mu}$. His abstract (p.~1) states: ``Basically the inequalities guarantee that a system in a box that is sufficiently insensitive to boundary conditions has a non-zero mass gap.'' The hypothesis there is a ratio of twisted to untwisted partition functions, not the doubling defect, and the conclusion is a decay in a box of period $L$, not a bound on the transfer matrix of every spatial torus (\S\ref{sec:compare}); the family of the argument (reflection positivity, Schwarz inequality, doublings) is the same as ours. The absence of these papers from the list of version~1.0 was a gap of the anteriority search.

\emph{Incremental content.} The contribution is predominantly a routine transposition of the doubling mechanism of \cite{OAI26}.
It consists of the higher-dimensional bookkeeping and the explicit bounds $256$ at $2^{-9}$ in $d=4$ and $64$ at $2^{-7}$ in $d=3$; the
packaging of the conditional implication against the Clay-type $q$-form; an elementary thermal-trace converse with counterexamples to omitted
hypotheses; and the witness analysis. The torus--cylinder comparison follows the insertion argument of \cite{OAI26}, with an additional
diagonal qualification when renormalised lattice observables depend on the cutoff. The reported $SU(2)$ interval up to $\beta=0.005$ is conditional on a
finite enumeration and a rational comparison that are not independently certified. We claim no new mass-gap mechanism, no optimality of the constants, and no
exhaustive priority determination: the anteriority search of \S\ref{sec:compare} is limited and recorded. In \cite{OAI26} it found neither
a converse nor a statement in dimension $d>2$. In \cite{BB26v5}, read on 12 November 2026 after the first draft of version~1.0 had been written, it found the
four-dimensional defect, its bracket and a one-scale criterion; the part of the content listed above that is also there, and what we did
not find there, are stated in \S\ref{sec:bb26v5}.

\subsection{Relation to version 5.0.0 of Bhattacharjee et al.}\label{sec:bb26v5}
Version~5.0.0 of the series of Bhattacharjee et al.\ \cite{BB26v5} has a Zenodo record created on 11 November 2026 at 23:29:13 UTC, before
the release of this note. We read it on 12 November 2026, after the first draft of version~1.0 had been written, in its abstract, its
introduction and its Sections~7.8 and~7.11; for version~1.2 we also read the lemmas and propositions cited in Section~\ref{sec:reverse} and here (7.3, 7.12, 7.27, 7.33, 7.44, 7.45, 7.48, 7.51, 7.56, 7.57, 7.59) and searched the whole text by keywords. For a spatial torus $s$ write $\widehat Z_s(M)=\sum_{i\ge0}r_i^M$, where $1=r_0\ge r_1\ge\dots\ge0$
are the eigenvalues of the transfer operator of $s$ divided by the largest one (their relative partition function, the thermal trace at
time $M$), and $V(s)$ for the number of sites of $s$. Those passages contain the following \st{read}.
\begin{itemize}
\item[(a)] \emph{The object and its bracket} (their Lemma~7.33 and Proposition~7.34). For a four-dimensional torus $n=(n_1,\dots,n_4)$ with
all sides at least $3$: the doubling defect $\delta(n)=16\log Z(n)-\log Z(2n)$, its invariance under $Z(n)\mapsto e^{a|n|}Z(n)$, its
decomposition over the four directions with the weights $8,4,2,1$, and
\begin{multline*}
8\log\widehat Z_{(n_1,n_2,n_3)}(n_4)\ \le\ \delta(n)\ \le\ 16\log\widehat Z_{(n_1,n_2,n_3)}(n_4)+8\log\widehat Z_{(n_1,n_2,2n_4)}(n_3)\\
+4\log\widehat Z_{(n_1,2n_3,2n_4)}(n_2)+2\log\widehat Z_{(2n_2,2n_3,2n_4)}(n_1).
\end{multline*}
On cubic tori $\delta(n)$ is our $\Delta_\beta(n)$, the decomposition is \eqref{eq:decomp}, the upper bound is that of
Proposition~\ref{prop:bracket}, and the lower bound is the first term of the lower bound of that proposition.
\item[(b)] \emph{A one-scale criterion} (their Theorem~7.56 and Proposition~7.57). Hypothesis (B$^*$), for an integer $T_0\ge3$:
$\log\widehat Z_s(T_0)\le\gamma V(s)$ for every spatial torus $s$ with all its sides in $[2T_0,48T_0)$, where $2\gamma(96T_0)^3\le\frac1{20}$.
Conclusion: a doubly exponential decrease of these thermal traces along the times $2^iT_0$, and a gap at least $0.172/T_0$ for the transfer
operator of every spatial torus whose smallest side $a$ is at least $2T_0$ and whose largest side is less than $24a$. Up to a factor $2$ in
$\gamma$, (B$^*$) is the family of inequalities $Z(s,T_0)^2\le e^{\gamma V(s)}Z(s,2T_0)$ between partition functions of finite tori, which
they read as near purity, $\Tr\rho^2\ge e^{-\gamma V(s)}$, of the thermal state.
\item[(c)] \emph{Scope remarks and a continuum statement.} Such a hypothesis contains a lattice gap (their Proposition~7.58); a gap and
the free energy per site do not control the doubling defect (their Proposition~7.35); and, for continuum families satisfying their
hypothesis (W), which they describe as ``the shape of the conclusion claimed by the unrefereed Lean~4 construction'' (their references
[45, 62]; their [45] is the Wilson Yang--Mills construction of OpenAI), (B$^*$) at one physical time extent $t_0$ gives a Hamiltonian with spectrum in $\{0\}\cup[0.172/t_0,\infty)$ and a simple eigenvalue $0$ (their Theorem~F).
\end{itemize}

\emph{What is theirs.} Items (a)--(c) are earlier than this note. They cover the object $\Delta$, the decomposition \eqref{eq:decomp}, the
bracket of Proposition~\ref{prop:bracket} in $d=4$, the purity reading, the doubly exponential mechanism in $d=4$, and the observation that
the criterion contains a finite-volume gap. We claim no priority for any of them; our internal notes \cite{internal} are unpublished and
establish none.

\emph{Their hypothesis implies ours.} Let $n=2T_0$. The four spatial tori that enter the upper bound of (a) for the cubic torus $n^4$ have
sides in $\{2T_0,4T_0\}$, hence in the window of (B$^*$), and volumes $2^jn^3$, $j=0,\dots,3$. Since $M\mapsto\widehat Z_s(M)$ is
non-increasing, (B$^*$) gives $\log\widehat Z_s(n)\le\log\widehat Z_s(T_0)\le\gamma V(s)$ for each of them, and
\begin{equation}\label{eq:bb26implies}
\Delta_\beta(2T_0)\ \le\ \gamma n^3\,(16+8\cdot2+4\cdot4+2\cdot8)\ =\ 512\,\gamma\,T_0^3\ \le\ \frac{512}{40\cdot96^3}\ =\ \frac1{69120}\ <\ 2^{-9}.
\end{equation}
\st{proved}, from the statements of \cite{BB26v5} as read. \emph{Conversely}, under (W1)--(W2), $\Delta_\beta(n)\le2^{-9}$ implies (B$^*$) at $T_0=16n$ with $\gamma=1/(31\cdot2^{23}T_0^3)$ (Theorem~\ref{thm:R}, Section~\ref{sec:reverse}). Version 1.0 and 1.1 of this note said that we did not know this; the proof was found on 12 November 2026. Hence, up to a change of scale (a factor $16$ in the reverse direction, not optimal; $32$ for a round trip), the two criteria are equivalent. Note that by Proposition~\ref{prop:bracket} alone the defect bounds the thermal traces of four boxes only; the extension to all tori of the window is obtained from the monotonicity of the pressure in each period (Step~1 of Theorem~\ref{thm:R}).

\emph{What differs.} Our hypothesis is one inequality between the partition functions of the two cubic tori $n^4$ and $(2n)^4$; theirs is a
family of inequalities, one for each spatial box with sides in a window. Their conclusion is wider than that of Theorem~\ref{thm:lattice}:
every torus of bounded aspect ratio above the base scale, not only the dyadic cubes $(2^jn)^3$. Theorem~\ref{thm:R} shows that this difference is one of scale, not of substance: the two hypotheses are equivalent up to a change of scale (a factor $16$ in the reverse direction, $32$ for a round trip), and the conclusion of Theorem~\ref{thm:lattice} extends to all integer cubes of side $\ge n$ and, with a worse constant, to all spatial tori with $\min_it_i\ge n$ and $\max_it_i\le24\min_it_i$. For the Wilson theory in $d=4$, the gap
statement \eqref{eq:gap} of Theorem~\ref{thm:lattice} is thus a variant of that of their Theorem~7.56; their recursion is on normalised bounds of the thermal traces, $y_{i+1}=16(e^{y_i}-1)^2$, not on the defect. In the passages read
we did not find: the recursion $\Delta(2n)\le256\,\Delta(n)^2$ for the cubic defect itself, and with it a gap on all dyadic sections from a
single inequality; the three-dimensional statement (Theorem~\ref{thm:3d}); a Lean formalisation of such a criterion
(Appendix~\ref{app:lean}); the witnesses of \S\ref{sec:witness}; the reported instance of Proposition~\ref{prop:instance}; and, by the complete reading and keyword search of 12 November 2026, the implication from a small defect at one size to (B$^*$) (Theorem~\ref{thm:R}), the comparison of $f(m^4)$ with the limiting pressure of a torus by telescoping the pressure differences along $m_j=2^jn$ (with the constant $2/31$), and the passage from a small $\Delta(n)$ at one size $n$ to a gap on all cubes of side $N\ge n$. This is not a proof of absence. What \emph{is} in \cite{BB26v5} and bears on this: the telescoping $\log\widehat Z(M)=\sum_{k\ge0}2^{-k-1}\psi(2^kM)$ for one transfer operator (their Lemma~7.33, a different telescoping, in the time), and, from Proposition~7.34(ii), the fact that a small defect $\delta(n)$ of the torus $n$ \emph{itself}, dyadic or not, bounds the thermal trace of its spatial torus $(n_1,n_2,n_3)$ at the time $n_4$ (our bracket, Proposition~\ref{prop:bracket}). So the quantity $\Delta$ and the non-dyadic case at the size where the defect is small are theirs; the new points are the implication to (B$^*$) (and (B$^{**}$), Remark~\ref{rem:Bss}) and the extension to all sizes $N\ge n$ from one size $n$. Their continuum
statement (c) and Theorem~\ref{thm:main} have different hypotheses, and we have not compared them; Corollary~\ref{cor:bridge} is conditional on (c) as we read it.

\emph{Limits of this comparison.} It was made by AI models (see ``Use of AI'') from the passages cited, on 12 November 2026; the rest
of the 157 pages was searched by keywords only. We have not examined in which version of that series each item first appeared, and no human
expert has checked the comparison.

\subsection{Status of the proof}\label{sec:status}
The abstract implication from $\Delta\le\eta$ to the $q$-form and, under the normalised reconstruction hypothesis, to \texttt{HasMassGap}
of ClayYM is formalised in Lean~4 (35 declarations, standard axioms only) under the stated positive-period spectral hypotheses. The paper proof on periods at least $2$ does not require the extension to periods $1$ described under Table~\ref{tab:hyp}; its identification with that particular Lean instance is conditional on that extension. The build was reproduced on 11 November 2026 by non-authors on a
clean machine, first from hashed copies of the sources and then from the git archive of the package tree (Appendix~\ref{app:lean}). The transfer-spectrum representation
(W2)--(W3), its identification with the lattice OS forms, the convergence (W4) and the reconstruction (W6) are named hypotheses and are not
formalised; (W4) is the existence of the continuum limit and is open. The insertion estimate and the quantitative torus--cylinder comparison
are paper arguments (Lean assumes that the difference tends to zero); the thermal-trace converse of Section~\ref{sec:converse} is a paper argument. The reverse implication of Section~\ref{sec:reverse} is a paper argument whose steps (a), (b) and the first half of (c) are also checked by the Lean kernel in an abstract form, with the symmetry, the spectral representation and the published doubling recursion as named hypotheses; the aspect-ratio statement (d), the logarithmic form of (c), the formula of Corollary~\ref{cor:bridge} and its link with \cite{BB26v5} are not formalised. The Maxwell model identification is
incomplete, the scalar crossings are numerical, and the strong-coupling interval is conditional on an enumeration and a rational comparison reported by one implementation. The original formalisation and draft were prepared within one model family. Read-only reviews by Qwen3.8 and GPT-6.1 Sol concerned the draft of 11 November, not this final text: they found the core implication sound subject to corrections and left the computational inputs uncertified; Claude Sonnet applied adaptations of their proposed corrections, and some proposed additions were not adopted. A read-only review of the first version of this text by GPT-6.1 Sol required sixteen corrections, which Claude Sonnet adapted. A further read-only check of this revision identified remaining corrections, recorded in \path{RELECTURE_SOL_TEXTE_CORRIGE.md} (one blocking and four non-blocking: the conditional status of Proposition~\ref{prop:instance} in the abstract and the list of results, the description of \texttt{SpecLink}, the concordance table, the attribution of the reviews, the wording of the rebuild descriptions); they are applied in this version. These reviews do not certify the computational inputs or endorse a subsequent revision. The text--Lean concordance was prepared within the author's model family. Appendix~\ref{app:lean} states, for each formalised statement, what the Lean declaration says exactly. \emph{Version 1.3.} Sections~\ref{sec:thooft} and~\ref{sec:fixed} are paper arguments; The decomposition and the equivalence (Proposition~\ref{prop:sectors}, Theorem~\ref{thm:twosub}) and Corollary~\ref{cor:L1} were read line by line by the pilot agent, Theorem~\ref{thm:L} was re-derived blind by GPT-6 Astra and attacked by GPT-6.1 Sol, and two reviewers of other families (DeepSeek V4 Pro, without tools) attacked the final text (\path{RELECTURES.md}); the abstract Lean files of Appendix~\ref{app:lean13} check only finite algebra. Hypothesis (C) for Wilson's theory rests on the paper proof of \cite{TBv12}, checked by models only. No human has checked the proofs.

%=====================================================================================================================
\section{Setting, hypotheses and main theorem}\label{sec:setting}

\subsection{Model}
For $\vec n=(n_1,\dots,n_4)$ with $n_\mu\ge2$, let $\Lambda(\vec n)=\prod_\mu\Z/n_\mu\Z$, with links $(x,\mu)$ and plaquettes
$p=(x;\mu,\nu)$, $\mu<\nu$; for $n_\mu\ge2$ every plaquette has four distinct links. The Wilson action and partition function are
\begin{equation}\label{eq:wilson}
S_\beta(U)=\beta\sum_p\bigl(1-\tfrac12\Tr U_p\bigr),\qquad Z_\beta(\vec n)=\int e^{-S_\beta(U)}\,dU,\qquad F_\beta(\vec n)=\log Z_\beta(\vec n),
\end{equation}
with $U_p=U_{x,\mu}U_{x+\mu,\nu}U_{x+\nu,\mu}^{-1}U_{x,\nu}^{-1}$, $dU$ the product of normalised Haar measures on $SU(2)$ and
$\beta=4/g^2$; $Z_0=1$. At lattice spacing $a$, a period $n$ is a physical length $na$. The three-dimensional model is defined in the
same way on $\prod_{\mu=1}^3\Z/n_\mu\Z$. For a profile $m$ (a multiset of four periods) we write $\tau_\mu m$ for the profile with the
$\mu$-th period doubled and
\begin{equation}\label{eq:p}
p_\mu(m)=2F(m)-F(\tau_\mu m).
\end{equation}

\subsection{Hypotheses}\label{sec:hyp}
Table~\ref{tab:hyp} lists every hypothesis used, with its status. The names (W1)--(W6) are used throughout, and in the Lean files.

\begin{table}[ht]
\centering\footnotesize
\begin{tabular}{@{}>{\raggedright\arraybackslash}p{0.06\textwidth}>{\raggedright\arraybackslash}p{0.50\textwidth}%
>{\raggedright\arraybackslash}p{0.36\textwidth}@{}}
\toprule
name & statement & status \\
\midrule
(W1) & $F_\beta$ is invariant under permutations of the four periods. & \st{proved}: the action and the Haar measure are invariant
under the hypercubic group. Lean: hypotheses \texttt{hs12}, \texttt{hs23}, \texttt{hs34}.\\
(W2) & For every spatial section $\sigma=(n_1,n_2,n_3)$ there is a self-adjoint, positive, trace-class operator $K_\sigma$ on the space
$\Hinv(\sigma)$ of gauge-invariant $L^2$ functions of the spatial links, with $Z_\beta(\sigma,N)=\Tr K_\sigma^N$ for every $N\ge2$ (Section~\ref{sec:reverse} uses this only for spatial periods and times at least $3$). &
\st{hypothesis}, standard \cite{Lus77}, \cite[\S2.3, Thm.~2.1]{OS78}; sketch in Remark~\ref{rem:W2}. Not formalised. Lean:
\texttt{HypSpecTC}. \emph{Positivity of $K_\sigma$ is indispensable} for Theorem~\ref{thm:R} (and for $\Delta\ge0$): without it $\Delta$ can be negative (Remark~\ref{rem:Rused}).\\
(W2$^+$) & In the configuration-space realisation, the normalised transfer is positivity preserving. After gauge projection, bounded
gauge-invariant slab insertions satisfy $|\Kh_Xf|\le\|X\|_\infty\Kh^r|f|$. & \st{hypothesis}, used only for the torus comparison
(Lemma~\ref{lem:torus}); an additional property of the path-integral realisation, not a consequence of positive semidefiniteness of an abstract
operator. Not formalised.\\
(W3) & The lattice OS form of the infinite cylinder $\Z\times\sigma$ is given by the same normalised operator $\Kh_\sigma=K_\sigma/\lambda_0$:
$B(F,G,ma)=\langle\psi(F),\Kh_\sigma^m\psi(G)\rangle$ with $\psi$ linear and $\psi(\Omega)$ the Perron vector. & \st{hypothesis},
standard (spectral theorem for the compact operator $\Kh_\sigma$ and transfer formula of the cylinder state). Not formalised. Lean:
\texttt{SpecLink}, \texttt{LienTransfert4}, which assume only linear square-summable coordinates $c$, $c(\Omega)=\delta_0$, the mixed identity $B(\Omega,F,0)=c_0(F)$ and the diagonal identities for $B(F,F,0)$ and $B(F,F,na)$ (Remark~\ref{rem:W2}), not the full transfer representation at arbitrary times.\\
(W4) & Along the specified non-trivial filter (sequence, subsequence; a double limit is an iterated limit, not the product filter), the
lattice OS forms $B_k(F,G,t)$ converge to $\mathrm{rpForm}\,R\,S$ for $F,G\in\cA_0\cup\{\Omega\}$ and $t\in\{0,\ell\}$. This includes matching the
lattice reflection convention to the continuum reflected form, including any vanishing one-step shifts. & \st{hypothesis}: the
\emph{existence of the continuum limit} and the identification of the OS forms; open. Lean: \texttt{ConvergesOn}.\\
(W5) & (Square tori.) For the family of square tori actually considered, the difference between its OS forms and those of the corresponding
cylinders tends to zero on the same families and at the same times. & Paper estimate under (W2$^+$) for one diagonal choice
$N_k=2^{j_k}n_k$ (Lemma~\ref{lem:torus}, Corollary~\ref{cor:tori}); for a previously specified family it is a separate obligation. Lean:
\texttt{convergesOn\_of\_proche} \emph{assumes} the vanishing difference.\\
(W6) & The limit admits the Osterwalder--Schrader reconstruction, and $S(\emptyset)=1$. & \st{hypothesis} of ClayYM
(\texttt{OSReconstruction}) \cite{OS73,OS75,ClayYM}; not formalised.\\
\bottomrule
\end{tabular}
\caption{Hypotheses. The paper arguments of Sections~\ref{sec:lattice} and later use periods at least $2$, and Section~\ref{sec:reverse} uses (W2) for spatial periods and times at least $3$ only. The Lean declarations assume a single symmetric function on natural periods, with \texttt{HypSpecTC} for every positive spatial period and every positive integer time. Applying them to Wilson's partition function therefore requires a symmetric extension to periods $1$ satisfying those spectral identities; this extension is an additional assumption here, not a consequence proved in this note. Values at a zero period may be set to zero. The Lean hypotheses on the spectral data are, in places, weaker than the table (\texttt{SpecLink} assumes linear square-summable coordinates $c$, $c(\Omega)=\delta_0$, the mixed identity $B(\Omega,F,0)=c_0(F)$, and the diagonal identities for $B(F,F,0)$ and $B(F,F,\ell)$; it does not assume the full transfer representation at arbitrary times. \texttt{ConvergesOn} records the specified convergence of the lattice forms), but the hypothesis on the function $F$ is stronger. The Lean
statement also assumes that the lattice forms are bilinear (\texttt{IsBilin}), which holds by construction for OS forms. (W2) does
\emph{not} assume that $\lambda_0$ is simple: degenerate eigenvalues appear as $r_i=1$ (Example~\ref{ex:z2}). Not assumed anywhere: uniqueness of the Gibbs state or of the periodic limit, uniqueness of the continuum limit, reflection positivity of
the limit as a separate hypothesis (it follows from (W3)--(W4)), $\beta_k\to\infty$.}
\label{tab:hyp}
\end{table}

Under (W2), let $\lambda_0\ge\lambda_1\ge\lambda_2\ge\dots\ge0$ be the eigenvalues of $K_\sigma$ with multiplicity, $r_i=\lambda_i/\lambda_0\in[0,1]$
($i\ge1$) and $s_\sigma(N)=\sum_{i\ge1}r_i^N$ (finite for $N\ge1$). Then
\begin{equation}\label{eq:FN}
F(\sigma,N)=N\log\lambda_0+\log\bigl(1+s_\sigma(N)\bigr).
\end{equation}

\begin{remark}[on (W2) and (W3)]\label{rem:W2}
For $\beta>0$, the character expansion $e^{(\beta/2)\Tr g}=\sum_jc_j(\beta)\chi_j(g)$ has $c_j(\beta)=2(2j+1)I_{2j+1}(\beta)/\beta>0$ (at
$\beta=0$ one uses the continuous limit); this gives the positive convolution factor of the usual temporal-gauge construction. Multiplication
by the spatial plaquette weights and projection onto gauge-invariant functions (the Polyakov loops inserted by fixing the temporal gauge on
$N-1$ slices) must then be combined with all per-slice scalar factors of \eqref{eq:wilson} to give the precise identity in (W2). Reflection
positivity and the positive transfer operator: \cite{Lus77} is cited for the standard transfer construction, not re-derived here; \cite[\S2.3]{OS78} \st{read}. The present paragraph is a sketch, not a
complete derivation of (W2)--(W3) with our conventions; for Lemma~\ref{lem:torus}, (W2$^+$) must also hold after gauge projection.
The centre transformations on the links crossing a spatial hyperplane commute with $K$, so that $\Hinv=\bigoplus_{e\in\Z_2^3}\mathcal
H_e$ (electric-flux sectors), and the periodic trace contains \emph{all} sectors. (W3) is the spectral theorem for the compact positive
operator $\Kh_\sigma$ on $\Hinv(\sigma)$ (eigenbasis $e_0=\Omega,e_1,\dots$, kernel included), together with the formula
$B(F,G,ma)=\langle\psi(F),\Kh^m\psi(G)\rangle$ for the reflected two-point function of the cylinder state; in coordinates
$c_i(F)=\langle e_i,\psi(F)\rangle$ it reads $B(F,F,0)=\sum_ic_i(F)^2$, $B(F,F,na)=c_0(F)^2+\sum_{i\ge1}r_i^nc_i(F)^2$, $B(\Omega,F,0)=c_0(F)$,
$c(\Omega)=\delta_0$. The spectrum $(r_i)$ is the same as in \eqref{eq:FN}.
\end{remark}

\subsection{Main theorem}
ClayYM \cite{ClayYM} writes Euclidean Schwinger kernels $S$ for the $21$ components of $\Tr F_{\mu\nu}F_{\rho\sigma}$, their
Osterwalder--Schrader form $\mathrm{rpForm}\,R\,S\,(F,G,t)=\langle F,e^{-tH}G\rangle$ on finite families $F,G$ of positive-time test
functions, the truncation $v_A=A-\langle\Omega,A\rangle\Omega$, and
\begin{equation}\label{eq:qgap}
\mathrm{QGap}(R,S,\cA_0,\ell,q):\quad 0<\ell,\ 0<q<1,\ \text{and}\ \langle v_A,e^{-\ell H}v_A\rangle\le q\,\langle v_A,v_A\rangle\ \text{for every }A\in\cA_0.
\end{equation}

\begin{theorem}[assembly, $d=4$]\label{thm:main}
Let $(\beta_k,n_k)_{k}$ be couplings and integers $n_k\ge2$, let $\ell>0$ and $a_k=\ell/n_k$, let $(j_k)$ be any integers $\ge0$, and
let $B_k$ be the lattice OS form of the cylinder $\Z\times(2^{j_k}n_k)^3$ at spacing $a_k$ and coupling $\beta_k$. Assume (W1)--(W3)
for every $k$, and
\begin{equation}\label{eq:unif}
\Delta_{\beta_k}(n_k)\le\eta\quad\text{for every }k,\qquad 0<\eta\le2^{-9}.
\end{equation}
If $B_k\to\mathrm{rpForm}\,R\,S$ along a non-trivial filter in the sense of (W4), with $\cA_0$ the set of all families, then
\begin{enumerate}
\item $\mathrm{QGap}(R,S,\text{all families},\ell,256\eta)$, with the same $q=256\eta\le\frac12$ for every $k$;
\item under (W6), $S$ has a mass gap $-\log(256\eta)/\ell\ge(\log2)/\ell$ (\texttt{HasMassGap} of ClayYM), and every vector fixed by
$e^{-\ell H}$ is a multiple of $\Omega$; if $\eta\le2^{-12}$, the gap is $\ge(\log16)/\ell$.
\end{enumerate}
The same conclusions hold for a convergent family of square-torus OS forms if (W5) holds for that family. Under (W2$^+$),
Corollary~\ref{cor:tori} establishes (W5) for its specified diagonal choice of $N_k=2^{j_k}n_k$; this choice does not, by itself, preserve a
continuum limit previously obtained along another volume sequence.
\end{theorem}

The physical time $\ell$ corresponds exactly to $n_k$ lattice steps. If $j_k\to\infty$, the spatial side $2^{j_k}\ell$ diverges, as
required by the lattice sequences of ClayYM. The closure argument applies along any specified non-trivial filter for which the required
convergence holds. An iterated limit (``volume, then spacing'') is not in general convergence along the product filter and the two must not
be identified: for an iterated limit, apply Lemma~\ref{lem:closure} successively, with the same $q$ at both stages. The conclusion is
compatible with a trivial limit: it proves neither the existence of a continuum limit nor a nonzero excited sector, and
\textbf{non-triviality is not addressed}.

%=====================================================================================================================
\section{The lattice part}\label{sec:lattice}

\begin{lemma}[squaring lemma]\label{lem:squaring}
Let $r_i\in[0,1]$ with $s(N)=\sum_ir_i^N<\infty$ for $N\ge1$, and $p(N)=2\log(1+s(N))-\log(1+s(2N))$. Then
(a)~$0\le\log(1+s(N))\le p(N)$, so that $r_i^N\le s(N)\le e^{p(N)}-1$; (b)~$p(2N)\le2\bigl(e^{p(N)}-1\bigr)^2$, without any smallness
condition. \st{proved} \st{Lean: \texttt{lemme\_fin\_tc}}
\end{lemma}
\begin{proof}
(a) $s(2N)=\sum_i(r_i^N)^2\le\sum_ir_i^N=s(N)$ since $r_i^N\le1$; hence $\log(1+s(2N))\le\log(1+s(N))$ and $p(N)\ge\log(1+s(N))\ge0$.
(b) $p(2N)=2\log(1+s(2N))-\log(1+s(4N))\le2s(2N)\le2s(N)^2\le2(e^{p(N)}-1)^2$, using $\log(1+x)\le x$, $\log(1+s(4N))\ge0$ and
$\sum_i(r_i^N)^2\le(\sum_ir_i^N)^2$.
\end{proof}

By (W1)--(W2) and \eqref{eq:FN}, every $p_\mu(m)$ is the quantity $p(N)$ of Lemma~\ref{lem:squaring} for the transfer operator in
direction $\mu$ on the section formed by the other three periods of $m$, with $N$ the $\mu$-th period (the eigenvalue $\lambda_0$ cancels).
In particular $p_\mu(m)\ge0$.

\begin{lemma}[rectangle identity and decomposition]\label{lem:rect}
For every function $F$ on profiles and $\mu\ne\nu$:
(R)~$p_\mu(\tau_\nu m)=2p_\mu(m)-2p_\nu(m)+p_\nu(\tau_\mu m)$. If $F$ is symmetric (W1), then
\begin{equation}\label{eq:decomp}
\Delta(n)=8a_0(n)+4a_1(n)+2a_2(n)+a_3(n),
\end{equation}
where $a_j(n)$ is $p$ at the profile $((2n)^j,n^{4-j})$, doubling one period $n$.
Under (W1)--(W2), $\Delta(n)\ge0$ and $a_j(n)\le\Delta(n)/2^{3-j}$. \st{proved} \st{Lean: \texttt{decomp\_4d}, which states $\Delta\ge0$, $a_0\le\Delta/8$ and $a_3\le\Delta$; the identities (R) and \eqref{eq:decomp} are steps inside the proofs of \texttt{decomp\_4d} and \texttt{doublement\_fin\_4d}}
\end{lemma}
\begin{proof}
(R): compute $F(\tau_\mu\tau_\nu m)=2F(\tau_\nu m)-p_\mu(\tau_\nu m)=2F(\tau_\mu m)-p_\nu(\tau_\mu m)$ and substitute
$F(\tau_\nu m)=2F(m)-p_\nu(m)$, $F(\tau_\mu m)=2F(m)-p_\mu(m)$. \eqref{eq:decomp}: double the four periods one after the other; the sum
telescopes, and (W1) identifies each intermediate profile.
\end{proof}

For $j=0,\dots,3$ let $G_j$ be $p$ at the profile $(n,(4n)^j,(2n)^{3-j})$ doubling the period $n$, and $A_j$ be $p$ at the profile
$(2n,(4n)^j,(2n)^{3-j})$ doubling one period $2n$ in the same direction. Then:
\begin{enumerate}
\item $G_0=a_3(n)\le\Delta(n)$;
\item $A_j\le f(G_j)$, $f(x)=2(e^x-1)^2$: Lemma~\ref{lem:squaring}(b) in the direction of the period $n$ of $G_j$;
\item $G_{j+1}\le2G_j+A_j$: (R) with $m=(n[\mu],2n[\nu],(4n)^j,(2n)^{2-j})$ gives $G_{j+1}=2G_j-2p_\nu(m)+A_j$, and $p_\nu(m)\ge0$;
\item $\Delta(2n)=8A_0+4A_1+2A_2+A_3$: \eqref{eq:decomp} at level $2n$.
\end{enumerate}

\begin{proposition}[doubling, $d=4$]\label{prop:doubling}
Under (W1)--(W2), if $\Delta(n)\le2^{-9}$ then $\Delta(2n)\le256\,\Delta(n)^2$. \st{proved} \st{Lean: \texttt{doublement\_fin\_4d}, \texttt{chaine\_num\_fin4}; the Lean theorem takes as hypotheses the two consequences of (W2) that the proof uses (positivity of $p$ and the squaring bound, from \texttt{hyp\_of\_specTC})}
\end{proposition}
\begin{proof}
Put $D=\Delta(n)\le2^{-9}$. On $[0,1]$, $e^x-1\le x+x^2$ (Mathlib: $|e^x-1-x|\le x^2$ for $|x|\le1$), so $f(x)\le2x^2(1+x)^2$. By items
1--3, $G_0\le D$, $G_1\le2D+2D^2(1+D)^2\le2.004\,D$, $G_2\le4.024\,D$, $G_3\le8.113\,D$, and by item~4
\begin{align*}
\Delta(2n)&\le8f(G_0)+4f(G_1)+2f(G_2)+f(G_3)\\
&\le(16.064+32.384+65.8+135.9)\,D^2=250.148\,D^2\le250.15\,D^2\le256\,D^2.
\end{align*}
\end{proof}
The coefficients above are rounded upward; the exact rational polynomial expression, evaluated at $D=2^{-9}$, is approximately $250.018080$; with the exact exponential the constant is $247.23$, and its limit
as $D\to0$ is $2^d(2^d-1)=240$ (Table~\ref{tab:const}).

\begin{theorem}[transfer gap on all sections]\label{thm:lattice}
Under (W1)--(W2), if $\Delta(n)\le2^{-9}$ then, for every $k\ge0$, $256\,\Delta(2^kn)\le(256\,\Delta(n))^{2^k}$, and for every $j\ge0$
every excited eigenvalue of the transfer operator on the section $(2^jn)^3$ satisfies
\begin{equation}\label{eq:gap}
(\lambda_i/\lambda_0)^n\le256\,\Delta(n)\le\tfrac12.
\end{equation}
If $\Delta(n)>0$ and $\lambda_i>0$ this gives $-\log(\lambda_i/\lambda_0)\ge-\log(256\Delta(n))/n\ge(\log2)/n$; zero eigenvalues satisfy the ratio
estimate directly, and if $\Delta(n)=0$ every excited ratio is zero. In particular $\lambda_0$ is simple. \st{proved} \st{Lean:
\texttt{iteration256}, \texttt{gap\_4d\_fin}, which states $r_i^n\le256\,\Delta(n)$; the bound $\le\frac12$, the logarithmic form and the
simplicity are immediate consequences not restated in Lean}
\end{theorem}
\begin{proof}
If $d_k:=\Delta(2^kn)\le2^{-9}$ then $256d_k\le\frac12$ and $d_{k+1}\le256d_k^2\le d_k/2$; by induction $256d_k\le(256d_0)^{2^k}$. On the
section $(2^jn)^3$ at time $2^jn$, Lemma~\ref{lem:squaring}(a) and \eqref{eq:decomp} give (with $a_0\le\Delta/8\le2^{-12}\le1$, so that $e^x-1\le2x$ applies) $r_i^{2^jn}\le e^{a_0(2^jn)}-1\le2a_0(2^jn)\le
\Delta(2^jn)/4\le(256\Delta(n))^{2^j}$; take the $2^j$-th root.
\end{proof}

\begin{theorem}[$d=3$]\label{thm:3d}
For the three-dimensional theory under the analogues of (W1)--(W2), $\Delta^{(3)}(n)=4a_0+2a_1+a_2$; if $\Delta^{(3)}(n)\le2^{-7}$ then
$\Delta^{(3)}(2n)\le64\,\Delta^{(3)}(n)^2$ and $(\lambda_i/\lambda_0)^n\le64\,\Delta^{(3)}(n)\le\frac12$ on every section $(2^jn)^2$.
\st{proved} \st{Lean: \texttt{doublement\_fin\_3d}, \texttt{gap\_3d\_fin} (conclusion $r_i^n\le64\,\Delta^{(3)}(n)$)}
\end{theorem}
\begin{proof}
Same chain with $G_0\le D\le2^{-7}$, $G_1\le2.016D$, $G_2\le4.098D$ and $4f(G_0)+2f(G_1)+f(G_2)\le60.68\,D^2\le64D^2$ (the polynomial expression evaluates numerically to approximately $60.655245$); on the section $(2^jn)^2$, $r_i^{2^jn}\le2a_0(2^jn)\le\Delta^{(3)}(2^jn)/2\le(64\Delta^{(3)}(n))^{2^j}$.
\end{proof}

\begin{table}[ht]
\centering\footnotesize
\begin{tabular}{@{}cccccc@{}}
\toprule
$d$ & threshold & constant used & polynomial chain & exact exponential \cite{internal} & $\lim_{D\to0}$ \\
\midrule
$2$ & $2^{-5}$ & $16$ & $13.93$ & $13.22$ & $12$ \\
$3$ & $2^{-7}$ & $64$ & $60.655$ & $59.22$ & $56$ \\
$4$ & $2^{-9}$ & $256$ & $250.018$ & $247.23$ & $240$ \\
$4$ & $2^{-12}$ & $256$ & $241.216$ & --- & $240$ \\
\bottomrule
\end{tabular}
\caption{Constants $C_d$ with $\Delta(2n)\le C_d\Delta(n)^2$ below the threshold. ``Polynomial chain'': $e^x-1\le x+x^2$, computed
with rational arithmetic (\texttt{reproduction/resultats/constantes\_poly.json}). Displayed decimal values are approximations to the computed expressions, not certified outward-rounded bounds; the proof uses the rational bounds $256$ and $64$. ``Exact exponential'': interval arithmetic, internal note of 10 November 2026. For
comparison, \cite[(2.12)]{OAI26} has $44$ (and $64$) at $2^{-10}$ in $d=2$, with the squaring bound $p(2n)\le4p^2$ valid for
$p\le1/4$.}
\label{tab:const}
\end{table}

%=====================================================================================================================
\section{The reverse implication: from the doubling test to a bound at one scale}\label{sec:reverse}

Section~\ref{sec:bb26v5} shows that the one-scale hypothesis (B$^*$) of \cite{BB26v5} implies $\Delta_\beta(2T_0)\le2^{-9}$. This section
proves the reverse implication, at a larger scale. We use the notation of \S\ref{sec:setting}. For a spatial torus $t=(t_1,t_2,t_3)$ with all sides at least $3$ put
$V(t)=t_1t_2t_3$, $s_t(M)=\sum_{i\ge1}r_i^M$ for the excited ratios $r_i=\lambda_i/\lambda_0$ of the transfer operator $K_t$ of (W2), and
\[
h_t(M)=\log\bigl(1+s_t(M)\bigr)=\log\frac{Z(t,M)}{\lambda_0(t)^M}\ \ge0,
\]
which is the logarithm of the relative partition function $\widehat Z_t(M)$ of \S\ref{sec:bb26v5}. Hypothesis (B$^*$) at $T_0$ with constant $\gamma$ reads
$h_t(T_0)\le\gamma V(t)$ for all $t$ with sides in $[2T_0,48T_0)$ and $2\gamma(96T_0)^3\le\frac1{20}$. Throughout this section, (W1) and (W2) are
assumed for all periods at least $3$, and, as everywhere in this note, \emph{positivity} of $K_t$ is part of (W2).

\begin{theorem}[reverse implication]\label{thm:R}
Assume (W1)--(W2). Let $n\ge3$ with $\Delta(n)\le2^{-9}$, put $q=256\,\Delta(n)\le\frac12$ and $m_j=2^jn$.
\begin{enumerate}
\item[(a)] (The condition $\Delta(m_j)\le2^{-9}$ needed to apply Proposition~\ref{prop:doubling} at each base is part of the induction of Step~2.) For every $j\ge0$ and every spatial torus $t$ with $\min_it_i\ge m_j$:
\[
h_t(m_j)\ \le\ \frac{2\,\Delta(m_j)}{31\,m_j^3}\,V(t),\qquad\text{and}\qquad \Delta(m_j)\le\frac{q^{2^j}}{256}.
\]
\item[(b)] (B$^*$) holds at $T_0=16n$ with $\gamma=1/(31\cdot2^{23}T_0^3)$, for all spatial tori with all sides at least $T_0$, hence in particular for those with sides in $[2T_0,48T_0)$. Here
$2\gamma(96T_0)^3=\frac{27}{3968}<\frac1{20}$.
\item[(c)] For every integer $N\ge n$, every excited ratio $r$ of the transfer operator of the cube $(N,N,N)$ satisfies $r^n\le q\le\frac12$; if $r>0$ the gap is at least
$\log(1/q)/n\ge(\log2)/n$.
\item[(d)] For every spatial torus with $\min_it_i\ge n$ and $\max_it_i\le24\min_it_i$, every excited ratio $r$ satisfies $r^n\le\theta:=e^{18/31}-1\approx0.787$, so the gap is at least $-\log\theta/n\approx0.2393/n$.
\end{enumerate}
\st{proved} on paper. \st{Lean}: (a), first inequality (statement \texttt{interpolation}), (b) (\texttt{reciproque\_Bstar}) and the first half of (c) ($r^n\le256\Delta(n)$, \texttt{trou\_cubes}), in the abstract form of Appendix~\ref{app:lean}; not (d), not the logarithmic form of (c), and not the passage from periods $\ge1$ to periods $\ge3$.
\end{theorem}

\begin{proof}
\emph{Step 1: monotonicity of the pressure.} Let $f(x)=\log Z(x)/|x|$ for $x\in\{3,4,\dots\}^4$, $|x|=x_1x_2x_3x_4$. Fix a spatial torus $t$ and write $f(t,M)=f(t_1,t_2,t_3,M)$. By \eqref{eq:FN},
$\log Z(t,M)=M\log\lambda_0+h_t(M)$, so
\[
f(t,M)=\frac{\log\lambda_0(t)}{V(t)}+\frac{h_t(M)}{M\,V(t)}.
\]
Since $0\le r_i\le1$ (this is where positivity of $K_t$ is used), $M\mapsto s_t(M)=\sum_{i\ge1}r_i^M$ is non-negative and non-increasing, hence so is $h_t$; and $M\mapsto1/M$ is positive and decreasing. The product of two non-negative non-increasing functions is non-increasing, so $M\mapsto f(t,M)$ is non-increasing. By (W1) the same holds in each of the four coordinates: \emph{$f$ is non-increasing in each period}, with the other three held fixed, with no divisibility condition. Moreover $0\le h_t(M)\le h_t(3)<\infty$, so $h_t(M)/M\to0$ and $f(t,M)\to\log\lambda_0(t)/V(t)$ as $M\to\infty$.

\emph{Step 2: the recursion at every base.} Since $f(m^4)-f((2m)^4)=\Delta(m)/(16m^4)$ by direct computation, and $(2m)^4$ is reached from $m^4$ by raising the four periods one after the other, Step 1 gives $\Delta(m)\ge0$ for all $m\ge3$. Write $\Delta_j=\Delta(m_j)$. By Proposition~\ref{prop:doubling}, applied at \emph{each} base $m_j$: if $\Delta_j\le2^{-9}$ then $\Delta_{j+1}\le256\Delta_j^2=(256\Delta_j)\Delta_j\le\Delta_j/2$. Since $\Delta_0\le2^{-9}$, induction gives $\Delta_j\le2^{-j}\Delta_0\le2^{-9}$ for all $j$, and $256\Delta_{j+1}\le(256\Delta_j)^2$, whence $256\Delta_j\le q^{2^j}$. This is the second inequality of (a).

\emph{Step 3: telescoping.} Fix $j$ and $J>j$. Since $m_i^4=16^{i-j}m_j^4$ and $\Delta_i\le2^{-(i-j)}\Delta_j$,
\[
f(m_j^4)-f(m_J^4)=\sum_{i=j}^{J-1}\frac{\Delta_i}{16\,m_i^4}\le\frac{\Delta_j}{16\,m_j^4}\sum_{k\ge0}32^{-k}=\frac{2\,\Delta_j}{31\,m_j^4}.
\]

\emph{Step 4: the pressure of a torus against the diagonal.} Let $t$ have $\min_it_i\ge m:=m_j$, and let $M\ge3$ be arbitrary. Choose $J\ge j$ with $t_1,t_2,t_3,M\le m_J$. By Step~1, $f(t,M)\ge f(m_J^4)$, and by Step~3, $f(m_J^4)\ge f(m^4)-2\Delta(m)/(31m^4)$. Hence
$\log\lambda_0(t)/V(t)+h_t(M)/(MV(t))\ge f(m^4)-2\Delta(m)/(31m^4)$ for every $M$; letting $M\to\infty$ (Step~1),
$\log\lambda_0(t)/V(t)\ge f(m^4)-2\Delta(m)/(31m^4)$. On the other hand $f(t,m)\le f(m^4)$ by Step~1, because $t_i\ge m$. Subtracting,
\[
\frac{h_t(m)}{m\,V(t)}=f(t,m)-\frac{\log\lambda_0(t)}{V(t)}\le\frac{2\,\Delta(m)}{31\,m^4},
\]
which is the first inequality of (a). No limit of free energies is interchanged with another, and no bound on the aspect ratio is used.

\emph{(b).} By Step~2, $\Delta_4=\Delta(16n)\le2^{-24}$ (the bounds $2^{-9}\to2^{-10}\to2^{-12}\to2^{-16}\to2^{-24}$, each step being $256u^2$). Apply (a) with $m=16n=T_0$: $h_t(T_0)\le2\cdot2^{-24}V(t)/(31T_0^3)=\gamma V(t)$. And $2\gamma(96T_0)^3=2\cdot96^3/(31\cdot2^{23})=27/3968$.

\emph{(c).} Choose $j$ with $m=m_j\le N<2m$, and put $p=2^j$, so $m=pn$. The cube has $V<8m^3$; by (a), $h\le16\Delta(m)/31\le q^p/496$, with $h=h_{(N,N,N)}(m)$. Every excited ratio satisfies $r^m\le s(m)=e^h-1\le2h\le q^p/248$ (as $h\le1$). Since $r^m=(r^n)^p$, this gives $r^n\le q$.

\emph{(d).} With $a=\min_it_i$ and $m=m_j\le a<2m$, $V(t)\le a\,(24a)^2=576a^3<4608\,m^3$. By (a), $h\le9216\Delta(m)/31\le\frac{36}{31}q^p\le\frac{18}{31}2^{1-p}$, using $q\le\frac12$. For $\alpha=2^{1-p}\in(0,1]$ convexity gives $e^{\alpha x}-1\le\alpha(e^x-1)$, so $s_t(m)\le2^{1-p}\theta$ with $\theta=e^{18/31}-1\in(\frac12,1)$. As $\theta\ge\frac12$, $2^{1-p}\le\theta^{p-1}$, so $r^m\le\theta^p$ and $r^n\le\theta$.
\end{proof}

\begin{remark}[what is used, what is not]\label{rem:Rused}
(i)~\emph{Positivity of $K_t$ is indispensable.} It gives $r_i\in[0,1]$ and hence Step~1. Without it $\Delta(n)$ can be negative, and the statement is empty or false. Example: $Z(x)=\prod_{\mu=1}^4(1+\rho^{x_\mu})^{|x|/x_\mu}$ with $\rho=-r$, $0<r<1$, is symmetric, and for each section $t$ equals $\operatorname{Tr}K_t^M$ for the self-adjoint trace-class operator $K_t=c_t\,(1\oplus\rho)^{\otimes V(t)}$ with a positive constant $c_t$ (all factors with $x_\mu\neq M$ are exponentials in $M$), which has an eigenvalue $-r\lambda_0$. For odd $n$, $\Delta(n)=64n^3\log(1-r^n)-32n^3\log(1+r^{2n})<0$. (ii)~The simplicity of $\lambda_0$ is not assumed: degenerate eigenvalues would be ratios $r_i=1$. It follows from the conclusions: for the cubes of (c), $r_i^n\le q\le\frac12$, and for the tori of (d), $r_i^n\le\theta<1$; for other sections the argument does not assert it. (iii)~The proof assumes (W1) and (W2), with positivity, for all integer periods at least $3$, with no upper cutoff: the recursion of Step~2 and the comparison of Step~4 use arbitrarily large dyadic cubes and rectangles. No parity restriction is imposed. For Wilson's $SU(2)$ theory, positivity of the transfer operator at all spatial sides is the point on which the application rests: Lemma~7.3(vii) of \cite{BB26v5} states $T\ge0$ for $\beta\ge0$; we read this statement, and have not checked its proof or its validity for odd sides. (iv)~The Lean version (Appendix~\ref{app:lean}) asks the spectral representation at all times $d\ge1$ and all sides $\ge1$, as the already published \texttt{HypSpecTC}; the paper argument uses only $d,\ t_i\ge3$ (replace $h(1)$ by $h(3)$ in Step~1).
\end{remark}

\begin{remark}[the factor 16 is not optimal]\label{rem:16}
With $T_0=8n$ the same proof would give $y_0=2\gamma(96T_0)^3=54/31>\frac1{20}$, so $16$ is the first dyadic factor that this proof allows. It is not optimal. On product-of-rings models the smallest ratio $T_0/n$ at which $\Delta(n)\le2^{-9}$ implies (B$^*$) lies, numerically, between $1.17$ and $1.75$ (the minimum depends on $n$ and on the family); this is a minimum over those models and not a bound. For massive free scalar fields at the threshold mass, a sufficient (not necessary) test certifies the ratios $1.85$, $2.05$ and $2.3$ at $n=6,8,12$; these are upper bounds on the smallest admissible ratio, not its value (\path{reproduction/reciproque/}). We do not know the optimal factor.
\end{remark}

\begin{remark}[the finite form (B$^{**}$)]\label{rem:Bss}
Proposition~7.57 of \cite{BB26v5} states that (B$^*$) with $\gamma$ implies the inequality (B$^{**}$): $2\log Z(s,T_0)-\log Z(s,2T_0)\le\gamma V(s)$ for all $s$ in the window, and that (B$^{**}$) with $\gamma$ implies (B$^*$) with $\gamma$ (their $\psi_s(T_0)$ is a defect of the doubling of the time only, for tori $(s,T_0)$ of all shapes in the window; ours doubles the four periods of a cube). By Theorem~\ref{thm:R}(b), $\Delta(n)\le2^{-9}$ gives (B$^*$) at $16n$, hence (B$^{**}$) at $16n$ with $2\gamma=2/(31\cdot2^{23}T_0^3)$: indeed $2h_s(T_0)-h_s(2T_0)\le2h_s(T_0)$ since $h_s\ge0$. Conversely (B$^*$), hence (B$^{**}$) up to the factor $2$ in $\gamma$, implies $\Delta(2T_0)\le2^{-9}$ by \eqref{eq:bb26implies}. So our cubic defect at one size controls their window-wide defect at $16$ times that size, and conversely.
\end{remark}

\begin{remark}[numerical checks, \st{computed}]\label{rem:Rnum}
Nothing in the proof depends on these. For massive free scalar fields on tori, on about $1.6\cdot10^5$ product-of-rings models (and direct sums of two of them) and on a pre-registered list of cubes and elongated tori, the inequality (a) and (B$^*$) at $16n$ were evaluated with $\Delta(n)\le2^{-9}$ computed exactly: no violation; the largest ratio of the left to the right side of (a) is $0.245$ (files in \path{reproduction/reciproque/}). These models are not general transfer matrices.
\end{remark}

\begin{corollary}[bridge to Theorem~F of \cite{BB26v5}; \emph{conditional on hypothesis (W) and Theorem F of \cite{BB26v5} as we read them; not verified by us}]\label{cor:bridge}
Let a continuum family in the sense of \cite[Definition~7.1]{BB26v5} for $G=SU(2)$ satisfy their hypothesis (W). Suppose there are $t_0>0$, $p_0$ and $N_0(p)$ such that for all $p\ge p_0$ and $N\ge N_0(p)$ there is an integer $n=n_{N,p}\ge3$ with $16n\le Nt_0$ and $\Delta_{\beta_{N,p}}(n)\le2^{-9}$, where both partition functions in $\Delta$ are taken at the \emph{same} coupling of the family at $(N,p)$, on periodic tori $n^4$ and $(2n)^4$, and (B$^*$) below is taken at that same physical coupling. \emph{Normalisation:} $\beta_{N,p}$ denotes the coupling of the family in the normalisation of \cite{BB26v5}, in which, for $SU(2)$, the Wilson coupling is $\beta_W=4\beta$ (their introduction, footnote on the unit quaternions); the defect $\Delta$ of this note is taken at $\beta_W=4\beta_{N,p}$. No numerical constant of \cite{BB26v5} that depends on $\beta$ is used in this note, and the statements of versions 1.0 and 1.1 contain no numerical comparison involving their $\beta$. Then, by Theorem~\ref{thm:R}(b), their hypothesis (B$^*$) holds at that coupling with $T_0=16n_{N,p}$, $3\le T_0\le Nt_0$, and $y_0=27/3968\le\frac1{20}$, $\gamma=y_0/(2(96T_0)^3)$; their Theorem~F then gives $\sigma(H)\subset\{0\}\cup[\Delta_{\mathrm{phys}},\infty)$ with $0$ simple and $\Delta_{\mathrm{phys}}\ge0.172/t_0$.
\end{corollary}
\noindent\emph{Comment.} We read in \cite{BB26v5} that Theorem~F asks for (W), $t_0>0$, $y_0\le\frac1{20}$, $p_0$, $N_0(p)$ and, for all $p\ge p_0$, $N\ge N_0(p)$, an integer $T_0$ with $3\le T_0\le Nt_0$ such that (B$^*$) holds at $\beta_{N,p}$ with $\gamma=y_0/(2(96T_0)^3)$; (B$^*$) concerns the tori with sides in $[2T_0,48T_0)$, for which Theorem~\ref{thm:R}(b) gives the bound, and (W) is a double limit that this corollary does not touch. The corollary says nothing about the truth of the hypothesis $\Delta\le2^{-9}$ along a family, in particular for the construction of OpenAI, and it is not a proof of the mass gap.

\medskip\noindent\emph{Consequence.} Under (W1)--(W2), the two one-scale criteria are equivalent up to a change of scale (a factor $32$ for a round trip, (B$^*$) at $T_0$ $\Rightarrow$ (B$^*$) at $32T_0$): (B$^*$) at $T_0$ implies $\Delta(2T_0)\le2^{-9}$ (\eqref{eq:bb26implies}); $\Delta(n)\le2^{-9}$ implies (B$^*$) at $16n$. In particular the conclusion of Theorem~\ref{thm:lattice} extends from the dyadic cubes $(2^jn)^3$ to all integer cubes $N\ge n$ and, with a worse constant, to all spatial tori with $\min_it_i\ge n$ and $\max_it_i\le24\min_it_i$. This does not answer the question left open in \cite{BB26v5} (their data there are a gap plus a bound $X_J\le\frac1{10}$, which is not our hypothesis $\Delta\le2^{-9}$).

%=====================================================================================================================
\section{Assembly}\label{sec:assembly}

\begin{lemma}[$q$-form of each lattice form]\label{lem:qform}
If $B$ has the representation (W3) at time $\ell=na>0$ and $r_i^n\le q<1$ for every $i\ge1$, then $B(v_A,v_A,\ell)\le q\,B(v_A,v_A,0)$ for
every $A$, where $v_A=A-B(\Omega,A,0)\Omega$ (and $q>0$). \st{proved} \st{Lean: \texttt{qGapOn\_of\_specLink}}
\end{lemma}
\begin{proof}
By linearity and $c(\Omega)=\delta_0$, $c_0(v_A)=c_0(A)-B(\Omega,A,0)=0$. Hence $B(v_A,v_A,\ell)=\sum_{i\ge1}r_i^nc_i(v_A)^2\le
q\sum_{i\ge1}c_i(v_A)^2=q\,B(v_A,v_A,0)$.
\end{proof}

\begin{lemma}[closure {\cite[\texttt{QGapOn.of\_tendsto}]{ClayYM}}]\label{lem:closure}
Let $B_k$, $B$ be bilinear forms such that $B_k(F,G,t)\to B(F,G,t)$ along a non-trivial filter for $F,G\in\cA_0\cup\{\Omega\}$ and
$t\in\{0,\ell\}$. If eventually every $B_k$ satisfies the $q$-form on $\cA_0$ at time $\ell$ with the same $q$, so does $B$.
\st{proved} \st{Lean, in ClayYM v1.1}
\end{lemma}
\begin{proof}
By bilinearity $B_k(v^{(k)}_A,v^{(k)}_A,t)=B_k(A,A,t)-c_kB_k(A,\Omega,t)-c_kB_k(\Omega,A,t)+c_k^2B_k(\Omega,\Omega,t)$ with
$c_k=B_k(\Omega,A,0)$: every term converges, and a non-strict inequality passes to the limit.
\end{proof}

\begin{proof}[Proof of Theorem~\ref{thm:main}]
For each $k$, Theorem~\ref{thm:lattice} on the section $(2^{j_k}n_k)^3$ gives $r_i^{n_k}\le256\Delta_{\beta_k}(n_k)\le256\eta\le\frac12$;
Lemma~\ref{lem:qform}, with (W3) at time $\ell=n_ka_k$, gives the $q$-form of $B_k$ with $q=256\eta$; Lemma~\ref{lem:closure} with (W4)
gives $\mathrm{QGap}(R,S,\text{all},\ell,256\eta)$. Under (W6), the $q$-form on a set whose truncated classes are dense in
$\Omega^\perp$ (automatic for all families) gives $\|e^{-\ell H}(1-P_\Omega)\|\le q$, hence the gap $-\log q/\ell$ and the uniqueness of
the vacuum \cite[\texttt{QGap.hasMassGap\_univ}, \texttt{QGap.vacuum\_unique}]{ClayYM}; $-\log(256\eta)\ge\log2$ since $256\eta\le\frac12$.
The Lean statement is \texttt{clay\_trou\_de\_critere\_4d} (Appendix~\ref{app:lean}). The positivity of the limit form is not assumed:
$B(F,F,0)=\lim_k\sum_ic_i^{(k)}(F)^2\ge0$.
\end{proof}

The $d=3$ version (\texttt{clay\_trou\_de\_critere\_3d}): under the three-dimensional counterparts of the hypotheses, $\ell>0$ and $\Delta^{(3)}_{\beta_k}(n_k)\le\eta$ for every $k$ with $0<\eta\le2^{-7}$, one gets $q=64\eta$ and the
gap $-\log(64\eta)/\ell$ (Lean), which is $\ge(\log2)/\ell$ since $64\eta\le\frac12$ (arithmetic not restated in Lean in $d=3$) for a limit $S$ of the three-dimensional theory (the layer \texttt{Gap}/\texttt{GapQ} of ClayYM is
generic in $d$; the Clay problem itself is four-dimensional). In $d=3$, $\beta_k=4/(g_3^2a_k)$ at fixed $g_3^2$ (super-renormalisable).

\begin{lemma}[torus versus cylinder]\label{lem:torus}
Let $N=2^jn$ with $j\ge1$, $\sigma=N^3$, and let $X$ be a bounded gauge-invariant observable supported in $r\le N/2$ consecutive time
slices. Let $\omega_N$ be the periodic torus state and $\omega_\infty$ the Perron-cylinder state on $\Z\times\sigma$. Assume (W1)--(W2),
the insertion property (W2$^+$), and $\Delta(n)\le2^{-9}$. Then
\[
|\omega_N(X)-\omega_\infty(X)|\le2\|X\|_\infty\bigl(e^{\Delta(N/2)}-1\bigr),\qquad \Delta(N/2)\le\tfrac1{256}(256\Delta(n))^{2^{j-1}}.
\]
\st{proved} on paper under (W2$^+$); not in Lean
\end{lemma}
\begin{proof}
Let $\Kh_X$ be the normalised insertion operator across the $r$ slices. By (W2$^+$), on the gauge-invariant configuration-space realisation,
$|\Kh_Xf|\le\|X\|_\infty\Kh^r|f|$; since $\|\Kh^r\|=1$ and $\||f|\|_2=\|f\|_2$, we get $\|\Kh_X\|\le\|X\|_\infty$. This step uses positivity
preservation and insertion domination, in addition to positive semidefiniteness; it is the insertion argument of \cite[proof of Prop.~2.2]{OAI26}.
In the eigenbasis of $\Kh$ on $\Hinv(\sigma)$, with $\mu_0=1$, $\mu_i=r_i$,
\[
\omega_N(X)=\frac{\Tr(\Kh^{N-r}\Kh_X)}{\Tr\Kh^N}=\frac{\langle\Omega,\Kh_X\Omega\rangle+\sum_{i\ge1}\mu_i^{N-r}\langle e_i,\Kh_Xe_i\rangle}{1+s_\sigma(N)},
\qquad\omega_\infty(X)=\langle\Omega,\Kh_X\Omega\rangle,
\]
where $\lambda_0$ is simple by Theorem~\ref{thm:lattice}, so that the cylinder state is the Perron vector state, and the gauge
projector commutes with $\Kh$ and $\Kh_X$ ($X$ is gauge invariant). Hence
$|\omega_N(X)-\omega_\infty(X)|\le\|X\|_\infty(s_\sigma(N-r)+s_\sigma(N))\le2\|X\|_\infty s_\sigma(N/2)$. By Lemma~\ref{lem:squaring}(a),
$s_\sigma(N/2)\le e^{p}-1$ with $p$ the doubling of the period $N/2$ at the profile $(N^3,N/2)$, which is $a_3(N/2)\le\Delta(N/2)$; the
last bound is Theorem~\ref{thm:lattice}.
\end{proof}

\begin{corollary}[a diagonal family of square tori, (W5)]\label{cor:tori}
Assume the hypotheses of Lemma~\ref{lem:torus} at each cutoff. Suppose that the lattice OS forms of a pair of families $(F,G)$ are finite
combinations of expectations of products of at most $D(F)+D(G)$ renormalised plaquette fields, each bounded by a constant $c_k<\infty$ at step
$k$ (the group is compact), the sum of the sup-norm bounds being at most $C(F)C(G)\max(1,c_k)^{D(F)+D(G)}$, and that for fixed $k$ these
bounds do not depend on the subsequently chosen spatial size. Set $N_k=2^{j_k}n_k$ and choose $j_k\ge1$ such that
$2(e^{\Delta_{\beta_k}(N_k/2)}-1)\max(1,c_k)^{2k}\le1/k$ and $N_ka_k\ge k\ell$. Then, for fixed $F,G$ and $t\in\{0,\ell\}$, the difference between the
forms of the tori $N_k^4$ and of the cylinders $\Z\times N_k^3$ is at most $C(F)C(G)/k$ from some $k$ on. Consequently either family of forms
converges to a specified form $B$ if and only if the other does, along this selected diagonal. \st{proved} on paper; the Lean reduction
\texttt{convergesOn\_of\_proche} \emph{assumes} the vanishing difference.
\end{corollary}
\begin{proof}
For fixed $F,G$, eventually $D(F)+D(G)\le2k$ and the fixed physical supports fit within half the temporal period; apply
Lemma~\ref{lem:torus} term by term. The choice of $j_k$ is possible because the upper bound for $\Delta_{\beta_k}(N_k/2)$ decreases doubly
exponentially in $j_k$. This construction does not show that an arbitrarily prescribed $j_k\to\infty$ suffices, nor that changing a previously
specified volume sequence preserves its limit.
\end{proof}

\begin{remark}[the meaning of (W4)]\label{rem:W4}
(W4) is the precise sense of ``the Schwinger functions converge'' used here: pointwise convergence of the OS forms of smeared families at
the two times $0$ and $\ell$, not convergence of kernels. It is the role of hypothesis (E) of \cite{BB26}. One technical point: the
plaquette field of ClayYM (\texttt{plaqObs}, plaquette based at the site $x$) is not exactly covariant under the site reflection of the
lattice (temporal plaquettes are shifted by one step); with the lattice sequences of ClayYM, (W4) therefore also requires that this
$O(a)$ shift disappear in the limit (true for locally uniform convergence, or with the clover field, which is exactly covariant). We
include it in (W4).
\end{remark}

%=====================================================================================================================
\section{The converse and why it needs a finite-volume gap}\label{sec:converse}

For $j=0,\dots,d-1$ let $\sigma_j=((2n)^j,n^{d-1-j})$ and $s_j(n)=\Tr'\Kh_{\sigma_j}^n=\sum_{i\ge1}r_i(\sigma_j)^n$ (excited trace). In
physical units these are the spatial tori with sides in $\{\ell,2\ell\}$, at Euclidean time $\ell$.

\begin{proposition}[bracket]\label{prop:bracket}
Under (W1)--(W2), for every $n$,
\[
\sum_{j=0}^{d-1}2^{d-1-j}\log\bigl(1+s_j(n)\bigr)\ \le\ \Delta_d(n)\ \le\ 2\sum_{j=0}^{d-1}2^{d-1-j}\log\bigl(1+s_j(n)\bigr)\ \le\ 2\sum_{j=0}^{d-1}2^{d-1-j}s_j(n).
\]
Consequently $\Delta_d(n)\le\eta$ implies $s_j(n)\le e^\eta-1$ for every $j$, and $\max_js_j(n)\le\eta/(2(2^d-1))$ implies
$\Delta_d(n)\le\eta$ (the constant is $30$ in $d=4$, $14$ in $d=3$). \st{proved}
\end{proposition}
\begin{proof}
$a_j=2\log(1+s_j(n))-\log(1+s_j(2n))$ with $0\le s_j(2n)\le s_j(n)$, so $\log(1+s_j(n))\le a_j\le2\log(1+s_j(n))$; sum with the weights of
\eqref{eq:decomp}.
\end{proof}
Thus \emph{the scalar criterion is, up to the factor $2(2^d-1)$, the smallness of the excited trace at time $\ell$ on $d$ spatial tori with
sides in $\{\ell,2\ell\}$.}

\begin{corollary}[a sufficient converse estimate]\label{cor:converse}
Fix $\ell>t_0>0$, $m>0$ and $B<\infty$, and let $n_ka_k=\ell$. Suppose that, uniformly in $k$ and on the $d$ sections $\sigma_j(n_k)$:
\begin{itemize}
\item[(G$_{\mathrm{fv}}$)] \emph{finite-volume gap, flux sectors included}: $r_i(\sigma_j)\le e^{-ma_k}$ for every $i\ge1$;
\item[(T)] \emph{trace bound at a time $t_0<\ell$}: $\Tr'\Kh_{\sigma_j}^{n_0}\le B$, with $n_0=\lceil t_0/a_k\rceil$.
\end{itemize}
Then $s_j(n_k)\le Be^{-m(\ell-t_0-a_k)}$ and $\Delta_d(n_k)\le2(2^d-1)\,B\,e^{-m(\ell-t_0-a_k)}$; hence $\Delta_d(n_k)\le\eta$ as soon as the right-hand
side is at most $\eta$. If $a_k\le a_*$, the condition $2(2^d-1)Be^{-m(\ell-t_0-a_*)}\le\eta$ is sufficient uniformly in $k$; if $a_k\to0$, the strict
inequality $2(2^d-1)Be^{-m(\ell-t_0)}<\eta$ suffices for $k$ large. \st{proved}
\end{corollary}
\begin{proof}
Since $t_0<\ell=n_ka_k$, $n_0\le n_k$ and $a_k(n_k-n_0)\ge\ell-t_0-a_k$. Hence $r_i^{n_k}=r_i^{n_0}r_i^{n_k-n_0}\le r_i^{n_0}e^{-ma_k(n_k-n_0)}$; sum over $i$ and
apply Proposition~\ref{prop:bracket}.
\end{proof}
This is a statement at a fixed physical size: if $\ell$ varies, $B$, $m$ and $t_0$ may vary with it, and eventual smallness as $\ell\to\infty$
requires, for instance, $\log B(\ell)-m(\ell-t_0(\ell)-a_*(\ell))\to-\infty$; uniformity in the cutoff alone does not give this.
An independent derivation by another model family (Qwen; see the statement on the use of AI) gives the same statement in the form
$\Delta_4(n)\le30A_0e^{m_0r_0}e^{-m_0n}$ under a uniform gap $(\lambda_2/\lambda_1)^r\le e^{-m_0(r-r_0)}$ and a trace bound $A_0$.

Neither hypothesis can be dropped.
\begin{example}[decoupled copies: (T) is needed]\label{ex:copies}
$N$ decoupled copies of a theory satisfying (W1)--(W2) have the same transfer gap and $\log Z$ multiplied by $N$, so
$\Delta^{(N)}=N\Delta^{(1)}$, unbounded in $N$. (T) says that the weighted number of states of the box below energy $\sim1/t_0$ is bounded
uniformly in the spacing: the convergence, at fixed physical size and temperature $1/t_0$, of the finite-volume free energy normalised by
the vacuum. This is physically reasonable but not proved for Yang--Mills (the same difficulty as (W4), in finite volume). \st{proved}
\end{example}

\begin{example}[$\Z_2$ gauge theory at $\beta=\infty$: a gap of the limit is not enough]\label{ex:z2}
For the four-dimensional $\Z_2$ gauge theory at $\beta=\infty$, only flat configurations contribute; their number on a torus with $V$
sites is $|B^1|\cdot|H^1(T^4;\Z_2)|=2^{V-1}\cdot2^4$. Hence $F(\vec n)=cV+3\log2$ on every rectangular torus, every $p$ equals $3\log2$,
and
\[
\Delta_4(n)=15\times3\log2=45\log2\approx31.19\qquad\text{for every }n.
\]
(W1)--(W2) hold: $K$ is a positive scalar multiple (depending on the spatial volume) of the rank-$8$ projector onto the holonomy sectors $H^1(T^3;\Z_2)$, so that $\Kh=K/\lambda_0$ is that projector. Every local gauge-invariant observable is
constant (all contractible loops equal $1$), so the local limit is the trivial theory: unique vacuum and infinite gap. The criterion
fails at every $n$ although the limit is gapped; it fails because (G$_{\mathrm{fv}}$) fails (seven exactly degenerate sectors) \st{proved}. The
limit is trivial, hence excluded by \texttt{FiniteMass} of ClayYM; the analogous phenomenon at large finite $\beta$ (deconfined gapped
phase with eight quasi-vacua split by tunnelling) is a conjecture here, not proved.
\end{example}

\begin{example}[two-dimensional Ising model below $T_c$]\label{ex:ising}
For the Ising model at $K=0.5$ (two coexisting phases), $\Delta_2(n)\to3\log2=(2^2-1)\log2$ \st{computed}, from Kaufman's exact solution
\cite{Kau49} checked against dense transfer matrices (internal note of 10 November 2026). This is the value given by a phase-coexistence
expansion of periodic partition functions of the type studied by Borgs and Koteck\'y \cite{BK90}, $\Delta_d\to(2^d-1)\log(\#\text{phases})$
(our reading; \cite{BK90} do not state it). The periodic limit is a mixture without unique vacuum, so this is not a counterexample in the
sense of ClayYM, but it shows again that (G$_{\mathrm{fv}}$) fails when two quasi-degenerate states exist in finite volume. The same
example was obtained independently by Qwen.
\end{example}

\begin{remark}[what the criterion encodes]\label{rem:flux}
The periodic trace tests every electric-flux sector $e\in\Z_2^3$ (Remark~\ref{rem:W2}), not only the states visible to local observables in
an infinite-volume limit; the bracket therefore controls both excited energies and their thermal multiplicities on the specified finite
boxes. Example~\ref{ex:z2} shows that a gapped local limit, in a broader class of lattice theories, need not imply this finite-volume
criterion; its local limit is trivial, so it does not establish strictness within the full non-trivial Clay class or within pure $SU(2)$
Yang--Mills theory, and no such strictness theorem is claimed. In a torelon model with energy $\approx\sigma L$ the criterion would also constrain
$\sigma\ell^2$; this interpretation, and the estimate $n\approx10$ ($\ell\sqrt\sigma\approx3.7$) obtained at $\beta=2.2986$ from the string tension
and glueball masses of \cite{AT21} with a free-particle and Nambu--Goto torelon model, are \st{heuristic}, not consequences proved for the
Wilson theory.
\end{remark}

%=====================================================================================================================
\section{Flux sectors and 't~Hooft twists}\label{sec:thooft}

\textbf{Status of this section.} The section is about the electric-flux sectors of the periodic trace (Remark~\ref{rem:W2}). Proposition~\ref{prop:sectors} and Theorem~\ref{thm:twosub} are \st{proved} on paper; they were read line by line by the pilot agent of the project (Claude Fable~5.2) on 12 November 2026. Theorem~\ref{thm:twistall} is \st{proved} on paper and rests on Theorem~\ref{thm:R}. Proposition~\ref{prop:wilsonC}, which provides hypothesis (C) below for Wilson's $SU(2)$ theory, rests on the paper proof of \cite{TBv12}, itself checked by AI models only. Proposition~\ref{prop:scale} is \st{conditional} on Theorem~7.59 of \cite{BB26v5}, which we read and re-did step by step. Seven abstract finite statements are checked by the Lean kernel (Appendix~\ref{app:lean13}); that appendix says what they state and what they do not. No human has checked any of it.

The identity of Proposition~\ref{prop:sectors}(a) is classical ('t~Hooft~\cite{tH79}, not read by us; \cite[(2.9) and App.~I]{TY85}; \cite[(8)--(9)]{Kan}); (d) is a rewriting of it. Part of the material on dyadic cubes was obtained earlier in unpublished internal reports of 11 November 2026 \cite{internal}. What we believe is new here is the exact additive decomposition (d) with the vacuum lemma (c), the extension to all cubes, all large times and tori of bounded aspect ratio through Theorem~\ref{thm:R}, the equivalence of Theorem~\ref{thm:twosub} with the criterion, and the scale bound of Proposition~\ref{prop:scale}; we have not searched the literature exhaustively. \emph{Nothing in this section is proved at weak coupling, and nothing in it gives an area law in infinite volume.}

\subsection{Centre operators and twisted partition functions}\label{subsec:C}
\begin{enumerate}
\item[(C)] For every spatial torus $t$ there are three commuting self-adjoint unitaries $C_1,C_2,C_3$ on $\Hinv(t)$ commuting with $K_t$; for $k\in\Z_2^3$ put $C^k=C_1^{k_1}C_2^{k_2}C_3^{k_3}$ and $Z_k(t,M)=\Tr(C^kK_t^M)$, and assume $Z_k(t,M)\ge0$.
\end{enumerate}
(We call $t$ a spatial torus with all sides at least $3$ when (W2) is used, as in Section~\ref{sec:reverse}.) By Remark~\ref{rem:W2}, the centre transformations of the links crossing a spatial hyperplane commute with $K$; (C) asks that they be implemented by unitaries on $\Hinv$ with non-negative twisted traces.

\begin{proposition}[(C) for Wilson's $SU(2)$ theory; \st{proved} on paper, resting on \cite{TBv12}]\label{prop:wilsonC}
Let $t$ be a spatial torus with all sides at least $2$, $\beta>0$, and $T=(MPC_\beta M)|_{\Hinv}$ the transfer operator of \cite{TBv12} (the operator of the proposition on the transfer representation there), so that $K_t=e^{-6\beta V(t)}T$ satisfies $Z_\beta(t,M)=\Tr K_t^M$ for $M\ge2$ and is positive and trace class. For $i=1,2,3$ let $C_i$ be the multiplication by $-1$ of all links of direction $i$ that cross a fixed plane $x_i=c_i$ (at all transverse positions). Then the $C_i$ are self-adjoint unitaries that commute with each other and with $K_t$, and for $M\ge2$ and $k\in\Z_2^3$, $\Tr(C^kK_t^M)$ is the Wilson partition function (same normalisation as $Z_\beta$) in which the sign of $\frac12\Tr U_p$ is reversed on the temporal plaquettes of the planes $(0,i)$ lying above the cut $i$, on one time layer, for each $i$ with $k_i=1$ (a 't~Hooft twist in these time--space planes; the cuts of the different directions are on a single common time layer). In particular $Z_k\ge0$.
\end{proposition}
\begin{proof}
For one direction this is the proof of items (i)--(ii) of the flux-sector theorem of \cite{TBv12}. The $C_i$ act on disjoint sets of links, so they commute. Each preserves the Haar measure and commutes with the gauge transformations (multiplication by a central element), it preserves every spatial plaquette weight (a plaquette contains $0$ or $2$ links of direction $i$ crossing the cut, at the same $x_i$) and the kernel $c_\beta$ (since $(-U)(-V)^{-1}=UV^{-1}$); hence it commutes with $M$, $P$, $C_\beta$ and $T$. In the layer integral of \cite{TBv12}, replacing $V$ by $C^kV$ in the kernel of the last step changes the sign of $h=\frac12\Tr$ on exactly the temporal plaquettes $(0,i)$ whose link in direction $i$ on layer $0$ lies in the cut, for the $i$ with $k_i=1$, because $h(-g)=-h(g)$, and touches no other plaquette; the Mercer expansion gives $\int s_M(U,C^kU)\,dU=\sum_j\lambda_j^M\langle\varphi_j,C^k\varphi_j\rangle=\Tr(C^kT^M)$. The factor $e^{-6\beta V(t)M}$ is common: the action has $6$ plaquettes per site, so $e^{-S_\beta}=e^{-6\beta V(t)M}\exp(\beta\sum_p\frac12\Tr U_p)$ on the lattice of volume $V(t)M$. $Z_k$ is the integral of a positive weight.
\end{proof}
\emph{Scope.} The proofs of \cite{TBv12} are written there for one long direction and a square fibre $F=(\Z/L)^2$; the positivity proof (multinomial expansion of $\exp(\beta\langle Q(U),Q(V)\rangle)$ in the unit quaternions) uses neither parity nor equality of the sides, and the composition of several commuting cuts is the same substitution; we did not rewrite those proofs for rectangular fibres. This is the point on which the application of Section~\ref{sec:reverse} to Wilson's theory rests (Remark~\ref{rem:Rused}(iii)), and the status of \cite{TBv12} is that of a paper proof checked by models only.

\subsection{Decomposition into sectors}\label{subsec:T1}
\begin{proposition}[sector decomposition; \st{proved}]\label{prop:sectors}
Assume (W2) and (C) for a torus $t$ (omitted from the notation). Let $P_e=\frac18\sum_k(-1)^{k\cdot e}C^k$ for $e\in\Z_2^3$, $\mathcal H_e=P_e\Hinv$, and $\Tr_e(M)=\Tr(P_eK^M)\ge0$. Let $\lambda_0=\|K\|>0$.
\begin{enumerate}
\item[(a)] $Z_k=\sum_e(-1)^{k\cdot e}\Tr_e$, $\Tr_0=\frac18\sum_kZ_k$ and $\sum_{e\ne0}\Tr_e=\frac18\sum_{k\ne0}(Z_0-Z_k)$.
\item[(b)] For $k\ne0$: $1-Z_k/Z_0=2\sum_{e:\,k\cdot e\ \mathrm{odd}}\Tr_e/Z_0\in[0,\min(2D,1)]$, where $D=\sum_{e\ne0}\Tr_e/Z_0=\frac18\sum_{k\ne0}(1-Z_k/Z_0)\in[0,1]$.
\item[(c)] (neutral vacuum) $\lambda_0$ is an eigenvalue of $K|_{\mathcal H_0}$; if $\lambda_0$ is simple, its eigenvector lies in $\mathcal H_0$.
\item[(d)] (exact decomposition) Let $s^0(M)=\Tr_0(K/\lambda_0)^M-1\ge0$ (the excited trace of the neutral sector, one copy of $\lambda_0$ removed, which is legitimate by (c)) and $\varphi(M)=\log\bigl(8Z_0/\sum_kZ_k\bigr)=-\log(1-D)\ge0$. Then, with $h(M)=\log(1+s(M))$ as in Section~\ref{sec:reverse},
\[
h(M)=\log\bigl(1+s^0(M)\bigr)+\varphi(M).
\]
\item[(e)] Let $s^{\mathrm{fl}}(M)=\sum_{e\ne0}\Tr_e(K/\lambda_0)^M$. Then $s^{\mathrm{fl}}\le s$, $D=s^{\mathrm{fl}}/(1+s)$, and $0\le1-Z_k/Z_0\le2s(M)/(1+s(M))\le2s(M)$; for $M'\ge M$, $s(M')\le s(M)\rho^{M'-M}$ with $\rho=\max_{i\ge1}r_i$.
\end{enumerate}
\end{proposition}
\begin{proof}
(a) The $C^k$ commute with $K$ and with each other, so they are simultaneously diagonalisable with $K$, and $\Hinv=\bigoplus_e\mathcal H_e$ with $C^k=(-1)^{k\cdot e}$ on $\mathcal H_e$. Hence $Z_k=\sum_e(-1)^{k\cdot e}\Tr_e$, and $\Tr_e=\frac18\sum_k(-1)^{k\cdot e}Z_k$ because $\sum_k(-1)^{k\cdot e}=8\delta_{e,0}$; the third identity is the sum of the second over $e\ne0$ (using $Z_0=\sum_e\Tr_e$). $\Tr_e\ge0$ since $K\ge0$.

(b) $Z_0-Z_k=\sum_e\bigl(1-(-1)^{k\cdot e}\bigr)\Tr_e=2\sum_{k\cdot e\text{ odd}}\Tr_e$; the sum runs over $e\ne0$, hence is at most $2D\,Z_0$. Since $Z_k\ge0$, $1-Z_k/Z_0\le1$. For $e\ne0$ exactly four of the seven $k\ne0$ have $k\cdot e$ odd; summing the identity over $k\ne0$ gives $\sum_{k\ne0}(1-Z_k/Z_0)=8D$.

(c) $K$ is positive and trace class, so $\lambda_0$ is an eigenvalue and the eigenvalues below $\lambda_0$ are at most $r_*\lambda_0$ with $r_*<1$. Let $g_e$ be the multiplicity of $\lambda_0$ in $\mathcal H_e$; $\sum_eg_e\ge1$. For every $e$, $\Tr_e(M)/\lambda_0^M=g_e+\sum_{\lambda<\lambda_0}(\lambda/\lambda_0)^M\to g_e$ as $M\to\infty$, the tail being at most $r_*^{M-1}\,\Tr(K/\lambda_0)$. Hence $Z_k(M)/\lambda_0^M\to G_k:=\sum_e(-1)^{k\cdot e}g_e$, and $G_k\ge0$ because $Z_k\ge0$. Now $\sum_kG_k=\sum_eg_e\sum_k(-1)^{k\cdot e}=8g_0$. If $g_0=0$, all the $G_k\ge0$ have sum $0$, so $G_0=\sum_eg_e=0$, contradicting $\sum_eg_e\ge1$. Hence $g_0\ge1$. If $\lambda_0$ is simple, $\sum_eg_e=1$ and $g_0=1$.

(d) $\Tr_0=Z_0-\sum_{e\ne0}\Tr_e=(1-D)Z_0$, and $D<1$ since $\Tr_0\ge\lambda_0^M>0$ by (c). Dividing by $\lambda_0^M$: $\Tr(K/\lambda_0)^M=1+s(M)$ (definition of $s$, one copy of $\lambda_0$ removed) and $\Tr_0(K/\lambda_0)^M=1+s^0(M)$, so $(1+s^0)=(1-D)(1+s)$; take logarithms. Finally $\Tr_0=\frac18\sum_kZ_k$ by (a), so $\varphi=-\log(1-D)=\log(8Z_0/\sum_kZ_k)$; $s^0\ge0$ because $\Tr_0(K/\lambda_0)^M\ge1$ by (c).

(e) By (c) one copy of $\lambda_0$ lies in $\mathcal H_0$, so the eigenvalues of $K$ on $\bigoplus_{e\ne0}\mathcal H_e$ form a sub-multiset of $\{\lambda_i\}_{i\ge1}$, and $s^{\mathrm{fl}}\le s$. Dividing $D\,Z_0=\sum_{e\ne0}\Tr_e$ by $\lambda_0^M$ gives $D=s^{\mathrm{fl}}/(1+s)$, and (b) gives the bound on $1-Z_k/Z_0$. Finally $r_i^{M'}=r_i^Mr_i^{M'-M}\le r_i^M\rho^{M'-M}$.
\end{proof}
Neither the simplicity of $\lambda_0$ nor any parity of the periods is used. The additive split (d) does \emph{not} pass term by term to the doubling curvatures: with $p^0(M)=2\log(1+s^0(M))-\log(1+s^0(2M))\ge0$ one has $p(M)=p^0(M)+[2\varphi(M)-\varphi(2M)]$, but the bracket can be negative (take at $M=1$ the neutral eigenvalues $1$ and $0.1$ with multiplicity $20$ and one eigenvalue $1$ in a flux sector: the bracket is $2\log\frac43-\log\frac{2.2}{1.2}\approx-0.031$). This is why Theorem~\ref{thm:twosub} below goes through the bracket of Proposition~\ref{prop:bracket} and not through this identity.

\subsection{The criterion forces twist disorder at every larger scale}\label{subsec:T2}
\begin{theorem}[\st{proved} under (W1), (W2), (C)]\label{thm:twistall}
Assume (W1), (W2) and (C) for all periods at least $3$, let $n\ge3$ with $\Delta(n)\le2^{-9}$, and put $q=256\,\Delta(n)\le\frac12$.
\begin{enumerate}
\item[(a)] For every integer $N\ge n$ let $m=2^jn\le N<2m$. For $t=(N,N,N)$, every $M\ge m$ and every $k\ne0$,
\[
s_t(M)\le\frac{q^{M/n}}{248},\qquad 0\le1-\frac{Z_k(t,M)}{Z_0(t,M)}\le\frac{q^{M/n}}{124};
\]
in particular $0\le1-Z_k/Z_0\le q^{N/n}/124\le2^{-N/n}/124$ on the hypercube $N^4$. If $q>0$, every sector $e\ne0$ of the cube has $-\log(\|K|_{\mathcal H_e}\|/\lambda_0)\ge\log(1/q)/n$.
\item[(b)] If $a=\min_it_i\ge n$, $\max_it_i\le24a$ and $m=2^jn\le a<2m$, then for $M\ge m$: $s_t(M)\le\theta^{M/n}$ and $1-Z_k/Z_0\le2\theta^{M/n}$, with $\theta=e^{18/31}-1\approx0.787$.
\item[(c)] For Wilson's $SU(2)$ theory on $N^4$ (so that (W1) extends to the twisted partition functions by permutations of the axes), the bound of (a) holds for every twist class $e_\mu\wedge b\in H^2(T^4;\Z_2)$, $b$ in the span of the other three axes: the six single-plane twists, the spatial (``magnetic'') ones included, and $22$ of the $63$ non-zero classes in all.
\end{enumerate}
In the normalisation of \cite{TY85,KT02}, $e^{-F_{\mathrm{el}}}=\frac12(1-Z_k/Z_0)$ and $e^{-F_{\mathrm v}}=Z_k/Z_0$ for a single-plane twist; on $N^4$ this reads
\[
F_{\mathrm{el}}\ \ge\ \frac Nn\log\frac1q+\log248\ \ge\ \frac Nn\log2+5.51,\qquad 0\le F_{\mathrm v}\le\frac{q^{N/n}}{123}.
\]
\end{theorem}
\begin{proof}
(a) Apply Theorem~\ref{thm:R}(a) at the base $m=m_j$ to $t=(N,N,N)$, for which $\min_it_i=N\ge m$ and $V(t)=N^3<8m^3$: $h_t(m)\le16\Delta(m)/31$, and $\Delta(m)\le q^p/256$ with $p=2^j=m/n$; so $h_t(m)\le q^p/496\le1$ and $s_t(m)=e^{h_t(m)}-1\le2h_t(m)\le q^p/248$. Each excited ratio satisfies $r^m\le s_t(m)\le q^p$, hence $r^n\le q$ and $\rho\le q^{1/n}$. By Proposition~\ref{prop:sectors}(e), for $M\ge m$, $s_t(M)\le s_t(m)\rho^{M-m}\le q^{p}q^{(M-m)/n}/248=q^{M/n}/248$, and $1-Z_k/Z_0\le2s_t(M)$. The spectrum of $K|_{\mathcal H_e}$, $e\ne0$, is at most $\rho\lambda_0$ (Proposition~\ref{prop:sectors}(e)).

(b) As in the proof of Theorem~\ref{thm:R}(d): $s_t(m)\le2^{1-p}\theta\le\theta^p$ with $p=m/n$, hence $r^n\le\theta$, $\rho\le\theta^{1/n}$ and $s_t(M)\le s_t(m)\rho^{M-m}\le\theta^{M/n}$.

(c) A permutation $\pi$ of the axes sending $\mu$ to the time axis maps $N^4$ to itself, preserves the Haar measure and the Wilson action, and maps the stack of twisted plaquettes of the class $e_\mu\wedge b$ (the twisted weight is the Wilson weight with the signs of that stack reversed) to the stack of the time--space twist $k=\pi(b)\in\Z_2^3\setminus\{0\}$, to which (a) applies with $M=N$. Count: the decomposable non-zero classes of $H^2(T^4;\Z_2)\cong\Lambda^2\Z_2^4$ are the $35$ two-dimensional subspaces of $\Z_2^4$; those that contain a coordinate vector are $6$ (spanned by two coordinate vectors), $12$ (a coordinate vector and the sum of two others) and $4$ (a coordinate vector and the sum of the three others), that is $22$; the other $13$ (for instance $(e_1+e_3)\wedge(e_2+e_3)$) contain none. The remaining $28$ classes have non-zero Pfaffian and are not decomposable.

The free energies: $\frac12(1-Z_k/Z_0)\le q^{N/n}/248$ gives the bound on $F_{\mathrm{el}}$ ($q\le\frac12$, $\log248=5.513$). With $x=1-Z_k/Z_0\le1/124$, $0\le-\log(1-x)\le x/(1-x)\le\frac{124}{123}x\le q^{N/n}/123$.
\end{proof}
\emph{What Theorem~\ref{thm:twistall} does not say.} For a time $M<m(N)$ the bound degrades (use a base $m_i\le M$, with a factor $(N/m_i)^3$). It says nothing on the classes of twist that need a magnetic background on the spatial torus ($41$ of the $63$), because the criterion sees tori without twist only. On dyadic hypercubes a bound with the exponent $2N/n$ instead of $N/n$ follows from the planar reflection inequality of the internal reports \cite{internal}; Theorem~\ref{thm:twistall} adds the non-dyadic sizes, all times $M\ge m$ and the tori of bounded aspect ratio.

\begin{remark}[a length, not an area]\label{rem:length}
The lower bound $F_{\mathrm{el}}\ge(N/n)\log(1/q)+\log248$ grows like a \emph{length}, $N/n$, and not like an area $N^2$. Nothing proved here gives a growth of $F_{\mathrm{el}}$ faster than linear in $N$, and no area law in infinite volume follows from Theorem~\ref{thm:twistall}. The exponent comes from the criterion through the excited trace, which the criterion controls extensively in the volume, not through any property of the transverse size of the box. The conclusion of Theorem~\ref{thm:twistall} fails for the deconfined $\Z_2$ theory and for $U(1)$ (table below), so it is not empty, but it is not a statement of confinement.
\end{remark}

\subsection{Two one-scale conditions}\label{subsec:T3}
For $j=0,\dots,3$ let $\sigma_j=((2n)^j,n^{3-j})$ and $\Sigma_j=\sigma_j\times n$ (time $n$), and define
\[
\varepsilon_1(n)=\max_j\max_{k\ne0}\Bigl(1-\frac{Z_k(\Sigma_j)}{Z_0(\Sigma_j)}\Bigr),\qquad
\varepsilon_2(n)=\max_j\Bigl(\frac{\frac18\sum_kZ_k(\Sigma_j)}{\lambda_0(\sigma_j)^n}-1\Bigr)=\max_js^0_{\sigma_j}(n).
\]
(V1) \emph{one-scale twist disorder}: $\varepsilon_1(n)$ is small; (V2) \emph{one-scale neutral bound}: $\varepsilon_2(n)$ is small.
\begin{theorem}[\st{proved} under (W1), (W2), (C)]\label{thm:twosub}
Let $n\ge3$.
\begin{enumerate}
\item[(a)] $\Delta(n)\le30\Bigl[\varepsilon_2+\dfrac{7\varepsilon_1/8}{1-7\varepsilon_1/8}\Bigr]$, hence $\Delta(n)\le30(\varepsilon_1+\varepsilon_2)$ if $\varepsilon_1\le\frac18$.
\item[(b)] $\varepsilon_2\le e^{\Delta(n)}-1$ and $\varepsilon_1\le2(1-e^{-\Delta(n)})\le2\Delta(n)$. On the cube $\Sigma_0=n^4$, which has the weight $8$, even $s^0\le e^{\Delta/8}-1$ and $1-Z_k/Z_0\le2(1-e^{-\Delta/8})$.
\item[(c)] If $\varepsilon_1+\varepsilon_2\le2^{-9}/30$ then $\Delta(n)\le2^{-9}$; if $\Delta(n)\le2^{-9}$ then $\varepsilon_1\le2^{-8}$ and $\varepsilon_2\le1.001\cdot2^{-9}$.
\item[(d)] With $Z^0(\sigma,M)=\frac18\sum_kZ_k(\sigma,M)$ and $p^0_\sigma(n)=2\log Z^0(\sigma,n)-\log Z^0(\sigma,2n)$: $\log(1+s^0_\sigma(n))\le p^0_\sigma(n)\le2\log(1+s^0_\sigma(n))$. So (V2) is, up to a factor $2$, the smallness of a doubling defect \emph{in the time} of the neutral sector, written with twisted partition functions only.
\end{enumerate}
\end{theorem}
\begin{proof}
(a) By Proposition~\ref{prop:bracket}, $\Delta(n)\le2\sum_j2^{3-j}\log(1+s_{\sigma_j}(n))$, and by Proposition~\ref{prop:sectors}(d) at $M=n$, $\log(1+s_{\sigma_j}(n))=\log(1+s^0_{\sigma_j}(n))+\varphi_j$ with $\log(1+s^0_{\sigma_j})\le s^0_{\sigma_j}\le\varepsilon_2$ and $\varphi_j=-\log(1-D_j)\le D_j/(1-D_j)$. By Proposition~\ref{prop:sectors}(b), $D_j=\frac18\sum_{k\ne0}(1-Z_k/Z_0)\le\frac78\varepsilon_1$, and $x\mapsto x/(1-x)$ is increasing. Since $\sum_j2^{3-j}=15$, the sum is at most $30[\varepsilon_2+(7\varepsilon_1/8)/(1-7\varepsilon_1/8)]$. If $\varepsilon_1\le\frac18$, then $7\varepsilon_1/8\le7/64$ and $(7/8)/(1-7/64)=56/57<1$.

(b) The left half of Proposition~\ref{prop:bracket} and Proposition~\ref{prop:sectors}(d) give $\sum_j2^{3-j}[\log(1+s^0_{\sigma_j})+\varphi_j]\le\Delta(n)$. Both terms are $\ge0$ and every weight is $\ge1$, so $\log(1+s^0_{\sigma_j})\le\Delta$ and $\varphi_j\le\Delta$, that is $D_j\le1-e^{-\Delta}$; Proposition~\ref{prop:sectors}(b) gives $1-Z_k/Z_0\le2D_j$. On the cube the weight is $8$.

(c) Computation: $\varepsilon_1\le2(1-e^{-2^{-9}})\le2^{-8}$ and $(e^{2^{-9}}-1)\cdot2^9=1.00098\ldots$.

(d) Apply Lemma~\ref{lem:squaring}(a) to the positive trace-class block $K|_{\mathcal H_0}$, whose largest eigenvalue is $\lambda_0$ by Proposition~\ref{prop:sectors}(c): $Z^0(\sigma,M)=\lambda_0^M(1+s^0(M))$, $\lambda_0$ cancels in $p^0_\sigma(n)$, and the upper bound holds because $\log(1+s^0(2n))\ge0$.
\end{proof}
Thus the criterion is equivalent, up to constants, to (V1) \emph{and} (V2) at one scale, and Theorem~\ref{thm:twistall} says that together they propagate to all larger scales. The two failure modes are separate:
\begin{center}\footnotesize
\begin{tabular}{@{}>{\raggedright\arraybackslash}p{0.27\textwidth}>{\raggedright\arraybackslash}p{0.2\textwidth}>{\raggedright\arraybackslash}p{0.2\textwidth}>{\raggedright\arraybackslash}p{0.25\textwidth}@{}}
\toprule
model & (V1) & (V2) & criterion\\
\midrule
$\Z_2$, $d=4$, $\beta=\infty$ (Example~\ref{ex:z2}) & fails: $Z_k=0$ for $k\ne0$, $\varepsilon_1=1$ & holds: $s^0=0$ & fails, $\Delta_4=15\log8=45\log2$, all of it from $\varphi=\log8$\\
$U(1)$ with weight $\varphi(z^2)$ \cite[Prop.~11.29]{BB26v5}, as read & holds exactly: $Z_k\equiv Z_0$ & fails (photons) & fails, as for $U(1)$\\
$U(1)$, Coulomb phase (stipulated spectrum, \S\ref{sec:witness}) & fails (light fluxes) & fails (photons) & fails\\
\bottomrule
\end{tabular}
\end{center}
The first line is \st{proved} once (C) is granted for that model (the multiplication of the links by $-1$ across a hyperplane is a symmetry of the rank-$8$ projector $K$ of Example~\ref{ex:z2}, and a non-trivial twist class admits no configuration with all twisted plaquettes satisfied, so $Z_k=0$): then $\Tr_e=Z_0/8$ for every $e$, so each of the eight sectors contains exactly one state, $s^0=0$, $\varepsilon_1=1$ and $\varphi=\log8$. The second line uses only that $U\mapsto-U$ leaves the weight invariant and that the image measure is the $U(1)$ Wilson measure, so that $\Delta$ is that of $U(1)$ (as we read \cite[Prop.~11.29]{BB26v5}); the failure of the criterion for $U(1)$ and the third line are conditional on the stipulated spectrum of Section~\ref{sec:witness}. In particular neither an area law nor (V1) alone implies the criterion, in agreement with \cite[Prop.~11.29]{BB26v5}. In the language of \cite[l.~6681--6686]{BB26v5} (``the hypothesis asks that the twisted sectors be disordered at the scale $\ell/2$''), (V1) is the periodic, partition-function form of such a disorder; their hypothesis (M) concerns Dirichlet boxes uniformly in the boundary condition and their (D) is a Poincar\'e inequality, and no implication between either of them and (V1) or (V2) is proved here.

\subsection{A necessary scale}\label{subsec:T5}
We use the following statement of \cite{BB26v5}, \st{read} and re-done by us, not verified in the sense of a full re-derivation. In their normalisation the Wilson weight on $SU(2)=$ unit quaternions is $\exp(c\sum_p\Re\Tr Q_p)$ with $c=2\beta_{\mathrm{BB}}$, and $\beta_W=4\beta_{\mathrm{BB}}$; with our $\beta=\beta_W$ this is $c=\beta/2$, and $c\ge\frac1{32}$ is $\beta\ge\frac1{16}$. \emph{Theorem~7.59 as read}: let $\epsilon(c)=12(\log c+5.42)/c$; for every spatial torus $t=(\ell_1,\ell_2,\ell_3)$ with sides at least $3$ and $(\ell_1\ell_2)^2\epsilon(c)<3$, the largest excited ratio $r_1(t)$ of their transfer operator satisfies $1-r_1(t)\le2(\ell_1+\ell_2)^2\epsilon/\bigl(1-(\ell_1\ell_2)^2\epsilon/3\bigr)$. Their $r_1$ is our largest excited ratio: both transfer operators have the traces $Z(t,M)$ for all $M\ge3$, and a positive summable family is determined by its power sums of order $M\ge3$ (the signed measure $\sum_i\lambda_i^3\delta_{\lambda_i}-\sum_j\mu_j^3\delta_{\mu_j}$ on a compact interval has all its moments equal to zero, hence vanishes by Weierstrass), so the non-zero spectra coincide with multiplicities. Put $\varepsilon:=\epsilon(\beta/2)=24(\log\beta-\log2+5.42)/\beta\le24(\log\beta+4.7269)/\beta$.

\begin{proposition}[necessary scale; \st{conditional} on Theorem~7.59 of \cite{BB26v5} as read]\label{prop:scale}
Assume (W1)--(W2) for Wilson's $SU(2)$ theory at $\beta\ge\frac1{16}$. If $\Delta_\beta(n)\le2^{-9}$ for some $n\ge3$ then, for every $\kappa\in(0,1)$,
\[
n\ \ge\ \min\Bigl\{\Bigl(\frac{3\kappa}{\varepsilon}\Bigr)^{1/4},\ \Bigl(\frac{(1-\kappa)\log2\,(1-\frac{\log2}6)}{8\varepsilon}\Bigr)^{1/3}\Bigr\}.
\]
With $\kappa=\frac12$ and $L_\beta=\log\beta+4.727$ this gives $n\ge\min\{(\beta/(16L_\beta))^{1/4},(\beta/(627L_\beta))^{1/3}\}$; and $n\ge(1-o(1))\bigl(\beta/(8L_\beta)\bigr)^{1/4}$ as $\beta\to\infty$.
\end{proposition}
\begin{proof}
By Theorem~\ref{thm:lattice} with $j=0$ (or Theorem~\ref{thm:R}(c) with $N=n$), $r_1^n\le256\Delta(n)\le\frac12$ on the cube $n^3$, so $1-r_1\ge1-2^{-1/n}\ge(\log2/n)(1-\log2/(2n))\ge\log2\,(1-\log2/6)/n$ for $n\ge3$ ($1-e^{-x}\ge x-x^2/2$). If $n^4\varepsilon\le3\kappa$, the theorem of \cite{BB26v5} at $\ell_1=\ell_2=n$ gives $1-r_1\le8n^2\varepsilon/(1-\kappa)$, whence $n^3\ge(1-\kappa)\log2\,(1-\log2/6)/(8\varepsilon)$. Otherwise $n^4\varepsilon>3\kappa$. The simplified form: for $\kappa=\frac12$ the two constants are $24/(3\kappa)=16$ and $192/(\frac12\log2(1-\log2/6))=626.35\ldots\le627$, with $\varepsilon$ replaced by its upper bound $24L_\beta/\beta$. For the asymptotics, as $\beta\to\infty$ the second branch ($\asymp(\beta/L_\beta)^{1/3}$) exceeds the first ($\asymp(\beta/L_\beta)^{1/4}$) for any fixed $\kappa$; let $\kappa\to1$ slowly.
\end{proof}
The smallest integer $n\ge3$ compatible with both inequalities of the proof (computed exactly, with $1-2^{-1/n}$ and the exact $\epsilon$, not from the analytic minimum, which is smaller) is $3$ for $\beta\le10^5$, $6$ at $10^6$, $22$ at $10^8$, $75$ at $10^{10}$, $242$ at $10^{12}$ and $2334$ at $10^{16}$ (\path{calculs/constantes_v1_3.py}); the bound is therefore empty for $\beta\le10^5$, polynomial in $\beta$, and far below the scale $n\sim e^{3\pi^2\beta/11}$ suggested by asymptotic freedom (\st{heuristic}); it is consistent with, and much weaker than, that expectation. Composing \cite[Cor.~7.60]{BB26v5} with Theorem~\ref{thm:R}(b) instead would give a bound $38$ times weaker.

\emph{Reading of the theorem of \cite{BB26v5}.} As we read its proof, the test functions are traces of Polyakov loops, odd under the centre flips: the theorem concerns the sectors of electric flux, and says that at a fixed size an electric-flux state becomes degenerate with the vacuum as $\beta\to\infty$. Since $1-Z_k/Z_0\ge2r_e^n/(1+s(n))$ for the sectors $e$ with $k\cdot e$ odd, this contradicts (V1) and (V2) \emph{together} at fixed size. That (V1) alone fails there is proved, independently of \cite{BB26v5}, in Section~\ref{sec:fixed}.

%=====================================================================================================================
\section{The criterion fails at every fixed lattice size as $\beta\to\infty$}\label{sec:fixed}

\textbf{Status of this section.} Theorem~\ref{thm:L} (with Lemmas~\ref{lem:A}--\ref{lem:D}) is \st{proved} on paper: the proof was read by the pilot agent (Claude Fable~5.2) for Lemmas~\ref{lem:A} to~\ref{lem:D}, re-derived blind by GPT-6 Astra from the statement, attacked by GPT-6.1 Sol (who corrected the wording of Lemma~\ref{lem:A} and proved the factor $1/\log\beta$ of Lemma~\ref{lem:C}(c)), and found again by an independent chain of the project. The constants $C$ and $\beta_0$ are \emph{not} explicit. Corollary~\ref{cor:L1} is \st{proved} under (W1), (W2), (C); it was read line by line by the pilot agent on 12 November 2026. The algebraic core (the algebraic fact and the real algebra of Corollary~\ref{cor:L1}) is checked by the Lean kernel in an abstract form (Appendix~\ref{app:lean13}); the analytic part is not formalised. Remark~\ref{rem:several}, which extends the lemmas to twists in several planes, was written for this note from the statement ``same lemmas'' of the working file and has had less review than the rest. No human has checked any of it.

\emph{Relation to earlier work (literature search of 12 November 2026; what was read and what was not).} That twisted boundary conditions admit configurations of zero action (``twist eating'' configurations), given by constant transition matrices with $[\Omega_\mu,\Omega_\nu]$ equal to the twist, is stated for $SU(N)$ on the continuum hypertorus by van~Baal \cite[Thm~4.1 and the lines that follow it]{vB82}, who adds that uniqueness up to constant gauge transformations is ``not hard to show'', and who attributes their discovery on the lattice to Groeneveld and collaborators \cite{GJK81} and, in the continuum, to Ambj{\o}rn and Flyvbjerg \cite{AF80} (neither obtained by us). These ``twist eaters'' and their algebra are reviewed in \cite[Sections~4.2--4.3 and 8.1]{GA98}. For $SU(2)$ on a spatial three-torus in the Hamiltonian formulation, the ``vacuum valley'' of constant abelian potentials and the fact that the potential is quartic in all constant potentials at its origin, $V=-(L^3/4g^2)\Tr([A_i,A_j]^2)$, are described by Koller and van~Baal \cite[Section~1]{KvB86}, who refer to L\"uscher \cite{Lu83} (not obtained by us) and note that a magnetic flux removes these quartic modes. The consequence for the small-volume expansion is stated in van~Baal's review \cite[Section~3.4]{vB00}: with twisted boundary conditions ``the classical vacuum at $A=0$ is isolated'' and the small-volume behaviour is ``a perturbative series in $g^2(L)$, as opposed to $g^{2/3}(L)$ in absence of twist''. That the integral over the constant modes of $SU(2)$ is marginally divergent in dimension four is known for the bosonic Yang--Mills matrix integral \cite{KS98}, \cite[eq.~(2.37)]{AW01}; this is the origin of the factor $\log\beta$ in Lemma~\ref{lem:C}(c). The ingredients of Theorem~\ref{thm:L} are therefore classical, and \emph{we do not claim novelty for the counting of zero modes}. In the texts we could read we did not find the bound $Z_k/Z_0\le C\beta^{-(d-1)/2}$ (with $1/\log\beta$ for $d=4$) for the lattice twist ratio at fixed size, its contrast with $U(1)$ and $\Z_2$, or the consequence for a finite-size criterion. We could not obtain \cite{tH79,AF80,GJK81,GAO83,Lu83,CGJK85} (three searches on 12 November 2026 found no copy that an automated client may fetch); the power counting may appear there, in particular in \cite{GJK81,CGJK85}, and the reader should not take its absence from our reading as evidence of novelty. Numerically, Kov\'acs and Tomboulis \cite{KT00} measured the ratio of the twisted to the untwisted partition function on hypercubic lattices at $\beta=2.3,2.4,2.5$ and found that, in physical units, the three couplings fall on one curve, with a sharp rise around $0.7$~fm and an approach to unity at larger sizes; de~Forcrand and von~Smekal \cite{dFvS02} measured the twisted partition functions across the deconfinement transition. Those are measurements at fixed physical size; Theorem~\ref{thm:L} concerns a fixed lattice size as $\beta\to\infty$, where the physical size shrinks.

\emph{What this section says, and does not say.} It says where the criterion fails: at every fixed lattice size $n$, $\Delta_\beta(n)$ stays above a constant much larger than $2^{-9}$ as $\beta\to\infty$. It is not a mass-gap result, it proves nothing at weak coupling about the continuum limit, and it says nothing about how fast $n$ must grow with $\beta$.

\subsection{Setting}
Let $\Lambda=\prod_{\mu=0}^{d-1}\Z/L_\mu$, $d\ge2$, $L_\mu\ge2$, $V=|\Lambda|$, with links $(x,\mu)$ and plaquettes $p=(x;\mu,\nu)$. For $G=SU(2)$ and signs $\epsilon_p\in\{\pm1\}$ put
\[
S_\epsilon(U)=\sum_p\bigl(1-\epsilon_p\tfrac12\Tr U_p\bigr)\ge0,\qquad Z_\epsilon=\int e^{-\beta S_\epsilon(U)}\,dU
\]
with $dU$ the normalised Haar measure and $\beta=\beta_W=4/g^2$ as in~\eqref{eq:wilson}. $Z_0:=Z_{\epsilon\equiv1}$. A \emph{single-plane twist} is $\epsilon_p=-1$ exactly on the plaquettes $(x;0,1)$ with $(x_0,x_1)$ fixed (all values of the other $d-2$ coordinates, so $\prod_{\mu\ge2}L_\mu$ plaquettes) and $\epsilon_p=+1$ elsewhere; $Z_k$ denotes its partition function and $R=Z_k/Z_0$. (This is the twist of Section~\ref{sec:thooft}, time being direction $0$; by translation invariance the position $(x_0,x_1)$ is immaterial, and we take $x_0=L_0-1$, $x_1=L_1-1$.) The gauge group $\mathcal G=SU(2)^\Lambda$ acts by $(g\cdot U)_{(x,\mu)}=g_xU_{(x,\mu)}g_{x+\mu}^{-1}$ and leaves $S_\epsilon$ invariant. For $U(1)$ and $\Z_2$ (on the cube $L_\mu\equiv L$) the action is $1-\epsilon_p\Re U_p$, same signs.

\begin{theorem}[fixed-size twist ratio]\label{thm:L}
For $SU(2)$ there are $C=C(\Lambda)$ and $\beta_0=\beta_0(\Lambda)$ such that for $\beta\ge\beta_0$ and a single-plane twist
\[
\frac{Z_k}{Z_0}\ \le\ C\,\beta^{-(d-1)/2};\qquad\text{for }d=4:\quad \frac{Z_k}{Z_0}\le C\,\frac{\beta^{-3/2}}{\log\beta}.
\]
On the cube $(\Z/L)^d$, for the same twist, $U(1)$ satisfies $Z_k/Z_0\le C'\beta^{q}e^{-2\beta L^{d-4}}$ (for some $C'$, $q$ depending on $L$, $d$) and $\Z_2$ satisfies $Z_k/Z_0\le2^{(d-1)(V-1)}e^{-2\beta L^{d-2}}$. The constants $C,\beta_0,C'$ are not explicit. \st{proved}
\end{theorem}
\emph{Warning on the form of the bound.} In dimension four the theorem gives $Z_k/Z_0=o(\beta^{-3/2})$; hence no equivalent of the form $C\beta^{-3/2}$ with $C>0$ can hold, and the exponent $(d-1)/2$ is an upper bound on the decay, not proved to be its exact order (the exact order is not claimed; a lower bound by a power of $\beta$ was found by Astra, $Z_k/Z_0\ge c_k\beta^{-3(d-1)V/2}$, so the ratio is squeezed between powers of $\beta$).

\subsection{The lemmas}
\begin{lemma}[zero sets]\label{lem:A}
Let $N_k=\{U:S_\epsilon(U)=0\}$ for the single-plane twist, and $N_0=\{S_0=0\}$ (flat configurations). (a) $N_k$ is the union of $2^{d-2}$ orbits of $\mathcal G$, each with stabiliser $\{\pm1\}$ (the constants), hence a compact smooth submanifold of dimension $3V$. (b) $N_0$ contains a smooth submanifold $N_0^{\mathrm{gen}}$ of dimension $3V-1+d$ (flat configurations whose holonomies commute and are not all central).
\end{lemma}
\begin{proof}
$S_\epsilon(U)=0$ iff $\epsilon_pU_p=1$ for every $p$. \emph{(a)} Open the torus along the hyperplanes $x_\mu=L_\mu-1\to0$. The interior plaquettes are flat (they are not twisted: the twisted plaquettes have $x_0=L_0-1$ and $x_1=L_1-1$, so they contain links crossing the cuts), and the box is simply connected, so a gauge transformation sets to $1$ all links that do not cross a cut. The remaining links are those crossing a cut, $U_{(x,\mu)}$ with $x_\mu=L_\mu-1$. A plaquette $(x;\mu,\nu)$ with $x_\mu=L_\mu-1$ and $x_\nu<L_\nu-1$ contains two such links, at $x$ and $x+e_\nu$, and two links that do not cross a cut; it is untwisted when $(\mu,\nu)\ne(0,1)$ or $x_1<L_1-1$, and then flatness says $U_{(x,\mu)}=U_{(x+e_\nu,\mu)}$. Hence each crossing link is constant along the transverse directions, equal to some $\Omega_\mu\in SU(2)$. The plaquettes at the intersection of cuts are then commutators: $\Omega_0\Omega_1\Omega_0^{-1}\Omega_1^{-1}=-1$ (the twisted ones) and $[\Omega_\mu,\Omega_\nu]=1$ for every other pair. Conversely every such $d$-tuple gives a configuration of $N_k$, and two tuples that are conjugate give gauge-equivalent configurations.

The relation $\Omega_0\Omega_1=-\Omega_1\Omega_0$ gives, taking traces in $\Omega_1\Omega_0\Omega_1^{-1}=-\Omega_0$, $\Tr\Omega_0=0$, and likewise $\Tr\Omega_1=0$: both are pure unit quaternions $u,v$, and $uv=-vu$ iff $u\perp v$. $SO(3)$ acts transitively on orthonormal pairs, so $(\Omega_0,\Omega_1)$ is conjugate to $(i\sigma_1,i\sigma_2)$. For $\mu\ge2$, $\Omega_\mu$ commutes with $i\sigma_1$ and $i\sigma_2$, hence with the whole quaternion algebra, so $\Omega_\mu=\pm1$. This gives $2^{d-2}$ classes. In the gauge where the non-crossing links equal $1$, the residual gauge transformations are the constants $g$ with $g\Omega_\mu g^{-1}=\Omega_\mu$ for all $\mu$, that is $g=\pm1$. The classes are distinct because $\Omega_\mu=\pm1$ ($\mu\ge2$) is gauge invariant. So $N_k$ is a union of $2^{d-2}$ orbits $\mathcal G/\{\pm1\}$, of dimension $3V$.

\emph{(b)} A commuting $d$-tuple with a non-central element is conjugate into the maximal torus $\{e^{i\theta\sigma_3}\}$; near such a point the commuting $d$-tuples form a manifold of dimension $2+d$ (the axis in $S^2$ and the $d$ angles). The gauge-fixed description as in (a) shows that the stabiliser of such a configuration is the centraliser of the tuple, a copy of $U(1)$, so the orbit has dimension $3V-1$, and the family of orbits has the $d$ angles as parameters: $\dim N_0^{\mathrm{gen}}=3V-1+d$.
\end{proof}

\begin{lemma}[Morse--Bott along $N_k$]\label{lem:B}
At every point of $N_k$ the Hessian of $S_\epsilon$ (for the bi-invariant metric $\langle X,Y\rangle=-\frac12\Tr XY$) has as kernel exactly the tangent space of the gauge orbit, and is positive definite on its orthogonal complement.
\end{lemma}
\begin{proof}
Write $U=\exp(A)\bar U$ with $\bar U\in N_k$, $A_e\in\mathfrak{su}(2)$. Each plaquette is $U_p=\epsilon_p\exp(F_p)$ with $F_p=(D_1A)_p+O(|A|^2)$, where $D_1$ is the lattice exterior derivative twisted by the adjoint representation of the flat background $\bar U$, and $1-\epsilon_p\frac12\Tr U_p=1-\cos|F_p|=\frac12|F_p|^2+O(|F_p|^4)$. Hence $S_\epsilon=\frac12\|D_1A\|^2+O(|A|^3)$: the Hessian is $D_1^*D_1$ with no curvature term, because $S_\epsilon$ vanishes with its gradient at $\bar U$. Its kernel is $\ker D_1=\operatorname{im}D_0\oplus H^1(\Lambda;\mathrm{ad}\bar U)$, and $\operatorname{im}D_0$ is the tangent space of the orbit. It remains to show $H^1=0$. Under $\mathrm{Ad}(i\sigma_1)$, $\mathrm{Ad}(i\sigma_2)$ and $\mathrm{Ad}(\pm1)$, $\mathfrak{su}(2)$ splits into three flat real line bundles: $\sigma_1$ (holonomy $-1$ along direction $1$), $\sigma_2$ (holonomy $-1$ along direction $0$), $\sigma_3$ (holonomy $-1$ along directions $0$ and $1$); the holonomy along $\mu\ge2$ is $+1$. The cohomology of a rank-one local system on the torus whose holonomy is not trivial vanishes in all degrees (K\"unneth: $H^*(S^1;\mathcal L)=0$ if the holonomy is $\ne1$), and the cell complex of the lattice torus ($L_\mu\ge2$) computes it. Hence $H^1=0$.
\end{proof}

\begin{lemma}[Laplace]\label{lem:C}
(a) $Z_k\le C_1\beta^{-3(d-1)V/2}$. (b) $Z_0\ge c\,\beta^{-(3(d-1)V+1-d)/2}$. Hence $Z_k/Z_0\le C\beta^{-(d-1)/2}$. (c) For $d=4$, $Z_0\ge c\,\beta^{-(9V-3)/2}\log\beta$, hence $Z_k/Z_0\le C\beta^{-3/2}/\log\beta$.
\end{lemma}
\begin{proof}
(a) By Lemmas~\ref{lem:A} and~\ref{lem:B} and compactness there are a tubular neighbourhood $\mathcal T$ of $N_k$, $c>0$ and $\delta>0$ with $S_\epsilon\ge c\,\mathrm{dist}(\cdot,N_k)^2$ on $\mathcal T$ and $S_\epsilon\ge\delta$ outside. In normal coordinates, $Z_k\le e^{-\beta\delta}+C\,\mathrm{vol}(N_k)(\pi/(c\beta))^{\kappa/2}$ with $\kappa=3dV-3V$ the codimension of $N_k$ in the configuration space of dimension $3dV$; this is $\le C_1\beta^{-3(d-1)V/2}$.

(b) Only a Taylor bound is needed (no non-degeneracy): on a tubular neighbourhood of a compact set $K\subset N_0^{\mathrm{gen}}$ with non-empty interior in one smooth piece, away from the central strata, $S_0\le C'\mathrm{dist}(\cdot,N_0)^2$ because $S_0\ge0$ vanishes to second order on $N_0$. The Haar density is bounded below in normal coordinates, so $Z_0\ge c(K)\beta^{-\kappa'/2}$ with $\kappa'=3dV-(3V-1+d)=3(d-1)V+1-d$. The quotient has the exponent $\bigl[3(d-1)V-(3(d-1)V+1-d)\bigr]/2=(d-1)/2$. (The central stratum of $N_0$ is \emph{not} Morse--Bott; we do not use a Gaussian expansion there.)

(c) (Sol.) In the gauge fixed at a base point, near the trivial field there are $12$ constant coordinates $X_\mu\in\mathfrak{su}(2)$ ($\mu=0,\dots,3$) and $9(V-1)$ further coordinates $q$ (total $12V-3(V-1)=9V+3$), and by Taylor expansion of the plaquette holonomies (smooth, vanishing when $q=0$ and the $X_\mu$ commute) $S_0\le C\bigl(\sum_{\mu<\nu}\|[X_\mu,X_\nu]\|^2+\|q\|^2\bigr)$. Consider the region $X_0=ru$ ($u$ a unit vector), $X_i=t_iu+z_i$ with $z_i\perp u$ ($i=1,2,3$), $\beta^{-1/4}\le r\le r_0$, $|t_i|\le cr$, $|z_i|\le c\beta^{-1/2}/r$, $\|q\|\le c\beta^{-1/2}$. There $\|[X_0,X_i]\|\lesssim r|z_i|\le c^2\beta^{-1/2}$ and $\|[X_i,X_j]\|\lesssim|t_i||z_j|+|t_j||z_i|+|z_i||z_j|\lesssim\beta^{-1/2}$ (using $r\ge\beta^{-1/4}$), so $\beta S_0$ stays bounded on the region. Its volume is, with $X_0=ru$ of density $r^2dr\,du$, $\int r^2dr\prod_{i=1}^3\bigl(r\cdot\beta^{-1}r^{-2}\bigr)\cdot\beta^{-9(V-1)/2}\asymp\beta^{-3-9(V-1)/2}\int_{\beta^{-1/4}}^{r_0}\frac{dr}{r}\asymp\beta^{-(9V-3)/2}\log\beta$. Hence $Z_0\ge c\,\beta^{-(9V-3)/2}\log\beta$, while by (a) with $d=4$ ($3(d-1)V/2=9V/2$) one gets $Z_k/Z_0\le C\beta^{-3/2}/\log\beta$.
\end{proof}

\begin{lemma}[$U(1)$ and $\Z_2$]\label{lem:D}
On the cube $(\Z/L)^d$ with the same twist: for $U(1)$, $\min S_\epsilon\ge2L^{d-4}$ and $Z_0\ge c\beta^{-q}$; for $\Z_2$, $\min S_\epsilon=2L^{d-2}$ and $Z_0\ge2^{-(d-1)(V-1)}$.
\end{lemma}
\begin{proof}
Fix the $d-2$ coordinates other than $0,1$: the $L^2$ plaquettes $(x;0,1)$ of this section form a closed $2$-torus, exactly one of them twisted, and there are $L^{d-2}$ sections. \emph{$U(1)$.} The sum of the plaquette angles $\theta_p$ over a section lies in $2\pi\Z$ (each link appears twice with opposite signs). With $\phi_p=\theta_p+\pi\mathbf 1[p\text{ twisted}]$, $\sum_p\phi_p\equiv\pi\pmod{2\pi}$. Since $1-\cos\phi\ge(2/\pi^2)\mathrm{dist}(\phi,2\pi\Z)^2$, $\sum_p\mathrm{dist}(\phi_p,2\pi\Z)\ge\pi$ and by Cauchy--Schwarz over $L^2$ plaquettes the cost of a section is at least $(2/\pi^2)\pi^2/L^2=2/L^2$; summing over $L^{d-2}$ sections gives $2L^{d-4}$. $Z_0\ge c\beta^{-q}$ by the Taylor argument of Lemma~\ref{lem:C}(b) on the smooth manifold of flat $U(1)$ configurations. \emph{$\Z_2$.} In each section the product of the $U_p$ is $1$, but exactly one plaquette is twisted, so at least one plaquette of the section has $\epsilon_pU_p=-1$ and costs $2$; $U=1$ attains $2L^{d-2}$. The flat configurations are $|B^1|\cdot|H^1(T^d;\Z_2)|=2^{V-1+d}$ out of $2^{dV}$, so $Z_0\ge2^{V-1+d-dV}=2^{-(d-1)(V-1)}$.
\end{proof}

\begin{proof}[Proof of Theorem~\ref{thm:L}]
$SU(2)$: Lemma~\ref{lem:C}(b) and (c). $U(1)$ and $\Z_2$: $Z_k\le e^{-\beta\min S_\epsilon}$ (normalised measure) and Lemma~\ref{lem:D}.
\end{proof}

\begin{remark}[several planes]\label{rem:several}
For a twist in the planes $(0,i)$, $i\in S\subset\{1,\dots,d-1\}$, $S\ne\emptyset$ (all with the same time layer), the same lemmas hold with the same exponent. The relations become $\Omega_0\Omega_i=-\Omega_i\Omega_0$ ($i\in S$) and all other pairs commuting. Then $\Omega_0=u$ and each $\Omega_i$ ($i\in S$) is a pure unit quaternion orthogonal to $u$, and two commuting such elements are equal up to sign, so $\Omega_i=\pm v$ with a fixed $v\perp u$; the other $\Omega_\mu$ are $\pm1$. The stabiliser is again $\{\pm1\}$, the zero set is a union of orbits $\mathcal G/\{\pm1\}$ of dimension $3V$, and the three line bundles of Lemma~\ref{lem:B} have holonomies, along the direction $0$ and along $i\in S$, $(+,-)$, $(-,+)$ and $(-,-)$ respectively, never all trivial since $S\ne\emptyset$, so $H^1=0$. Lemma~\ref{lem:C} then gives $Z_k/Z_0\le C\beta^{-(d-1)/2}$ for every $k\ne0$ (the denominator $Z_0$ is the same). This remark is used for the $7$ twists $k\in\Z_2^3\setminus\{0\}$ of Corollary~\ref{cor:L1}(ii)--(iii) and for the rectangular tori $\sigma_j\times n$; the proofs of the lemmas are written for general sides $L_\mu\ge2$ and use the shape only through $V$.
\end{remark}

\subsection{The algebraic fact}
Two elements of $SU(2)$ can anticommute ($i\sigma_1$ and $i\sigma_2$); in an abelian group, or in any commutative ring of characteristic $\ne2$ without zero divisors, two non-zero elements never do. A twist is therefore classically free for $SU(2)$ (there are flat $SO(3)$ bundles of the twist class considered; not every class $w_2\ne0$ has one, for instance $n_{01}=n_{23}=1$, whose Pfaffian is non-zero) and costs an action of order $\beta$ for $U(1)$ and $\Z_2$ (Lemma~\ref{lem:D}). This is the source of the difference of Theorem~\ref{thm:L}: for $SU(2)$ the suppression of the twist is only entropic (isolated twisted vacua against a moduli space of untwisted flat connections of positive dimension), a power of $\beta$; for $U(1)$ and $\Z_2$ it is energetic, $e^{-c\beta}$. Lean checks the algebraic statements in Appendix~\ref{app:lean13} (anticommuting pair of $SU(2)$, the commutant, an explicit pair, and the abelian cases); the lattice geometry is not formalised.

\subsection{Consequence for the criterion}
\begin{corollary}[the criterion fails at fixed size; \st{proved} under (W1), (W2), (C)]\label{cor:L1}
Let $n\ge3$ and assume (W1), (W2), (C) for the tori $\sigma_j\times n$, $j=0,\dots,3$ (for Wilson's $SU(2)$ theory, (W1) is proved and (W2), (C) are given by Proposition~\ref{prop:wilsonC}). Then:
\begin{enumerate}
\item[(i)] There are $C(n)$ and $\beta_1(n)$ with, for $\beta\ge\beta_1(n)$, $\Delta_\beta(n)\ge8\log\dfrac{8}{5+3C(n)\beta^{-3/2}}>3.7$.
\item[(ii)] $\displaystyle\liminf_{\beta\to\infty}\Delta_\beta(n)\ \ge\ 15\log8=45\log2\approx31.19$, the number $\Delta_4(n)=45\log2$ computed for the deconfined $\Z_2$ theory in Example~\ref{ex:z2} (the $SU(2)$ statement is an inequality, not an equality).
\item[(iii)] On the cube alone, with the seven twists: $\Delta_\beta(n)\ge8\log\dfrac{8}{1+7r}$ whenever the seven ratios $Z_k/Z_0$ ($k\ne0$) of the torus $n^4$ are at most $r$; in particular $\Delta_\beta(n)\ge8\log4\approx11.09$ as soon as $r\le1/7$.
\end{enumerate}
These are lower bounds, not limit values. They imply that $\varepsilon_1(n)\to1$ as $\beta\to\infty$ at every fixed $n$, so (V1) of Section~\ref{sec:thooft} fails there.
\end{corollary}
\begin{proof}
By Theorem~\ref{thm:twosub}'s ingredients (the left half of Proposition~\ref{prop:bracket} and Proposition~\ref{prop:sectors}(d)), $\sum_j2^{3-j}\bigl[\log(1+s^0_{\sigma_j}(n))+\varphi_{\sigma_j}(n)\bigr]\le\Delta(n)$, all terms being non-negative, so $\Delta(n)\ge\sum_j2^{3-j}\varphi_j$ with $\varphi_j=-\log(1-D_j)$ and $D_j=\frac18\sum_{k\ne0}(1-R_k)$, $R_k=Z_k/Z_0$ on $\Sigma_j$.

(i) On the cube ($j=0$, weight $8$) the three single-plane twists $(0,i)$ have the same ratio $R\le C(n)\beta^{-3/2}$ by Theorem~\ref{thm:L} and hypercubic symmetry; the four other terms of $D_0$ are $\ge0$ (Proposition~\ref{prop:sectors}(b): $0\le1-Z_k/Z_0$), so $D_0\ge\frac38(1-R)$, $1-D_0\le(5+3R)/8$, and $\Delta\ge8\varphi_0\ge8\log(8/(5+3R))$. For $\beta$ large $8\log(8/5)=3.760\ldots>3.7$.

(ii) For each $j$ the seven twists $k\ne0$ have ratios $R_k\le C_j\beta^{-3/2}$ (Remark~\ref{rem:several}), so $D_j\to7/8$ and $\varphi_j\to\log8$; and $\sum_j2^{3-j}=15$.

(iii) $1-D=(1+\sum_{k\ne0}R_k)/8\le(1+7r)/8$.
\end{proof}
Corollary~\ref{cor:L1} is independent of Theorem~7.59 of \cite{BB26v5} and of Proposition~\ref{prop:scale}, and it passes through (V1) alone; but it is qualitative in $n$ (it gives nothing on how large $n$ must be for a given $\beta$), whereas Proposition~\ref{prop:scale} is quantitative and conditional. That the gap of the transfer operator closes at fixed size as $\beta\to\infty$ is, as we read it, already in \cite{BB26v5}; Corollary~\ref{cor:L1} does not allow us to claim it.

\begin{remark}[comparison with $U(1)$ and $\Z_2$; physical reading]\label{rem:L-comparison}
In Theorem~\ref{thm:L} the twist ratio of $SU(2)$ is bounded by a power of $\beta$, that of $U(1)$ and $\Z_2$ by $e^{-c\beta}$. The only difference between $SU(2)$ and the abelian theories for (V1) at fixed size is therefore the nature of the deficit (a power of $\beta$ against an exponential). A proof of (V1) at weak coupling must see this non-abelian feature, and must show that the polynomial deficit disappears when the effective coupling at the scale $n$ becomes of order one as $n$ grows; nothing in this note does so. \st{heuristic}: a box of fixed lattice size is, at weak coupling, smaller than the deconfinement length in all four directions, and sees the four isolated twisted vacua of Lemma~\ref{lem:A} as degenerate with the untwisted one up to a power of $\beta$.
\end{remark}

%=====================================================================================================================
\section{Witnesses}\label{sec:witness}

\begin{proposition}[universal lower bound]\label{prop:lower}
Under (W1)--(W2), if the section $n^{d-1}$ carries $M$ excited states with $(\lambda/\lambda_0)^n\ge e^{-c}$, then
$\Delta_d(n)\ge2^{d-1}a_0(n)\ge2^{d-1}\log(1+Me^{-c})$. \st{proved}
\end{proposition}
\begin{proof}
Lemma~\ref{lem:squaring}(a) and \eqref{eq:decomp}.
\end{proof}

\paragraph{A photon spectral witness, conditional on its model identification.} Consider the transfer spectrum consisting of two photon
oscillators for each nonzero spatial momentum, three electric rotors (energies $k^2L_i^2/(2\beta V_s)$ per step) and three magnetic-flux energies
$2\pi^2\beta m^2L_l/(L_iL_j)$, $\beta>0$. This is the spectrum of the $U(1)$ Villain model \emph{without monopoles} (Gaussian gauge field with
compact holonomies, all integer fluxes) if the Euclidean partition function of that model equals $\Tr K^N$ for it with a trace constant equal
to zero, so that (W1)--(W2) hold. \emph{We do not derive this identification} from a fully specified Euclidean measure: the treatment of compact
holonomies, harmonic fluxes, gauge-orbit factors and the absolute trace normalisation is supplied only in a summary (Poisson summation over the
three temporal fluxes). A numerical check---the Euclidean formula and the transfer formula agree to $7\cdot10^{-16}$ relative on $12$ four-dimensional
tori, and the method reproduces the exact two-dimensional Villain partition function to $2\cdot10^{-41}$ \st{computed}---is a consistency
check, not a derivation. Given the identification, the lowest photon frequency on the section $n^3$ is $\omega_1=2\operatorname{arcsinh}(\sin(\pi/n))\le
2\pi/n$, independent of $\beta$; for $n\ge3$ the $6$ lowest nonzero momenta are distinct and, with $2$ polarisations, give twelve one-photon
states, so that Proposition~\ref{prop:lower} gives, for every $n\ge3$ and every $\beta>0$,
\[
\Delta_4(n)\ge8\log(1+12e^{-2\pi})=0.1773\approx90\times2^{-9}\qquad\text{\st{proved}, given the spectrum}.
\]
The exact values (Table~\ref{tab:maxwell}) are numerical; they are much larger, minimal at the self-dual point $\beta=1/(2\pi)$ (electric--magnetic
duality $\beta\leftrightarrow1/(4\pi^2\beta)$), and grow at both ends (quasi-degenerate electric fluxes at large $\beta$, magnetic fluxes at small
$\beta$). Under the stated identification, Maxwell theory never satisfies the criterion. The compact boson without vortices has the same
qualification ($\Delta_4\ge8\log(1+6e^{-2\pi})=0.0891$; numerically $\Delta_4\to25.9$--$42.4$).

\begin{table}[ht]
\centering\footnotesize
\begin{tabular}{@{}lcccccc@{}}
\toprule
$\beta$ & $0.02$ & $0.1$ & $1/(2\pi)$ & $0.5$ & $1$ & $5$\\
\midrule
$\Delta_4(4)$ & $77.4$ & $42.4$ & $38.2$ & $56.5$ & $72.1$ & $108.3$\\
$\Delta_4(32)$ & $72.2$ & $37.2$ & $33.0$ & $51.3$ & $66.9$ & $103.1$\\
\bottomrule
\end{tabular}
\caption{Numerical evaluation of the stipulated photon, rotor and flux spectrum, $\Delta_4(n;\beta)$ \st{computed}; the limit is reached to
$10^{-2}$ at $n=16$. Its identification with a fully normalised monopole-free Villain measure is an additional unproved input, and no uniform
error certificate is asserted.}
\label{tab:maxwell}
\end{table}

\paragraph{Massive free scalar at fixed physical size.} With $x=M\ell$, lattice mass $x/n$, and ``continuum'' values from the spectrum
$\sqrt{M^2+k^2}$ on the tori of sides $\ell$ and $2\ell$ \st{computed}:
\begin{center}\footnotesize
\begin{tabular}{@{}lccccc@{}}
\toprule
$x=M\ell$ & $0.1$ & $1$ & $5$ & $10$ & $15$\\
\midrule
$\Delta_4$, $n=8$ & $48.4$ & $14.7$ & $0.653$ & $0.0204$ & $9.8\cdot10^{-4}$\\
$\Delta_4$, $n=64$ & $48.0$ & $14.27$ & $0.506$ & $0.00724$ & $8.6\cdot10^{-5}$\\
$\Delta_4$, continuum & $48.0$ & $14.26$ & $0.504$ & $0.00711$ & $8.2\cdot10^{-5}$\\
$\Delta_3$, continuum & $21.9$ & $6.21$ & $0.175$ & $0.00194$ & $1.9\cdot10^{-5}$\\
\bottomrule
\end{tabular}
\end{center}
Numerically, the massive scalar passes: $\Delta_4^{\mathrm{cont}}(x)\le2^{-9}$ for $x\ge x_4^*=11.47$ ($2^{-12}$: $13.79$), and
$\Delta_3^{\mathrm{cont}}\le2^{-7}$ for $x\ge x_3^*=8.47$. These crossings are numerical: the spectrum is truncated (frequencies below $80$ in the
script \path{reproduction/scripts/temoins.py}) and no tail certificate is given. At $x=x_4^*$ Theorem~\ref{thm:lattice} gives a gap $\ge(\log2)/\ell=M/16.5$: the
bound is a minorant, loose by a factor $\approx16$. In the massless limit $\Delta^{\mathrm{cont}}(x)\approx-(2^d-1)\log x+\mathrm{const}$
(from $x=0.1$ to $0.01$ the increase is $34.53=15\ln10$ in $d=4$ and $16.11=7\ln10$ in $d=3$).

\emph{Not treated:} the compact $U(1)$ theory (Villain or Wilson, with monopoles) in its four-dimensional Coulomb phase, and
weakly coupled $SU(2)$ in five dimensions; the criterion is expected to fail there for the same reason (massless photons), but we have
no computation.

\paragraph{A non-admissible witness.} The massless \emph{non-compact} scalar violates (W2) (its spatial zero mode has continuous
spectrum). With the convention ``four-dimensional zero mode removed'', $\Delta_4(n)=10.46,\ 1.81,\ -4.74,\ -12.5,\ -17.3,\ -23.6,\
-28.0$ for $n=2,3,4,6,8,12,16$ \st{computed}: negative values, impossible under (W1)--(W2) since $\Delta$ is a positive combination of
$p$'s. Prescriptions that remove zero modes therefore test models outside the hypotheses.

\paragraph{Dimension three.} For the same kind of Gaussian transfer spectrum, assuming its identification with the periodic partition function,
the four distinct lowest photon momenta give, for $n\ge3$ and $\beta>0$, $\Delta_3(n)\ge4\log(1+4e^{-2\pi})=0.0298\approx3.8\times2^{-7}$ \st{proved}
given the spectrum; numerically $\Delta_3\ge6.46$ (minimum at $e^2\ell\approx8.0$, $n=32$) \st{computed}. This Gaussian calculation is not a
calculation of the compact three-dimensional $U(1)$ theory with monopoles; reference \cite{GM82} concerns the confinement of static charges, and
no further spectral statement from that paper is asserted here.

%=====================================================================================================================
\section{Strong coupling and a reported computational instance}\label{sec:strong}

\paragraph{Thermodynamic integration.} Since $\partial_\beta F_\beta(\vec n)=-\sum_p\langle1-\frac12\Tr U_p\rangle=-6|\Lambda|(1-\langle
P\rangle_{\vec n})$ with $\langle P\rangle$ the mean plaquette and $F_0=0$,
\begin{equation}\label{eq:thermo}
\Delta_\beta(n)=96\,n^4\int_0^\beta\bigl(\langle P\rangle_{n^4}-\langle P\rangle_{(2n)^4}\bigr)\,d\beta'\qquad\text{\st{proved}}.
\end{equation}

\paragraph{Leading terms.} With $t=I_2(\beta)/I_1(\beta)\approx\beta/4$, the $6n^2$ flat wrapping sheets and the $12n^3$ unit-cross-section wrapping
tubes have fundamental weights $t^{n^2}$ and $t^{4n}$ (a sheet is a full $2$-torus of $n^2$ plaquettes with Euler characteristic $0$; a tube has $4n$ plaquettes). Their
contributions to $16\log Z(n^4)$ are respectively $96\,n^2t^{n^2}$ and $192\,n^3t^{4n}$. At $n=2$ the sheets contribute $384\,t^4$ and the tubes contribute $1536\,t^8$;
the latter is not the complete coefficient of $t^8$ in $\Delta_4(2)$, which is not determined here. At $n=4$ these two families together contribute $13\,824\,t^{16}$.
Establishing the complete Taylor coefficients and the stated remainder orders requires the classification of the remaining graphs, representations and logarithmic
cluster contributions; that derivation is not supplied here. A Monte Carlo computation of \eqref{eq:thermo} on $2^4$ and $4^4$ for $\beta\le1.5$ is consistent with the
leading coefficient $384$ to $2.5\%$ (fitted $A=0.997\pm0.025$ in $\Delta_4(2)/(384t^4)=A+O(t^4)$) \st{computed, internal note of 10 November 2026}; it does not
determine a Taylor coefficient. At leading order, $n=2$ would suffice for $\beta\le0.190$ and $n=4$ up to $\beta\approx1.65$; but the leading term underestimates
$\Delta_4(2)$ by $14\%$ at $\beta=1$ and $72\%$ at $\beta=1.5$, so these values are \st{heuristic}. A Monte Carlo estimate never certifies the criterion: at weak coupling it
would require the plaquette to a relative accuracy of order $10^{-13}$.

\paragraph{A reported computational instance.} Assuming the two computational inputs specified below, the following analytic argument gives an interval on which the scalar hypothesis holds for the $SU(2)$ Wilson partition function.

\begin{proposition}[reported computational instance, $n=2$]\label{prop:instance}
Assume the reported finite combinatorial input (ENC): the complete counts of connected closed plaquette supports through size six, as listed in the proof. Assume also
the exact rational comparison (RAT) specified after the proof. Then, for the four-dimensional $SU(2)$ Wilson theory and $0\le\beta\le\beta_*=\BETASTAR$,
$\Delta_\beta(2)\le\DELTABOUND<2^{-9}$. Under (W1)--(W2), the transfer-ratio consequence follows from Theorem~\ref{thm:lattice}: on every spatial torus $(2^{j+1})^3$,
$(\lambda_i/\lambda_0)^2\le256\,\Delta_\beta(2)$ for $\beta\le\beta_*$. The analytic estimate is proved below conditionally on (ENC) and (RAT); these computational inputs
are reported from one implementation and are not independently certified in this note. \st{proved} conditionally on \st{computed} inputs
\end{proposition}
\begin{proof}
Write $e^{(\beta/2)\Tr U}=c_0(1+f(U))$ with $c_0=\int e^{(\beta/2)\Tr U}dU$. Since the tori $2^4$ and $4^4$ have $96$ and $1536=16\cdot96$
plaquettes, the constants $c_0$ and $e^{-\beta N_p}$ cancel in $\Delta$, and $\Delta_\beta(2)=16\log\widehat Z(2^4)-\log\widehat Z(4^4)$,
with $\widehat Z(\Lambda)=\sum_{P}w(P)$ over sets $P$ of plaquettes and $w(P)=\int\prod_{p\in P}f(U_p)\,dU$.

\emph{Activity.} Under Haar measure $X=\frac12\Tr U$ has a symmetric law on $[-1,1]$, and $(\beta/2)\Tr U=\beta X$. Jensen's inequality
gives $c_0=\mathbb E\,e^{\beta X}\ge1$, and the convexity bound $e^{\beta x}\le\cosh\beta+x\sinh\beta$ on $[-1,1]$ gives $c_0\le\cosh\beta$.
Hence $f\le e^\beta/c_0-1\le e^\beta-1$ and $f\ge e^{-\beta}/c_0-1\ge e^{-\beta}/\cosh\beta-1=-\tanh\beta\ge-\beta$, so that
$|f|\le\varepsilon:=e^\beta-1\le\beta+\beta^2$ for $0\le\beta\le1$, and $|w(P)|\le\varepsilon^{|P|}$.

\emph{Closed sets.} If a link of $P$ belongs to only one plaquette $p\in P$, integrating that link first gives $\int f(U_lA)\,dU_l=\int
f=0$ by invariance of the Haar measure (the four links of $p$ are distinct). Hence $w(P)=0$ unless $P$ is
\emph{closed} (every link of $P$ in at least two plaquettes of $P$). A closed $P$ is the disjoint union, in links, of its connected
components $\gamma$ (connected through shared links), each closed, and $w(P)=\prod_\gamma w(\gamma)$. Therefore, with
$S_\Lambda=\sum_\gamma\varepsilon^{|\gamma|}$ over connected closed sets,
\[
\sum_{P\ne\emptyset}|w(P)|\le\sum_{r\ge1}\frac{S_\Lambda^r}{r!}=e^{S_\Lambda}-1,\qquad\log\widehat Z(2^4)\le S_{2^4},\qquad
\widehat Z(4^4)\ge2-e^{S_{4^4}}.
\]
\emph{Counting.} Let $N_m(\Lambda)$ be the number of connected closed sets of $m$ plaquettes. By translation and axis permutations every
plaquette is equivalent to a fixed root, so $N_m=|\mathcal P|\,c_m/m$ with $c_m$ the number of such sets containing the root. An exact
enumeration (Redelmeier's algorithm on the plaquette adjacency graph, with a pruning that is exact for closed sets; \S\ref{sec:repro})
gives \COUNTS. For larger $m$ we use, on $2^4$, $N_m\le\binom{96}{m}$, and on $4^4$, $N_m\le\frac{1536}{m}\binom{20m}{m-1}\le
\frac{1536}{m}\,e\,(20e)^{m-1}$: the adjacency graph has degree $20$, and a connected set of size $m$ containing the root is determined
by the $0/1$ string, of length at most $20m$ with $m-1$ ones, that records which neighbour slots are new members in a breadth-first
exploration. The numbers listed above, including the absence of all other supports of size at most six, are the input (ENC), not a classification proved on paper here. ``Closed'' means that every incident link belongs to at least two selected plaquettes, allowing branched supports.

\emph{Conclusion.} With $y=e^{S_{4^4}}-1\le S_{4^4}+S_{4^4}^2$ and $-\log(1-y)\le y/(1-y)$,
\[
\Delta_\beta(2)\le16\,S_{2^4}+\frac{y}{1-y},
\]
and, for $0\le\beta\le1/200$, $\varepsilon\le u$. The counting bounds give $S_{2^4}\le A$ and $S_{4^4}\le B$, with $A$ and $B$ defined in (RAT). Assuming (RAT), $B\le1$ and $e^{S_{4^4}}-1\le B+B^2=Y<1$. Thus $\Delta_\beta(2)\le16A+Y/(1-Y)=R\le1741/10^6<1/512$. Under (W1)--(W2), $\Delta\ge0$ by Lemma~\ref{lem:rect}.
\end{proof}
\noindent\textbf{(RAT).} With $u=201/40000$ ($=\beta+\beta^2$ at $\beta=1/200$) and $E=2719/1000$, put
\begin{align*}
A&=24u^4+64u^6+\sum_{m=7}^{96}\binom{96}{m}u^m,\qquad B=1024u^6+\frac{1536}{7}\,E(20E)^6\frac{u^7}{1-20Eu},\\
Y&=B+B^2,\qquad R=16A+\frac{Y}{1-Y}.
\end{align*}
The comparison to be certified in exact rational arithmetic is $B\le1$, $Y<1$ and $R\le1741/10^6$. The program reports $R\approx0.0017400504$, $A\approx1.0362\cdot10^{-6}$ and
$B\approx0.0017176$; but the stored output does not contain the exact fraction (the field is empty), and these decimals are not the certificate. (ENC) and (RAT) are the two
computational inputs of Proposition~\ref{prop:instance}. The same implementation also reports that this upper bound exceeds $2^{-9}$ at $\beta=0.0051$, and reports $\Delta_\beta(2)\le2^{-12}$ for $\beta\le0.0038$ \st{computed}; these additional computations are not certified by the stated input (RAT).

\begin{remark}
The exact counts $c_m$ for $m\le6$ and the exact rational evaluation come from a single program (\path{reproduction/instance/instance_n2_M6.py}) and
have not been independently recounted. The bound is crude: the activity $\varepsilon\approx\beta$ is four times the true expansion parameter $t\approx\beta/4$, and no
cancellation between the two tori is used, so $\beta_*$ is about \BETARATIO{} times smaller than the leading-order value $0.190$. A
cluster expansion with the true activities and the cancellation of non-winding clusters should reach $\beta\approx0.1$--$0.2$ for $n=2$ \st{heuristic},
and $n=4$ up to $\beta\approx1$ would require a computer-assisted expansion with rigorous tails; we have not done either. In this regime
exponential decay and a lattice mass gap are already known from the cluster expansion of Osterwalder and Seiler \cite{OS78}, with better
bounds: Proposition~\ref{prop:instance} is a reported check conditional on (ENC) and (RAT), not a new result on the gap.
\end{remark}

%=====================================================================================================================
\section{Discussion}\label{sec:discussion}

\subsection{What this note does not prove}\label{sec:notprove}
\begin{itemize}
\item \textbf{No weak-coupling estimate.} The theorem requires a scalar bound at some fixed physical size (of order $\ell\sim2$\,fm
for the physical theory), uniformly along the chosen cutoff trajectory. No estimate here establishes this uniform scalar hypothesis along a weak-coupling continuum trajectory (the ``desert'' between strong and weak coupling); establishing it remains an open sufficient route, and equivalence with the Yang--Mills existence and mass-gap problem is not proved. The strong-coupling interval reported in Proposition~\ref{prop:instance}, conditional on (ENC) and (RAT), lies in a regime where the gap is already known.
\item \textbf{No existence and no non-triviality of the continuum limit.} (W4) is assumed; a trivial limit satisfies the conclusion
for every $q$.
\item \textbf{No reconstruction.} (W6) is assumed, as in ClayYM.
\item \textbf{The criterion contains additional finite-volume information.} It controls excited thermal traces, including flux
sectors, which are not determined by the local infinite-volume gap alone in the broader class illustrated by Example~\ref{ex:z2}. That example
has a trivial local limit: strictness within the full Clay class or within pure $SU(2)$ is not established.
\item \textbf{The transfer structure (W2)--(W3) is not formalised} (it would require Haar integration on $SU(2)^{\text{links}}$ and the
compact spectral theorem in Lean), and the comparison of tori and cylinders is on paper only.
\item \textbf{The relation with hypothesis (D) of \cite{BB26} is not established in either direction.} In particular, (D) does not
obviously give (G$_{\mathrm{fv}}$) on the flux sectors. Hypothesis (D) is that of version~4.45.0. For hypothesis (B$^*$) of version~5.0.0
\cite{BB26v5} the two criteria are equivalent up to the scale: (B$^*$) at $T_0$ implies $\Delta_\beta(2T_0)\le2^{-9}$ by \eqref{eq:bb26implies}, and $\Delta_\beta(n)\le2^{-9}$ implies (B$^*$) at $16n$ (Theorem~\ref{thm:R}). Neither implies that the hypothesis holds along a weak-coupling trajectory, and Corollary~\ref{cor:bridge} is conditional and unverified by us.
\item \textbf{Sections~\ref{sec:thooft} and~\ref{sec:fixed} prove nothing at weak coupling in favour of the criterion.} Theorem~\ref{thm:twistall} gives a length, not an area, and no area law in infinite volume; (V1) and (V2) are not proved at any weak coupling; Corollary~\ref{cor:L1} shows that the criterion is \emph{false} at every fixed size $n$ as $\beta\to\infty$, so the scale $n$ must grow with $\beta$ (Proposition~\ref{prop:scale} gives a conditional polynomial lower bound, far below the scale suggested by asymptotic freedom), and nothing here controls that growth. No statement of this note is presented as a proof of the mass gap.
\item The gap bound $(\log2)/\ell$ is a sufficient lower bound, not an optimal estimate; the comparison with the massive scalar (a factor
$\approx16$) uses the uncertified numerical crossing and is only illustrative.
\end{itemize}

\subsection{Comparison with related conditional statements}\label{sec:compare}
\begin{itemize}
\item \emph{OpenAI's two-dimensional criterion \cite{OAI26}} \st{read} (Lemma~2.1, Proposition~2.2 and the rectangle identity, in
\path{sections/finite-volume.tex} of the repository). Two-dimensional spin models with a positive, self-adjoint, trace-class transfer operator,
threshold $2^{-10}$, recursion $44\Delta^2\le64\Delta^2$, gap $(\log2)/n$ of every local limit of square tori, and of the lattice theory if the
periodic state is unique; the finite-volume section supplies the excited-trace estimate, the rectangle identity, the quadratic iteration and
the insertion comparison between square tori and cylinders. Theorem~\ref{thm:lattice} transposes this argument to $d=4$ and $d=3$, using the
weights $2^{d-1-j}$ of \eqref{eq:decomp} and the full gauge-invariant trace, flux sectors included. The explicit constants and the Lean packaging
are incremental additions; the continuum passage is a conditional closure of the $q$-form, and the converse is an elementary thermal-trace
estimate. We found no converse and no $d>2$ statement in the passages of \cite{OAI26} that were read (the passages of the finite-volume section and introduction recorded in the anteriority audit, lines 36--135 and 222--262 respectively, together with bounded keyword searches elsewhere in the source).
These additions do not constitute a substantially new mass-gap mechanism, and the priority assessment is limited to the sources and passages
actually read.
\item \emph{Bhattacharjee et al., Theorem~B \cite{BB26}} \st{read}. ``Let $a_k\to0$ and $\beta_k\ge0$ be couplings, and suppose that there
are $\varpi>0$, $\ell>0$ and integers $n_k\ge2$ with $n_ka_k\le\ell$ such that $\mathrm{BP}(n_k,\varpi)$ holds for the Wilson measure at
$\beta_k$ on all large tori. Then every continuum limit [\dots] has Hamiltonian $H$ with $\sigma(H)\subset\{0\}\cup[\Delta,\infty)$,
$\Delta=16\varpi/(625e^2\ell)$'' (v4.45.0, \texttt{p1\_intro.tex}). $\mathrm{BP}(n,\varpi)$ is a variance (Poincar\'e) inequality for the
dynamics that resamples cubes of side $n$, on all large tori. Same form as Theorem~\ref{thm:main}; the hypothesis is a spectral gap of a
dynamics, not an inequality between two partition functions, and the limit is a separate hypothesis (E). Neither hypothesis is known to
imply the other.
\item \emph{Bhattacharjee et al., version~5.0.0, Theorem~7.56 and Theorem~F \cite{BB26v5}} \st{read} (abstract, introduction, Sections~7.8
and~7.11 only). The hypothesis is no longer a spectral gap of a dynamics: it is a family of inequalities between partition functions of
finite tori, on the same doubling defect as ours. Their record is earlier than version~1.0 of this note; their hypothesis implies ours, and ours implies theirs at a scale $16$ times larger (Theorem~\ref{thm:R}); their conclusion extends beyond the dyadic-cubes conclusion of Theorem~\ref{thm:lattice}, and Theorem~\ref{thm:R} supplies our extension to arbitrary integer cubes and bounded-aspect-ratio tori (\S\ref{sec:bb26v5}).
\item \emph{Rajabu, Theorem~5.3 \cite{Raj26}}, as summarised in the project's earlier reading note (the primary text was not re-read in the anteriority audit for this version). A gauge-covariant block hierarchy (H1)--(H3), weighted Lipschitz covariance and
transfer bounds at every fixed physical scale (C1), a landing bound (C2), time continuity (C3) and a non-zero excited sector (C4) give a
gap $\ge w$ for the transfer operators of limits of block observables. Four multi-scale hypotheses against one scalar inequality; none of
them is discharged here, and (C3)--(C4) have no counterpart (non-triviality is not treated here).
\item \emph{Tomboulis--Yaffe, Kov\'acs--Tomboulis, Kanazawa} \cite{TY85,KT02,Kan} \st{read} (\cite{TY85} pp.~317 and 332, \cite{KT02} eqs.~(8) and (11), \cite{Kan} Theorems~1, 2 and~5, read in the texts; page numbers are those of the printed paper for \cite{TY85} and of the arXiv versions for the others; \cite{TY85} was read on these pages only). Same family of argument in \cite{KT02,Kan}: reflection positivity, Schwarz inequality, repeated doublings of a loop or of a correlation until periodicity removes it, with divisibility conditions on the sides ($L_\mu=2^pl_\mu$ in the proof of \cite[Theorem~2, p.~9]{Kan}, whose footnote~10 says that the same condition was present in \cite{TY85}; reflection positivity is applied to a two-point function in \cite[p.~332]{TY85}), and an inequality between partition functions of finite periodic tori. Different hypothesis: insensitivity of the torus to a centre twist on a stack of plaquettes, $\langle\mathcal O[\mathcal V]\rangle=Z^{(-)}/Z\to1$ (equivalently $e^{-F_{\mathrm{el}}}=\frac12(1-Z^{(-)}/Z)$ small), not the doubling defect $16\log Z(n^4)-\log Z((2n)^4)$ of this note. Different conclusion: a bound on Wilson loops of a periodic box, $\langle W[C]\rangle\le2(e^{-F_{\mathrm{el}}})^{|A_C|/|A_{12}|}$ (and, in \cite{Kan}, on correlations of spins), and not a bound on the excited eigenvalues of the transfer operator on every spatial torus. We do not claim that either hypothesis implies the other. These papers also contain the precedent recalled in \S\ref{subsec:notnew}, and are not used in our proofs.
\item \emph{Knabe-type criteria} \cite{Kna88,GM16,Lem19}. Finite-size gap criteria for frustration-free quantum Hamiltonians: a local
gap above a threshold on a finite subsystem implies a global gap. Same logical form, different objects: local Hamiltonian gaps, no
partition functions, no reflection positivity, no continuum limit, one direction only.
\item \emph{Borgs--Koteck\'y} \cite{BK90}. A finite-size scaling theory of periodic partition functions at first-order transitions. With
such an expansion $\Delta_d\to(2^d-1)\log(\#\text{phases})$ (Example~\ref{ex:ising}): the test counts phases. This is our reading, not
a statement of \cite{BK90}, whose text we have not re-read.
\end{itemize}

\subsection{Reach}
The note supplies a higher-dimensional version of an existing finite-volume sufficient criterion and packages its conditional continuum
implication. It does not establish the scalar bound along a weak-coupling trajectory, construct the continuum limit, or identify an existing
construction with all the objects required by the criterion (Section~\ref{sec:openai-construction}). For square tori, the comparison must hold along
the actual convergent family. Under exact blocking, $Z_{\mathrm{block}}(n)=Z(bn)$ gives $\Delta_{\mathrm{block}}(n)=\Delta(bn)$: this covariance is
an identity, not a quantitative estimate that makes the missing smallness hypothesis easier.

%=====================================================================================================================
\section{Relation to the OpenAI construction}\label{sec:openai-construction}
\begingroup\small
OpenAI have released a Lean-formalised construction for four-dimensional $SU(2)$ lattice Yang--Mills theory. This comparison concerns the local
snapshot of that repository held by the author (folder dated 28 October 2026), inspected on 11 November 2026; file references are relative to it.
\textbf{The criterion does not apply to that construction as it stands}, for three reasons, and a fourth point limits what can be said.
\begin{enumerate}
\item \emph{The bare coupling may depend on the volume.} \path{YM/Source/Construction.lean:38} defines \verb|abbrev beta (N p) := C.path.beta N p N|.
The scalar test must then be evaluated at the common coupling $\beta_{N,p}$ of the target volume, $\Delta_{\beta_{N,p}}(n_N)$ with $n_N=\ell/a_N$;
a difference of two separately tuned couplings is not the defect controlled by the doubling argument. A volume-dependent coupling is not a
logical obstruction, but the required same-coupling estimate is an additional hypothesis.
\item \emph{The dyadic volumes are not exported.} \path{PublicWilson.lean:38} exports only the field \texttt{periodStrict} (a \texttt{StrictMono} statement on \texttt{(period p).val}).
The internal recipe contains \verb|4*(M+1)| and \verb|period i := page (M+i)| (\path{YM/Price/Waits.lean:87,102}), but the exported refinement
(\path{ChallengeSpec.lean:188}, \verb|T.Properties|) does not supply the hypothesis $L_{p_k}=2^{j_k}\ell$, $j_k\to\infty$ that our torus comparison uses.
\item \emph{The order of limits differs.} The fixed-volume limit uses \verb|C.smear N p v t| (\path{PublicWilson.lean:89}); the subsequent volume limit
uses \verb|row (volumeOrder k)| (line 90): first the mesh limit at fixed physical volume, then the volume limit. Our torus--cylinder estimate
(a diagonal choice of volumes as a function of the cutoff, Corollary~\ref{cor:tori}) changes the family and does not identify its limit with this
iterated limit.
\item \emph{Partition-function control exists but does not bound the defect.} \path{YM/FullVolume/Scalar.lean} contains the one-volume bound
\verb|partition G n b <= q^(3*Fintype.card (Site n)-3)| (line 46), \verb|raw_lost| (line 68) and a positive lower bound for \verb|stripped A N p m|
(line 98). These estimates do not give the same-coupling two-volume bound $\Delta\le\eta$ that the criterion requires; we assert no absence of
control of $\log Z$.
\end{enumerate}
An initial internal repair used an additional energy-bound hypothesis (EB). Later internal work reports a conditional four-copy transport estimate avoiding (EB), using moment bounds already available in the construction. These repairs are outside this note and do not establish the same-coupling scalar estimate or a complete identification with the objects used here; the present criterion therefore does not plug into that construction as it stands. Our searches in that repository for a bridge to our Lean objects (\texttt{ClayYM}, \texttt{rpForm}, \texttt{OSReconstruction}, \texttt{QGap}),
for a finite-volume spectral decomposition, and for a bound on $\Delta$ were bounded textual searches and found none; they are not an exhaustive audit.
\endgroup

%=====================================================================================================================
\section{Reproducibility}\label{sec:repro}
The package accompanying this note contains:
\begin{itemize}
\item \path{lean/CritereScalaire.lean} (this work) and \path{ClayYM_v1_1_extrait/ClayYM/} (six files of ClayYM v1.1, verbatim);
\path{lean/assembler_paquet.sh} builds the single file \path{main.lean} = \texttt{import Mathlib} + the six ClayYM files (imports
removed, each wrapped in a \texttt{section}) + \path{CritereScalaire.lean}; the SHA-256 of \path{lean/main.lean} is that of the original build
(\path{lean/MANIFEST.txt});
\item \path{lean/build_rejeu_2026-11-11.json}: the output of the recompilation from a clean copy by the author's side (\S\ref{app:lean}), 
\path{lean/reconstruction_non_auteur_2026-11-11/}: the log and project files of reconstruction (b), and \path{rejeu_depuis_arbre_git_2026-11-12/}: the script, scaffolding and log of reconstruction (c);
\item \path{reproduction/scripts/constantes_poly.py} (constants of Table~\ref{tab:const}), \path{reproduction/scripts/temoins.py} (witnesses), with their
outputs in \path{reproduction/resultats/};
\item \path{reproduction/instance/instance_n2.py} (Proposition~\ref{prop:instance}); the file actually run, \path{instance_n2_M6.py}, differs
only by fixing the maximal exact size to $6$; its output is \path{sortie_bac_M6_job_20261111-011302-d9004d.json} ($0.4$\,s);
\item \path{lean/Reciproque.lean} (Theorem~\ref{thm:R}, abstract form), \path{lean/main_v1_2.lean} (the published \path{main.lean} followed by the new file, byte for byte), \path{lean/reciproque/} (frozen statements, comparison script, mutants, certificate), and \path{reproduction/reciproque/} (the numerical checks of Remark~\ref{rem:Rnum}, with outputs);
\item version 1.3: \path{lean/thooft_critere/} and \path{lean/verrou_L/} (frozen statements, proofs, certificates and replay report of the two abstract files of Appendix~\ref{app:lean13}), \path{lean/sorties_rejeu_v1_3/} (sandbox outputs of 12 November 2026), \path{calculs/} (the constants quoted in Sections~\ref{sec:thooft} and~\ref{sec:fixed}, with outputs from the two sandboxes, and the comparison of frozen and proved statements);
\item \path{SHA256SUMS}.
\end{itemize}
All programs ran in a sandbox without network and without privileges (``qlab'', Python~3, Lean~4 v4.34.0 with Mathlib commit
\texttt{5ed2965}, tag v4.34.0). The command used for Lean is
{\footnotesize
\begin{verbatim}
bash lean/assembler_paquet.sh /path/to/main.lean
<sandbox runner> --kind lean --timeout 1200 /path/to/main.lean
\end{verbatim}}
\noindent and, for Python, the same runner with \texttt{--kind python --timeout S FILE}. Outside this sandbox, any Lean~4 v4.34.0 project with
Mathlib at that commit compiles \path{main.lean} with \texttt{lake env lean main.lean}.

%=====================================================================================================================
\appendix
\section{The Lean formalisation}\label{app:lean}
\paragraph{Files and build.} \path{main.lean} (2175 lines), SHA-256\\
{\footnotesize\texttt{944a99f6fc10f1696e6dfa5bc231b73dd9c57a101b48518c52e3e4ca944ef2c6}},\\
is assembled from \path{CritereScalaire.lean}, SHA-256\\
{\footnotesize\texttt{0c4264a7c3f665346d1e479f6ed99e8b531b3b23960efcc718095d4ed4fe8496}},\\
and the ClayYM v1.1 files \texttt{Basic}, \texttt{OS}, \texttt{Gap}, \texttt{Fake/Trivial}, \texttt{Transfer}, \texttt{GapQ} (commit
\texttt{93ac204}). Toolchain: Lean~4 v4.34.0, Mathlib \texttt{5ed2965}. No \texttt{axiom} is declared; the hypotheses are explicit arguments.
The following build checks, all on 11 November 2026. (a) On the author's side, in the sandbox (final build \texttt{build4}, job \texttt{20261111-002010-32231f}) and again
from a clean copy of the package (job \texttt{\REBUILDJOB}): return code $0$, $0$ errors, no \texttt{sorry}; the $35$ declarations of the namespace
\texttt{CritereScalaire} depend only on \texttt{propext}, \texttt{Classical.choice}, \texttt{Quot.sound}. (b) By a non-author (an AI agent, Claude
Sonnet~5.5, not the author of the Lean files), about 15:20 UTC, on a clean rented machine with the official Lean~4.34.0 toolchain and the
official Mathlib cache: \texttt{lake build} return code $0$ (8925 jobs; $187$\,s for \texttt{lake build}, following a $25$\,s cache check) on the file \path{main.lean} of the hash above, and the same $35$
declarations with standard axioms; this one started from hashed copies of the sources, not from a git archive.
(c) By a non-author (Claude Sonnet~5.5, not the author of the Lean files), from the \emph{git archive of the package}: the package directory
\path{papers/Paper_ScalarCriterion_SU2} was added to a scratch git repository (no commit in the public repository), its tree computed, and
\texttt{git archive} of that tree copied to the rented machine (archive sha256 \texttt{97cd82b6\dots01573c}; git tree \texttt{16d2a27b\dots} of the package at
that moment). The subtrees used, \path{lean/} (git id \texttt{0bc5d1f1\dots}) and \path{ClayYM_v1_1_extrait/} (\texttt{f6c358cb\dots}), are those of the final package; the
note, the metadata and the logs, which differ, are not used by the build. On the machine, in a new directory, under the unprivileged user, at
$20{:}03$--$20{:}05$ UTC, with Lean~4.34.0 and Mathlib \texttt{5ed2965}: the ten hashes of \path{lean/MANIFEST.txt} were checked on the extracted files (all OK);
the lake scaffolding (\texttt{lakefile.toml}, \texttt{lake-manifest.json}, \texttt{lean-toolchain}, a file of \texttt{\#print axioms}) is not in the archive and was taken from
reconstruction (b); the Mathlib package directory was copied from reconstruction (b) (official cache: ``no files to download''), but no build artefact of the note was
copied, and \path{main.lean} was taken from the archive (sha256 \texttt{944a99f6\dots}); \texttt{lake build} return code $0$ (8925 jobs, \texttt{main} built in $28$\,s; $33$\,s for \texttt{lake build}, following a $13$\,s cache check), the
same two linter warnings as before, and the same $35$ declarations with axioms \texttt{propext}, \texttt{Classical.choice}, \texttt{Quot.sound}, the list being identical to that of \texttt{build4}. The logs are
in the package (\path{rejeu_depuis_arbre_git_2026-11-12/}). These descriptions of rebuilds (a)--(c) concern version~1.0; version~1.2 adds separate Lean files, whose rebuilds are described below. Rebuilds (b) and (c) compiled the already assembled \path{main.lean}; they did not rerun the assembler from the component sources. They used the standard Lean kernel, not an independent kernel. The complete tree of the final package, which differs from the archived one only outside the two subtrees above, was not itself rebuilt. Version~1.1 (12 November 2026) changed the note, the README and the metadata files only. Version~1.2 adds files (below) and changes the note; \path{lean/CritereScalaire.lean}, \path{lean/main.lean}, \path{ClayYM_v1_1_extrait/} and \path{rejeu_depuis_arbre_git_2026-11-12/} are byte for byte those of version~1.0. An independent kernel check (nanoda) has \emph{not} been run for \path{main.lean}; it was run for the new file, as described next.

\paragraph{Version 1.2: the reverse implication in Lean.}
\path{lean/Reciproque.lean} (474 lines, SHA-256 \texttt{efa8870dc236ef5d5a35d92e45e7cd7408573b2229e18d633404c017ab9d4dbb}) contains four principal theorem statements of the namespace \texttt{Reciproque} together with auxiliary lemmas and proof theorems (seven declarations introduced by \texttt{theorem}); the four principal statements were frozen on 12 November 2026 before their Lean proofs were written (file \path{lean/reciproque/ENONCES.lean}, SHA-256 \texttt{6a2c0e9f\dots ec120af8}); a script of the package (\path{lean/reciproque/verif_statique.py}) checks that the four statements and the definitions \texttt{sT}, \texttt{Delta}, \texttt{Rep}, \texttt{HypSpec}, \texttt{HypRec} are identical, character for character after whitespace normalisation, to the frozen ones. The abstract setting: a function $F$ on four periods, symmetric (\texttt{hs12}, \texttt{hs23}, \texttt{hs34}); \texttt{Rep}: for each section $(a,b,c)$ with positive sides, ratios $r_i\in[0,1]$ with $\sum_ir_i^d<\infty$ and $F(a,b,c,d)=dL+\log(1+\sum_ir_i^d)$ for every $d\ge1$ (this is the already published \texttt{HypSpecTC}, with the positivity $r_i\ge0$ in it); \texttt{HypRec}: the published recursion, $\Delta(m)\le2^{-9}\Rightarrow\Delta(2m)\le256\Delta(m)^2$, as a named hypothesis. The file \path{lean/main_v1_2.lean} is the published \path{main.lean} followed by this file inside a \texttt{section}, plus the theorem \texttt{Reciproque.reciproque\_publie}, which \emph{derives} \texttt{HypRec} from the published theorem \texttt{CritereScalaire.doublement\_fin\_4d} for any $F$ with \texttt{CritereScalaire.HypSpecTC}; its first 114\,899 bytes are \path{main.lean} (checked by comparison). \path{CritereScalaire.lean} and \path{main.lean} are thus unchanged, and the new content is separate and checkable.
\begin{center}\footnotesize
\begin{tabular}{@{}>{\raggedright\arraybackslash}p{0.30\textwidth}>{\raggedright\arraybackslash}p{0.63\textwidth}@{}}
\toprule
declaration & what it states\\
\midrule
\texttt{Reciproque.facile} & the easy direction of \S\ref{sec:bb26v5} in abstract form: from $a_j\le\gamma2^jn^3$ ($n=2T_0$) and $\Delta\le16a_0+8a_1+4a_2+2a_3$, $2\gamma(96T_0)^3\le\frac1{20}$: $\Delta\le512\gamma T_0^3\le1/69120<2^{-9}$\\
\texttt{Reciproque.interpolation} & Theorem~\ref{thm:R}(a), first inequality: $\log(1+s_t(m))\le2\Delta(m)V(t)/(31m^3)$ for $t$ with sides $\ge m$, given $\Delta(m)\le2^{-9}$\\
\texttt{Reciproque.reciproque\_Bstar} & Theorem~\ref{thm:R}(b): $2\gamma(96T_0)^3\le\frac1{20}$ and $\log(1+s_t(16n))\le\gamma V(t)$ for all sides $\ge16n$, $\gamma=1/(31\cdot2^{23}T_0^3)$\\
\texttt{Reciproque.trou\_cubes} & Theorem~\ref{thm:R}(c), first half: $r_i^n\le256\Delta(n)$ for the cube $(N,N,N)$, every $N\ge n$\\
\texttt{Reciproque.reciproque\_publie} (in \path{main\_v1_2.lean} only) & \texttt{reciproque\_Bstar} for $F$ with \texttt{CritereScalaire.HypSpecTC}, symmetric, with $\texttt{CritereScalaire.Delta}\le\frac1{512}$ at $n$\\
\bottomrule
\end{tabular}
\end{center}
\emph{Not formalised}: (d) of Theorem~\ref{thm:R} (aspect ratio $\le24$, the constant $e^{18/31}-1$); the logarithmic form of (c); the Casimir-type bounds; the passage from sides $\ge3$ (paper) to sides $\ge1$ (Lean hypothesis, stronger); Corollary~\ref{cor:bridge}; the verification of (W2) for Wilson's theory.
\paragraph{Rebuilds of the new files (12 November 2026, by a non-author).} (1) Sandbox, Lean 4 v4.34.0, Mathlib \texttt{5ed2965}: \path{Reciproque.lean} (job \texttt{20261112-070619-0a612c}) and \path{main\_v1\_2.lean} (job \texttt{20261112-070619-f382b1}, with one extra \texttt{\#print axioms} line for \texttt{interpolation}): return code $0$, $0$ errors, no \texttt{sorry}; the four theorems, their auxiliary lemmas and \texttt{reciproque\_publie} depend only on \texttt{propext}, \texttt{Classical.choice}, \texttt{Quot.sound}; the sandbox reports no axiom outside these. The same two files were compiled a second time in the second sandbox (same toolchain, other machine; jobs \texttt{20261112-072756-20ad77} and \texttt{20261112-072756-55f333}): return code $0$, $0$ errors, standard axioms only. (2) Rented machine, \emph{Lean 4.32.0, Mathlib v4.32.0} (not the version of the sandbox), unprivileged user: Lean 4.32.0 (commit \texttt{8c9756b2}) with the Mathlib v4.32.0 of a project present on the machine, \texttt{lake env lean}, \texttt{nice 10}: both files, unmodified (same SHA-256), compile with return code $0$ and no error (linter warnings on unused variables only), and give the same standard-axiom lines. The statements were not adapted for this version. (3) Independent kernel \emph{nanoda}: for each file, the closure of the declarations (the four theorems; for \path{main\_v1\_2.lean} also \texttt{reciproque\_publie} and \texttt{CritereScalaire.doublement\_fin\_4d}) was exported by \texttt{lean4export} (binary SHA-256 \texttt{e5736998\dots}) and replayed by the independent kernel nanoda (binary \texttt{ce46eaff\dots}, unmodified, only the three standard axioms permitted, any other an error): \path{Reciproque.lean}, 1\,377\,205 export lines, ``Checked 16949 declarations with no errors''; \path{main\_v1\_2.lean}, 1\,433\,069 lines, ``Checked 17156 declarations with no errors''. This is the first nanoda check of any Lean file of this note; \path{main.lean} alone was not replayed, but \texttt{doublement\_fin\_4d} and its dependencies are in the second export. (4) Mutants of \path{Reciproque.lean}: constant $31\to30$ in the statement of \texttt{reciproque\_Bstar}; factor $16\to8$ in its window; the positivity $r_i\ge0$ removed from \texttt{Rep}; the threshold $\frac1{512}\to\frac1{256}$ in the hypothesis of \texttt{reciproque\_Bstar}: all four make the build fail (errors located inside the targeted declarations or their users; \path{lean/reciproque/mutants/RESULTATS.txt}). The rebuilder is an instance of Claude Sonnet~5.5 that did not write the proofs; it also wrote the comparison script.


\begin{center}\footnotesize
\begin{tabular}{@{}>{\raggedright\arraybackslash}p{0.30\textwidth}>{\raggedright\arraybackslash}p{0.63\textwidth}@{}}
\toprule
declaration & what it states, relative to the text\\
\midrule
\texttt{lemme\_fin\_tc} (l.\,42) & Lemma~\ref{lem:squaring}: $p\ge0$, $s\le e^p-1$, $r_i^n\le s$, $p(2n)\le2(e^p-1)^2$ (countable spectrum, $r_i\in[0,1]$, $s(n)<\infty$)\\
\texttt{decomp\_4d} (l.\,156) & $\Delta\ge0$, $a_0\le\Delta/8$, $a_3\le\Delta$, from symmetry and positivity of $p$; not the full identity \eqref{eq:decomp} nor $a_1,a_2$\\
\texttt{doublement\_fin\_4d} (l.\,186), \texttt{chaine\_num\_fin4} (l.\,125) & Proposition~\ref{prop:doubling}, with positivity and squaring bound of $p$ as hypotheses (supplied by \texttt{hyp\_of\_specTC} from \texttt{HypSpecTC})\\
\texttt{iteration256} (l.\,248), \texttt{gap\_4d\_fin} (l.\,299) & Theorem~\ref{thm:lattice}: $r_i^n\le256\,\Delta(n)$ on the section $(2^kn)^3$ given its spectral data; no simplicity, no logarithmic form\\
\texttt{doublement\_fin\_3d} (l.\,418), \texttt{gap\_3d\_fin} (l.\,513) & Theorem~\ref{thm:3d}\\
\texttt{SpecLink} (l.\,570), \texttt{qGapOn\_of\_specLink} (l.\,581) & (W3) in coordinates (linear square-summable $c$, $c(\Omega)=\delta_0$, $B(\Omega,F,0)=c_0(F)$, diagonal identities at times $0$ and $\ell$) and Lemma~\ref{lem:qform}\\
\texttt{LienTransfert4} (l.\,611), \texttt{qGapOn\_limite\_4d} (l.\,628) & abstract form of Theorem~\ref{thm:main}, any filter; no $a_k$ or equality of spacings appears in Lean: \texttt{SpecLink} relates the integer exponent $n_k$ to the real time $\ell$, and $\ell=n_ka_k$ is the assumed physical identification\\
\texttt{convergesOn\_of\_proche} (l.\,672) & transfer of convergence \emph{assuming} that the torus--cylinder difference tends to zero; the estimate and the diagonal choice are not formalised\\
\texttt{clay\_trou\_de\_critere\_4d} (l.\,700), \texttt{\_3d} (l.\,722) & Theorem~\ref{thm:main} with \texttt{QGap}, \texttt{HasMassGap} of ClayYM and, in $d=4$, $(\log2)/\ell\le-\log(256\eta)/\ell$; vacuum uniqueness is \texttt{QGap.vacuum\_unique} of ClayYM, the case $\eta\le2^{-12}$ is not restated\\
\texttt{non\_vacuite\_4d} (l.\,802) & the hypotheses of \texttt{qGapOn\_limite\_4d} hold simultaneously for the trivial theory\\
\bottomrule
\end{tabular}
\end{center}

\paragraph{Named hypotheses in Lean.} Symmetry \texttt{hs12}, \texttt{hs23}, \texttt{hs34} (W1); \texttt{HypSpecTC} (W2), for all
positive sizes; bilinearity \texttt{IsBilin} of the lattice forms; \texttt{LienTransfert4}/\texttt{SpecLink} (W3); \texttt{ConvergesOn} (W4); \texttt{OSReconstruction R S} and
\texttt{S [] = 1} (W6). The guard \texttt{non\_vacuite\_4d} instantiates the hypotheses of the abstract theorem \texttt{qGapOn\_limite\_4d} with $V=\R$ and $B(F,G,t)=FG$. It does not instantiate the reconstruction hypothesis, and it does not show that the hypotheses hold for Wilson's $SU(2)$ theory.

\paragraph{Mutation tests.} Eight mutants of \path{CritereScalaire.lean} must each make the build fail: constant $256\to200$; threshold
$2^{-9}\to2^{-7}$; a wrong symmetry in \texttt{decomp\_4d}; vacuum coordinate $c(\Omega)\equiv1$; a poisoned theorem $0=1$; constant $64\to56$ in
$d=3$; conclusion $q=128\eta$; spectral link without the weights $r_i^n$. \MUTANTS


\section{Lean files of version 1.3}\label{app:lean13}
Version 1.3 adds two Lean files, both in abstract finite form. \emph{They check small pieces of Sections~\ref{sec:thooft} and~\ref{sec:fixed} and nothing else; the analytic parts of Theorem~\ref{thm:L}, the identification of $Z_k$ with $\Tr(C^kK^M)$ for Wilson's theory, Theorem~\ref{thm:twistall}(b),(c), Proposition~\ref{prop:scale} and the passage from operators to finite spectra are not formalised.} Neither file is part of \path{main.lean}, and none of the files of versions 1.0 to 1.2 is changed.

\paragraph{Files, freezing and builds.} \path{lean/thooft_critere/ENONCES.lean} (7 statements with \texttt{sorry}, SHA-256 \texttt{7dd1a95c\dots}) and \path{lean/thooft_critere/PREUVES.lean} (SHA-256 \texttt{58a549d7\dots}); \path{lean/verrou_L/Enonces.lean} (15 statements with \texttt{sorry}, SHA-256 \texttt{06c73577\dots}) and \path{lean/verrou_L/Preuves.lean} (SHA-256 \texttt{07ee1c56\dots}). In each case the statements were frozen before any proof was written, the proof files contain the same statements and definitions (text comparison of the statements, \path{calculs/comparer\_enonces.py}: $7=7$ and $15=15$ identical; the only added declarations are lemmas named \texttt{aux\_*}). Lean~4 v4.34.0 with the Mathlib of the sandboxes; no \texttt{axiom} is declared; the proof files contain no \texttt{sorry}; \texttt{\#print axioms} gives only \texttt{propext}, \texttt{Classical.choice}, \texttt{Quot.sound} for the proofs of the first file and, for the second file, a subset of these (\texttt{c1\_int}: \texttt{propext}; \texttt{c2\_nat}, \texttt{c2\_int}: \texttt{propext}, \texttt{Quot.sound}; \texttt{a4\_unites\_commutatif}: \texttt{propext}). On 12 November 2026 (11h17--11h18 CET) the four files were compiled again by the preparer of this version (an instance of Claude Sonnet~5.5) in the two sandboxes (the one of the first machine, and the 8-core sandbox of the second machine): return code $0$, $0$ errors, no \texttt{sorry}, no axiom outside the three standard ones, in the four runs (jobs \texttt{20261112-111703-9c632b}, \texttt{-111703-ecda1d}, \texttt{-111708-74f1b3}, \texttt{-111717-883e64}; outputs in \path{lean/sorties\_rejeu\_v1\_3/}). Four mutants of the first file (a hypothesis removed, the constant $30$ replaced by $29$, the weight $8$ by $7$, the positivity of the $G_k$ removed) fail to compile; a replay by a non-author proved the negation of twelve mutated statements of the first file by explicit small counter-examples and checked that the unmutated hypotheses of each are satisfiable (non-vacuity of three statements). \emph{Not done:} an independent second kernel (nanoda), a build on Lean~4.32.x, a mutation campaign for the second file beyond its own abelian cases, and a rebuild from a git archive.

\paragraph{Who.} The first file was written by the agent (Claude Opus~5.5) that wrote the paper proofs of Section~\ref{sec:thooft}; the replay and the comparison of each statement with the prose were done by an instance of Claude Sonnet~5.5 that did not write it. The second file was formalised by an instance of Claude Sonnet~5.5 that did not write the paper proofs of Section~\ref{sec:fixed}.

\paragraph{First file: abstract finite form of Proposition~\ref{prop:sectors} and of the last links of Theorems~\ref{thm:twistall} and~\ref{thm:twosub}.} The data are a \emph{finite} family $r_i\ge0$ with a distinguished index $i_0$, $r_{i_0}=1$ (the vacuum), a sector label $e\in\Z_2^3$ for each $i$, and $Z_k(M)=\sum_i(-1)^{k\cdot e_i}r_i^M$.
\begin{center}\footnotesize
\begin{tabular}{@{}>{\raggedright\arraybackslash}p{0.2\textwidth}>{\raggedright\arraybackslash}p{0.43\textwidth}>{\raggedright\arraybackslash}p{0.30\textwidth}@{}}
\toprule
declaration & what it states & difference from the text\\
\midrule
\texttt{chi\_somme} & $\sum_k\chi_k(e)=8$ if $e=0$, $0$ otherwise (characters of $\Z_2^3$) & none\\
\texttt{trace\_neutre} & $\sum_kZ_k=8\Tr_0$ & finite spectrum\\
\texttt{defaut\_torsion} & if the vacuum lies in the sector $e=0$ and $r_i\ge0$: $0\le1-Z_k/Z_0\le2s(M)$ & \emph{coarse form only}: the refinements $\le\min(2D,1)$ and $2s^{\mathrm{fl}}/(1+s)$ of Proposition~\ref{prop:sectors}(b),(e) are not stated, and the positivity $Z_k\ge0$ is not used\\
\texttt{decomposition} & $\log Z_0=\log(1+s^0)+\log(8Z_0/\sum_kZ_k)$ with $s^0=\Tr_0-1$ & finite spectrum\\
\texttt{vide\_neutre} & if all $G_k=\sum_e\chi_k(e)g_e\ge0$ and $\sum_eg_e\ge1$ then $g_0\ge1$ & \emph{combinatorial core only}: the $G_k$ are taken as hypotheses; the limit $M\to\infty$ for a trace-class operator is not formalised\\
\texttt{propagation} & if $r_i^n\le q$ for $i\ne i_0$ and $s(m)\le A$ then $1-Z_k/Z_0\le2Aq^l$ at time $m+nl$ & last step of Theorem~\ref{thm:twistall}(a) only; $q=256\Delta$ and $A=q^p/248$ are not in the statement\\
\texttt{deux\_sous\_verrous} & $D_j\le7\varepsilon_1/8$, $s^0_j\le\varepsilon_2$, $\varepsilon_1\le\frac18$, and $\Delta\le2\sum_j2^{3-j}(\log(1+s^0_j)-\log(1-D_j))$ imply $\Delta\le30(\varepsilon_1+\varepsilon_2)$ & \emph{last link only}: the bracket of Proposition~\ref{prop:bracket} is a hypothesis, $D_j\le7\varepsilon_1/8$ is a hypothesis, and the conclusion is the coarse $30(\varepsilon_1+\varepsilon_2)$, weaker than the fine form of Theorem~\ref{thm:twosub}(a)\\
\bottomrule
\end{tabular}
\end{center}
Not formalised: the passage from trace-class operators to finite spectra, the limit of \texttt{vide\_neutre}, the identification of the twisted partition functions of Wilson's theory with $\Tr(C^kK^M)$, Theorem~\ref{thm:twistall}(b),(c), Theorem~\ref{thm:twosub}(b)--(d), Proposition~\ref{prop:scale}.

\paragraph{Second file: the algebraic fact and the real algebra of Corollary~\ref{cor:L1}.}
\begin{center}\footnotesize
\begin{tabular}{@{}>{\raggedright\arraybackslash}p{0.2\textwidth}>{\raggedright\arraybackslash}p{0.74\textwidth}@{}}
\toprule
declaration & what it states\\
\midrule
\texttt{a1\_anticommute} & for $U,V$ in \texttt{Matrix.specialUnitaryGroup (Fin 2) $\mathbb C$} with $UV=-VU$: $\Tr U=\Tr V=0$, $U^2=V^2=-1$ (the proof uses only $\det=1$)\\
\texttt{a2\_commutant} & if moreover $W\in SU(2)$ commutes with $U$ and with $V$, then $W=\pm1$\\
\texttt{a3\_existence} & $i\sigma_1$ and $i\sigma_2$ (explicit matrices) are in $SU(2)$ and anticommute\\
\texttt{a4\_unites\_\allowbreak commutatif}, \texttt{a4\_integre}, \texttt{a4\_cercle\_\allowbreak et\_Z2} & anticommutation is impossible for units of a non-trivial commutative ring in which $2$ is invertible, for non-zero elements of an integral domain with $2\ne0$, and for the elements of \texttt{Circle} and of $\mathbb Z^\times$\\
\texttt{a5\_mutant\_\allowbreak commuter\_\allowbreak avec\_U\_seul} & commuting with $U$ alone does not force $W=\pm1$ ($W=U$)\\
\texttt{b1\_cube} & $0\le R\le1$, $D\ge\frac38(1-R)$, $D<1$, $8(-\log(1-D))\le\Delta$ imply $\Delta\ge8\log(8/(5+3R))$ (the hypotheses on $R$ are declared and not all used)\\
\texttt{b2\_sept\_rapports}, \texttt{b2\_log4} & seven ratios $0\le R_k\le r$ with $1-D=(1+\sum_kR_k)/8$ give $-\log(1-D)\ge\log(8/(1+7r))$, and $\ge\log4$ if $r\le1/7$\\
\texttt{b3\_poids} & $8+4+2+1=15=\sum_{j<4}2^{3-j}$ and $15\log8=45\log2$\\
\texttt{c1\_nat}, \texttt{c1\_int}, \texttt{c2\_nat}, \texttt{c2\_int} & the dimension counts $(3dV-3V)-(3dV-3V+1-d)=d-1$ and $6+9(V-1)=9V-3$, in $\mathbb N$ (with hypotheses) and in $\mathbb Z$\\
\bottomrule
\end{tabular}
\end{center}
\emph{Not formalised:} all of the analytic part of Theorem~\ref{thm:L} (Lemma~\ref{lem:A}: gauge orbits and dimensions of the zero sets; Lemma~\ref{lem:B}: Morse--Bott; Lemma~\ref{lem:C}: Laplace and the exponents; Lemma~\ref{lem:D}: the $U(1)$ and $\Z_2$ bounds); the link between the algebraic fact and the lattice configurations (the cut gauge, the transitivity of $SO(3)$ on orthonormal pairs); the identification $D_j\to7/8$ on the tori $\sigma_j\times n$; (W1), (W2), (C); the left half of the bracket used in Corollary~\ref{cor:L1} (the second file formalises only the real algebra from the inequalities as hypotheses). The text and the Lean declarations were compared by the replayer; there is no second kernel.

%=====================================================================================================================
\section*{Use of AI}\label{sec:ai}
\addcontentsline{toc}{section}{Use of AI}
The work reported here was carried out with AI models, under the author's direction. The roles below are those recorded in the working
files of the project.
\begin{itemize}
\item \emph{Claude Opus~5.5 (Anthropic)} wrote, in several sessions on 10 and 11 November 2026, the transposition of the doubling mechanism to
$d=4$ and $d=3$, the chain of constants, the assembly with ClayYM, the converse and its examples, the witness programs, the Lean files, the
instance of Proposition~\ref{prop:instance} and the first draft of this note.
\item \emph{Qwen3.8 (Qwen team), run locally}, carried out the anteriority search of \S\ref{sec:compare} (sources read and not read are
recorded in the project files), derived independently the converse bound $\Delta_4(n)\le30A_0e^{m_0r_0}e^{-m_0n}$ and the Ising
counterexample, and \REVIEWSTATUS
\item \emph{GPT-6.1 Sol (OpenAI)} reviewed the draft of 11 November 2026 by reading, without running programs, and proposed scope and status
corrections, adapted and incorporated in this version (the hypotheses (W2$^+$), (W5), the qualification of the torus and filter statements, the
conditional status of the Maxwell witness, the formulation of the converse, the novelty and strictness claims, and the comparison with the OpenAI
construction); Sol also recommended that this material be folded into a broader note rather than stand alone; publishing it as a stand-alone note is the author's decision. On 12 November 2026 Sol reviewed, again by reading and without running programs, the first version of this text and required sixteen corrections (three blocking: the domain of the Lean specialisation, the status of Proposition~\ref{prop:instance}, and the order-eight coefficients), which Claude Sonnet adapted; a second read-only check of 12 November 2026 (\path{RELECTURE_SOL_TEXTE_CORRIGE.md}) found partial applications and asked for one blocking and four non-blocking corrections (the conditional status of Proposition~\ref{prop:instance} in the abstract and the list of results, the description of \texttt{SpecLink}, the concordance table, the attribution of the reviews and the wording of the rebuild descriptions), which are applied in this version; those earlier reports of Sol concerned the earlier texts; its review of version~1.2 is described below. These reviews do not certify the computational inputs or endorse a subsequent revision.
\item \emph{Claude Fable~5.2 (Anthropic)} acted as the project's pilot: it directed and arbitrated the work (the arbitrations are recorded in the project's registry). It wrote none of the proofs of this note, with one exception: on 12 November 2026 it wrote \S\ref{sec:bb26v5}, including the derivation \eqref{eq:bb26implies}, and the other sentences of this version that refer to \cite{BB26v5}, and it ran the final \texttt{pdflatex} compilation. That subsection was written after two independent readings of the same passages of \cite{BB26v5}, by \emph{DeepSeek V4 Pro (DeepSeek)} and by Qwen3.8, both without running programs. DeepSeek V4 Pro then checked the text of that subsection against the same passages (12 November 2026, without running programs): it confirmed the derivation \eqref{eq:bb26implies} and asked for two changes of wording (a literal quotation for hypothesis (W), and ``narrower'' instead of ``weaker'' for our conclusion), which are applied; it has not re-read the final wording.
\item \emph{Version 1.1 (12 November 2026).} A literature search was made on 12 November 2026, after the publication of version~1.0, by Claude Sonnet~5.5, which read the texts it obtained (scanned copies and arXiv files); the equation numbers and pages cited above for \cite{TY85,KT02,Kan} were checked on images of the pages. In the same search DeepSeek V4 Pro, without tools, compared excerpts of those sources with statements of the companion note on the twist ratio (not of the present note). The corrections of this version (Section~\ref{sec:intro}, \S\ref{sec:compare}, the bibliography and the status statements) were written by Claude Sonnet~5.5. One further review of this version by DeepSeek V4 Pro, without tools, on the text of the note and excerpts of the sources, is described in the file \texttt{RELECTURES.md} of the package. No statement, proof, constant or Lean file was changed.
\item \emph{GPT-6 Astra (OpenAI)} examined, blind, the electric-flux sectors under a scalar hypothesis (registry, 11 November 2026); that material is not part of this note. On 12 November 2026 GPT-6 Astra found the paper proof of the reverse implication (Theorem~\ref{thm:R}); GPT-6.1 Sol, DeepSeek V4 Pro and Claude Opus~5.5 read it the same day and reported no error (Sol with reservations on the wording of the statement, which are taken into account here); a first instance of Claude Sonnet~5.5 then checked each inequality, recomputed the constants, froze the four Lean statements and afterwards wrote their Lean proofs; the pilot (Claude Fable~5.2) redid the paper proof line by line and compared the frozen Lean statements with the theorem. The instance that prepared this version is a different one.
\item \emph{Version 1.2 (12 November 2026).} Claude Sonnet~5.5 prepared this version: it rebuilt the Lean files of the reverse implication as a non-author, on two toolchains and with an independent kernel, ran the mutants, wrote the section of Theorem~\ref{thm:R} from the paper proof (corrected where indicated in \path{RELECTURES.md}), rewrote \S\ref{sec:bb26v5}, ran a priority check against the current record of \cite{BB26v5} and the literature, and compiled the note. Two reviewers of other model families tried to break the final text, proof included, on 12 November 2026, in two rounds each: GPT-6.1 Sol (OpenAI), read-only, with file-reading tools (it ran no build, script or network request), and DeepSeek V4 Pro (DeepSeek), without tools. Neither found an error in Theorem R or in its constants; both asked for wording corrections (the minimum-side restriction in two summaries, the domains of (W2), the remark on degenerate eigenvalues, the round-trip factor, the attribution of the monotonicity, the provenance statements of this section), which are applied; their reports, and the answer to each objection, are in \path{RELECTURES.md}. These reviews concern the text of 12 November 2026, 07h--08h, not later revisions. In addition, a complete reading of the text of \cite{BB26v5} by DeepSeek V4 Pro, with quotations checked against the text by the preparer, was used for the comparison. These reviews do not certify later revisions. After these reviews, on 12 November 2026, the pilot (Claude Fable~5.2) corrected the description of who froze the Lean statements and who wrote the Lean proofs, removed a sentence about the public repository that was false, removed a conjecture on the optimal factor that our numerical checks do not support (Remark~\ref{rem:16}), added the output of the corresponding check (\path{reproduction/reciproque/t5.out}, rerun in a sandbox), and verified the reference \cite{BFS79} on the image of the page and its bibliographic data in Crossref; these last corrections were not reviewed by another model. No human has checked the proofs.
\item \emph{Version 1.3 (12 November 2026).} Sections~\ref{sec:thooft} and~\ref{sec:fixed} were written by Claude Opus~5.5 (statements and paper proofs), attacked by GPT-6.1 Sol and re-derived blind by GPT-6 Astra on the source text (Theorem~\ref{thm:L}), and read line by line by the pilot agent (Claude Fable~5.2). The seven abstract statements of the first Lean file of Appendix~\ref{app:lean13} were written by that same instance of Claude Opus~5.5, and were replayed, compared with the prose and tested against mutants by an instance of Claude Sonnet~5.5 that wrote none of them; the fifteen statements of the second file (the algebraic core of Section~\ref{sec:fixed}) were frozen and then proved by that instance of Claude Sonnet~5.5, and the frozen statements were read by the pilot agent. The version was assembled by an instance of Claude Sonnet~5.5, which also recomputed the constants, replayed the Lean files in two sandboxes and had the final text of the two new sections attacked by DeepSeek V4 Pro, without tools (four calls; objections and answers in \path{RELECTURES.md}). The other families (OpenAI) could not be consulted on the final text, the subscription being exhausted. The remarks of \S\ref{sec:fixed} on several planes (Remark~\ref{rem:several}) were written for the assembly and have been less reviewed. The literature search of \S\ref{sec:fixed} was made on 12 November 2026 by another instance of Claude Sonnet~5.5, which read the texts it could obtain from open sources and listed those it could not; the pilot agent checked each quotation used in that paragraph against the downloaded texts and wrote the paragraph. No human has checked the proofs.
\item \emph{Claude Sonnet~5.5 (Anthropic)} prepared version~1.0 of 12 November 2026 (this and the preceding paragraphs on the pilot and on the rebuilds of Appendix~\ref{app:lean} concern version~1.0; version~1.2 is described in the last paragraph of this list): it compared each formalised statement with the Lean
declarations, applied the corrections above to the text (no statement, proof, constant or Lean file was otherwise modified), ran the
\texttt{pdflatex} compilation, and ran reconstruction (c) of Appendix~\ref{app:lean} from the git archive of the package. Another instance of Claude Sonnet~5.5 performed
reconstruction (b) of that appendix.
\end{itemize}
These checks are not a substitute for refereeing: no human has checked the proofs.
The human author read the final text of this note before its release; he did not verify the proofs, and no human expert has checked them.

{\small
\begin{thebibliography}{99}\setlength{\itemsep}{1pt}
\bibitem{AT21} A.~Athenodorou, M.~Teper, \emph{SU(N) gauge theories in 3+1 dimensions: glueball spectrum, string tensions and
topology}, J.\ High Energy Phys.\ 2021, no.~12, 082, doi:10.1007/JHEP12(2021)082, arXiv:2106.00364.
\bibitem{BB26} D.~Bhattacharjee, U.~Bhattacharya, J.~M.~Lockwood, T.~Romanov, S.~Singha Roy, \emph{Yang--Mills theory and the mass
gap}, version 4.45.0 (10 November 2026), Zenodo, doi:10.5281/zenodo.24111624.
\bibitem{BB26v5} D.~Bhattacharjee, U.~Bhattacharya, J.~M.~Lockwood, T.~Romanov, I.~J.~Sessford, S.~Singha Roy, \emph{The Yang--Mills Mass
Gap from One Bound at One Scale}, version 5.0.0, Zenodo, doi:10.5281/zenodo.24147967 (record created on 11 November 2026 at 23:29:13 UTC;
publication date 12 November 2026; paper of 157 pages in the archive \path{creelie/4Y-M-v5.0.0.zip}). Read on 12 November 2026: the passages listed in \S\ref{sec:bb26v5}, the whole text by keyword searches, and (for version~1.2) a complete reading by DeepSeek V4 Pro with quotations checked against the text.
\bibitem{BFS79} D.~Brydges, J.~Fr\"ohlich, E.~Seiler, \emph{On the construction of quantized gauge fields. I. General results}, Ann.\ Phys.\ 121 (1979) 227--284 doi:10.1016/0003-4916(79)90098-8 (our scanned copy is the preprint IHES/P/78/237, June 1978; the argument cited is on its p.~38, page 39 of the scan).
\bibitem{BK90} C.~Borgs, R.~Koteck\'y, \emph{A rigorous theory of finite-size scaling at first-order phase transitions}, J.\ Stat.\
Phys.\ 61 (1990) 79--119, doi:10.1007/BF01013955.
\bibitem{GM82} M.~G\"opfert, G.~Mack, \emph{Proof of confinement of static quarks in 3-dimensional $U(1)$ lattice gauge theory for all
values of the coupling constant}, Commun.\ Math.\ Phys.\ 82 (1982) 545--606, doi:10.1007/BF01961240.
\bibitem{GM16} D.~Gosset, E.~Mozgunov, \emph{Local gap threshold for frustration-free spin systems}, J.\ Math.\ Phys.\ 57 (2016)
091901, doi:10.1063/1.4962337, arXiv:1512.00088.
\bibitem{Kan} T.~Kanazawa, \emph{Generalizing the Tomboulis--Yaffe inequality to $SU(N)$ lattice gauge theories and general classical spin systems}, Ann.\ Phys.\ 324 (2009) 1634--1665, arXiv:0808.3442v3 [math-ph] (page numbers cited here are those of the arXiv version).
\bibitem{Kau49} B.~Kaufman, \emph{Crystal statistics. II. Partition function evaluated by spinor analysis}, Phys.\ Rev.\ 76 (1949)
1232--1243, doi:10.1103/PhysRev.76.1232.
\bibitem{KT02} T.~G.~Kov\'acs, E.~T.~Tomboulis, \emph{Vortex waistlines and long range fluctuations}, Phys.\ Rev.\ D 65 (2002) 074501, arXiv:hep-lat/0108017v2 (page numbers cited here are those of the arXiv version).
\bibitem{Kna88} S.~Knabe, \emph{Energy gaps and elementary excitations for certain VBS-quantum antiferromagnets}, J.\ Stat.\ Phys.\ 52
(1988) 627--638, doi:10.1007/BF01019721.
\bibitem{Lem19} M.~Lemm, \emph{Finite-size criteria for spectral gaps in $D$-dimensional quantum spin systems}, arXiv:1902.07141 (2019).
\bibitem{Lus77} M.~L\"uscher, \emph{Construction of a selfadjoint, strictly positive transfer matrix for Euclidean lattice gauge
theories}, Commun.\ Math.\ Phys.\ 54 (1977) 283--292, doi:10.1007/BF01614090.
\bibitem{OAI26} OpenAI, \emph{Sharp mass bounds for the two-dimensional $O(4)$ model}, preprint dated 23 September 2026, repository
\url{https://github.com/openai/math} (family 215, preprint directory \path{Sharp-mass-bounds-for-the-two-dimensional-O4-model-September-23-2026},
\path{sections/finite-volume.tex}); read from a local copy of the sources, not re-fetched on 12 November 2026.
\bibitem{OS73} K.~Osterwalder, R.~Schrader, \emph{Axioms for Euclidean Green's functions}, Commun.\ Math.\ Phys.\ 31 (1973) 83--112,
doi:10.1007/BF01645738.
\bibitem{OS75} K.~Osterwalder, R.~Schrader, \emph{Axioms for Euclidean Green's functions II}, Commun.\ Math.\ Phys.\ 42 (1975)
281--305, doi:10.1007/BF01608978.
\bibitem{OS78} K.~Osterwalder, E.~Seiler, \emph{Gauge field theories on a lattice}, Ann.\ Phys.\ 110 (1978) 440--471,
doi:10.1016/0003-4916(78)90039-8.
\bibitem{Raj26} H.~T.~Rajabu, \emph{Multiscale covariance transfer and a conditional continuum mass gap for four-dimensional $SU(2)$
lattice Yang--Mills theory}, version 2.1 (29 October 2026), in the Zenodo record doi:10.5281/zenodo.23729328.
\bibitem{CGJK85} A.~Coste, A.~Gonz\'alez-Arroyo, J.~Jurkiewicz, C.~P.~Korthals Altes, \emph{Zero-momentum contribution to Wilson loops in periodic boxes}, Nucl.\ Phys.\ B 262 (1985) 67--94, doi:10.1016/0550-3213(85)90064-1 (not obtained).
\bibitem{tH79} G.~'t~Hooft, \emph{A property of electric and magnetic flux in non-Abelian gauge theories}, Nucl.\ Phys.\ B 153 (1979) 141--160, doi:10.1016/0550-3213(79)90595-9 (not read by us; cited after \cite{TY85}).
\bibitem{TBv12} K.~R\'emondi\`ere, \emph{An elementary twist-ratio bound for the $SU(2)$ Wilson lattice, formalised in Lean~4 on the objects of OpenAI's four-dimensional development and its extension to a cyclic centre}, note v1.2, Zenodo, doi:10.5281/zenodo.24153930 (Section on the flux-sector bound on a tube).
\bibitem{TY85} E.~T.~Tomboulis, L.~G.~Yaffe, \emph{Finite temperature $SU(2)$ lattice gauge theory}, Commun.\ Math.\ Phys.\ 100 (1985) 313--341 (read in a scanned copy with an OCR text layer).
\bibitem{vB82} P.~van Baal, \emph{Some results for $SU(N)$ gauge fields on the hypertorus}, Commun.\ Math.\ Phys.\ 85 (1982) 529--547, doi:10.1007/BF01403503 (read: Section~4).
\bibitem{AF80} J.~Ambj{\o}rn, H.~Flyvbjerg, \emph{'t~Hooft's non-abelian magnetic flux has zero classical energy}, Phys.\ Lett.\ B 97 (1980) 241--245, doi:10.1016/0370-2693(80)90592-4 (not obtained; cited after \cite{vB82}).
\bibitem{GJK81} J.~Groeneveld, J.~Jurkiewicz, C.~P.~Korthals Altes, \emph{Twist as a probe for phase structure}, Phys.\ Scr.\ 23 (1981) 1022--1031, doi:10.1088/0031-8949/23/5B/022 (not obtained; abstract only).
\bibitem{GAO83} A.~Gonz\'alez-Arroyo, M.~Okawa, \emph{The twisted Eguchi--Kawai model: a reduced model for large-$N$ lattice gauge theory}, Phys.\ Rev.\ D 27 (1983) 2397, doi:10.1103/PhysRevD.27.2397 (not obtained).
\bibitem{Lu83} M.~L\"uscher, \emph{Some analytic results concerning the mass spectrum of Yang--Mills gauge theories on a torus}, Nucl.\ Phys.\ B 219 (1983) 233--261, doi:10.1016/0550-3213(83)90436-4 (not obtained).
\bibitem{GA98} A.~Gonz\'alez-Arroyo, \emph{Yang--Mills fields on the 4-dimensional torus. Part I: Classical theory}, arXiv:hep-th/9807108 (read: Sections~3--5 and 8.1).
\bibitem{vB00} P.~van Baal, \emph{QCD in a finite volume}, arXiv:hep-ph/0008206 (read: Section~3, in particular the beginning of Section~3 and Section~3.4).
\bibitem{KvB86} J.~Koller, P.~van Baal, \emph{A rigorous nonperturbative result for the glueball mass and electric flux energy in a finite volume}, Nucl.\ Phys.\ B 273 (1986) 387--412 (read: Section~1).
\bibitem{KS98} W.~Krauth, M.~Staudacher, \emph{Finite Yang--Mills integrals}, Phys.\ Lett.\ B 435 (1998) 350--355, arXiv:hep-th/9804199 (read).
\bibitem{AW01} P.~Austing, J.~F.~Wheater, \emph{Convergent Yang--Mills matrix theories}, JHEP 04 (2001) 019, arXiv:hep-th/0103159 (read: abstract and Sections~2--3).
\bibitem{KT00} T.~G.~Kov\'acs, E.~T.~Tomboulis, \emph{Computation of the vortex free energy in $SU(2)$ gauge theory}, Phys.\ Rev.\ Lett.\ 85 (2000) 704--707, arXiv:hep-lat/0002004 (read).
\bibitem{dFvS02} Ph.~de~Forcrand, L.~von~Smekal, \emph{'t~Hooft loops, electric flux sectors and confinement in $SU(2)$ Yang--Mills theory}, Phys.\ Rev.\ D 66 (2002) 011504, arXiv:hep-lat/0107018 (read).
\bibitem{ClayYM} K.~R\'emondi\`ere, \emph{ClayYM: a candidate Lean~4 statement of the Yang--Mills existence and mass gap problem, with
compiled guards against fake models}, version 1.1.0 (25 October 2026), Zenodo, doi:10.5281/zenodo.23644079 (all versions).
\bibitem{internal} Internal notes of the project, 25 October and 10 November 2026 (unpublished).
\end{thebibliography}}

\end{document}
